% Generated mechanically from frozen input p09/ko/src/spaces-limits.tex.
% Do not edit this generated file; edit the bound locale source instead.

% OK, start here.
%

\setcounter{chapter}{69}
\chapter{대수공간의 극한}



\phantomsection
\label{spaces-limits-section-phantom}


\section{서론}
\label{spaces-limits-section-introduction}

\noindent
이 장에서는 대수공간의 극한에 관한 내용을 다룬다.
먼저 대수공간 $F$가 밑 $S$ 위에서 국소 유한 표시라는 것을
극한을 보존하는 함자라는 조건으로 특성화한다.
이어서 유향 집합(Categories, Definition \ref{categories-definition-directed-set})을
지표로 하고 전이 사상이 아핀인 역계의 극한을 연구한다.
\cite{CLO}를 따라 준콤팩트 준분리 대수공간의 절대 Noether 근사를 논의한다.
David Rydh의 또 다른 접근법(\cite{rydh_approx} 참조)은
특정 대수 스택의 절대 Noether 근사까지 다룬다.


\section{규약}
\label{spaces-limits-section-conventions}

\noindent
모든 스킴은 큰 fppf 사이트 $\Sch_{fppf}$에 속한다고 항상 가정한다.
또한 다루는 모든 환 $A$에 대해 $\Spec(A)$는 이 큰 사이트의
어떤 대상과 (동형인) 것으로 가정한다.

\medskip\noindent
$S$를 스킴, $X$를 $S$ 위의 대수공간이라 하자.
이 장과 이후의 장에서는 $X \times_S X$라는 표기를 쓴다.
이는 $X$와 자기 자신의 곱($S$ 위의 대수공간의 범주에서 취한 곱)에 대해
$X \times X$ 대신 쓰는 표기이다.











\section{유한 표시 사상}
\label{spaces-limits-section-finite-presentation}

\noindent
이 절에서는 Limits, Proposition
\ref{limits-proposition-characterize-locally-finite-presentation}을
대수공간의 사상으로 일반화한다.
다음 정의의 동기는 방금 인용한 명제에서 나온다.

\begin{definition}
\label{spaces-limits-definition-locally-finite-presentation}
$S$를 스킴이라 하자.
\begin{enumerate}
\item 함자 $F : (\Sch/S)_{fppf}^{opp} \to \textit{Sets}$가
{\it 극한을 보존한다} 또는 {\it 국소 유한 표시이다}라는 것은,
아핀 스킴 $T$가 $S$ 위에서 극한 $T = \lim T_i$로 주어지되,
이것이 아핀 스킴 $T_i$들의 $S$ 위의 유향 역계의 극한일 때마다
$$
F(T) = \colim F(T_i).
$$
가 성립한다는 뜻이다.
이때 $F$가 {\it $S$ 위에서 국소 유한 표시이다}라고도 한다.
\item $F, G : (\Sch/S)_{fppf}^{opp} \to \textit{Sets}$라 하자.
함자 사이의 변환 $a : F \to G$가
{\it 극한을 보존한다} 또는 {\it 국소 유한 표시이다}라는 것은,
모든 스킴 $T$가 $S$ 위에 주어졌을 때 모든 $y \in G(T)$에 대해 함자
$$
F_y : (\Sch/T)_{fppf}^{opp} \longrightarrow \textit{Sets}, \quad
T'/T \longmapsto \{x \in F(T') \mid a(x) = y|_{T'}\}
$$
가 $T$ 위에서 국소 유한 표시라는 뜻이다\footnote{Lemma
\ref{spaces-limits-lemma-characterize-relative-limit-preserving}의 특성화 (2)가
더 이해하기 쉬울 수 있다.}.
이때 $F$가 $G$ 위에서 {\it 상대적으로 극한을 보존한다}라고도 한다.
\end{enumerate}
\end{definition}

\noindent
함자 $F_y$는 어떤 의미에서 $a : F \to G$의 $y$ 위의 올이다.
다만 $T$의 큰 fppf 사이트 위의 전층이라는 차이가 있다.
이 함자는 다음 식으로 주어진다.
\begin{equation}
\label{spaces-limits-equation-fibre-map-functors}
F_y =
F|_{(\Sch/T)_{fppf}}
{\times}_{G|_{(\Sch/T)_{fppf}}}
*
\end{equation}
여기서 $*$는 $(\Sch/T)_{fppf}$ 위의 (전)층의 범주의 최종 대상이고
(Sites, Example \ref{sites-example-singleton-sheaf} 참조),
사상 $* \to G|_{(\Sch/T)_{fppf}}$는 $y$로 주어진다.
$j : (\Sch/T)_{fppf} \to (\Sch/S)_{fppf}$가 국소화 함자이면,
위 식은 $F_y = j^{-1}F \times_{j^{-1}G} *$가 되며
$j_!F_y$는 바로 올곱 $F \times_{G, y} T$라는 점에 주의하자.
(국소화에 대해서는 Sites, Section \ref{sites-section-localize}를,
특히 전층에 대한 $j_!$에 대해서는 Sites, Remark
\ref{sites-remark-localize-presheaves}를 참조하라.)

\medskip\noindent
현재로서는 사상 $X \to Y$가 $S$ 위의 대수공간의 사상으로서 국소 유한 표시라는
말에 일시적으로 두 가지 정의가 있다.
하나는 Morphisms of Spaces,
Definition \ref{spaces-morphisms-definition-locally-finite-presentation}에 따른
정의이고, 다른 하나는 $X \to Y$가 함자 사이의 변환이라는 사실을 이용하여
Definition \ref{spaces-limits-definition-locally-finite-presentation}을 적용하는 정의이다
(후자의 개념에는 가능한 한 ``극한 보존''이라는 용어를 쓸 것이다).
Proposition \ref{spaces-limits-proposition-characterize-locally-finite-presentation}에서
이 두 정의가 일치함을 보인다.

\begin{lemma}
\label{spaces-limits-lemma-characterize-relative-limit-preserving}
$S$를 스킴이라 하자.
$a : F \to G$를 함자
$(\Sch/S)_{fppf}^{opp} \to \textit{Sets}$ 사이의 변환이라 하자.
다음 조건들은 동치이다.
\begin{enumerate}
\item $a : F \to G$는 극한을 보존한다.
\item 아핀 스킴 $T$가 $S$ 위에서 극한 $T = \lim T_i$로 주어지되,
이것이 아핀 스킴 $T_i$들의 $S$ 위의 유향 역계의 극한일 때마다 집합의 가환도
$$
\xymatrix{
\colim_i F(T_i) \ar[r] \ar[d]_a & F(T) \ar[d]^a \\
\colim_i G(T_i) \ar[r] & G(T)
}
$$
는 올곱 가환도이다.
\end{enumerate}
\end{lemma}

\begin{proof}
(1)을 가정하자. (2)에서와 같은 $T = \lim_{i \in I} T_i$를 생각하자.
$(y, x_T)$를 올곱
$\colim_i G(T_i) \times_{G(T)} F(T)$의 원소라 하자.
그러면 $y$는 $y_i \in G(T_i)$에서 오며, 여기서 $i$는 적절히 택한다.
함자 $F_{y_i}$를 $(\Sch/T_i)_{fppf}$ 위에서 생각하자.
이는 Definition \ref{spaces-limits-definition-locally-finite-presentation}에서 정의한 함자이다.
$x_T \in F_{y_i}(T)$임을 알 수 있다.
또한 $T = \lim_{i' \geq i} T_{i'}$는 $T_i$ 위의 아핀 스킴의 유향계이다.
따라서 (1)에 의해 $x_T$는 유일한 원소 $x$의 상이며,
이 원소는 $\colim_{i' \geq i} F_{y_i}(T_{i'})$에 속한다.
그러므로 $x$는 $\colim F(T_i)$의 원소 중 쌍 $(y, x_T)$로 가는 유일한 원소이다.
이로써 (2)가 성립함을 보였다.

\medskip\noindent
(2)를 가정하자. $T$를 스킴이라 하고 $y_T \in G(T)$라 하자.
$F_{y_T}$가 극한을 보존함을 보여야 한다.
$T' = \lim_{i \in I} T'_i$를 $T$ 위의 아핀 스킴 $T'_i$의 유향 극한인
$T$ 위의 아핀 스킴이라 하자. $x_{T'} \in F_{y_T}$라 하자.
$i \in I$를 택하자. 이는 $I$가 유향 집합이므로 가능하다.
% P09-ERR-0050; P09-ERR-0051: canonical source notation retained.
$y_i \in F(T'_i)$로 $y_{T'}$의 상을 나타내자.
그러면 $(y_i, x_{T'})$가 올곱
$\colim_i G(T'_i) \times_{G(T')} F(T')$의 원소임을 알 수 있다.
따라서 (2)에 의해 유일한 원소 $x$를 $\colim_i F(T'_i)$에서 얻으며,
이는 $(y_i, x_{T'})$로 간다.
$x$는 $\colim_i F_y(T'_i)$의 원소로서 $x_{T'}$로 가는 원소를 정의함이 분명하므로
증명이 끝난다.
% P09-ERR-1454: y_{T'} and F_y retained as printed in the canonical source.
\end{proof}

\begin{lemma}
\label{spaces-limits-lemma-composition-locally-finite-presentation}
$S$를 $\Sch_{fppf}$에 속하는 스킴이라 하자.
$F, G, H : (\Sch/S)_{fppf}^{opp} \to \textit{Sets}$라 하자.
$a : F \to G$, $b : G \to H$를 함자 사이의 변환이라 하자.
$a$와 $b$가 극한을 보존하면
$$
b \circ a : F \longrightarrow H
$$
도 극한을 보존한다.
\end{lemma}

\begin{proof}
$T = \lim_{i \in I} T_i$가
Lemma \ref{spaces-limits-lemma-characterize-relative-limit-preserving}의 특성화 (2)에서와
같다고 하자. 다음 집합의 가환도를 생각하자.
$$
\xymatrix{
\colim_i F(T_i) \ar[r] \ar[d]_a & F(T) \ar[d]^a \\
\colim_i G(T_i) \ar[r] \ar[d]_b & G(T) \ar[d]^b \\
\colim_i H(T_i) \ar[r] & H(T)
}
$$
가정에 의해 두 정사각형은 올곱 정사각형이다.
따라서 바깥 직사각형도 올곱 가환도이므로 보조정리가 증명된다.
\end{proof}

\begin{lemma}
\label{spaces-limits-lemma-locally-finite-presentation-permanence}
$S$를 $\Sch_{fppf}$에 속하는 스킴이라 하자.
$F, G, H : (\Sch/S)_{fppf}^{opp} \to \textit{Sets}$라 하자.
$a : F \to G$, $b : G \to H$를 함자 사이의 변환이라 하자.
$b \circ a$와 $b$가 극한을 보존하면 $a$도 극한을 보존한다.
\end{lemma}

\begin{proof}
$T = \lim_{i \in I} T_i$가
Lemma \ref{spaces-limits-lemma-characterize-relative-limit-preserving}의 특성화 (2)에서와
같다고 하자. 다음 집합의 가환도를 생각하자.
$$
\xymatrix{
\colim_i F(T_i) \ar[r] \ar[d]_a & F(T) \ar[d]^a \\
\colim_i G(T_i) \ar[r] \ar[d]_b & G(T) \ar[d]^b \\
\colim_i H(T_i) \ar[r] & H(T)
}
$$
가정에 의해 아래 정사각형과 바깥 직사각형은 집합의 올곱이다.
따라서 위 정사각형도 올곱 정사각형이므로 보조정리가 증명된다.
\end{proof}

\begin{lemma}
\label{spaces-limits-lemma-base-change-locally-finite-presentation}
$S$를 $\Sch_{fppf}$에 속하는 스킴이라 하자.
$F, G, H : (\Sch/S)_{fppf}^{opp} \to \textit{Sets}$라 하자.
$a : F \to G$, $b : H \to G$를 함자 사이의 변환이라 하자.
다음 올곱 가환도를 생각하자.
$$
\xymatrix{
H \times_{b, G, a} F \ar[r]_-{b'} \ar[d]_{a'} & F \ar[d]^a \\
H \ar[r]^b & G
}
$$
$a$가 극한을 보존하면 그 밑변환 $a'$도 극한을 보존한다.
\end{lemma}

\begin{proof}
생략한다. 힌트: 형식적으로 따른다.
\end{proof}

\begin{lemma}
\label{spaces-limits-lemma-fibre-product-locally-finite-presentation}
$S$를 $\Sch_{fppf}$에 속하는 스킴이라 하자.
$E, F, G, H : (\Sch/S)_{fppf}^{opp} \to \textit{Sets}$라 하자.
$a : F \to G$, $b : H \to G$, $c : G \to E$를
함자 사이의 변환이라 하자. $c$, $c \circ a$, $c \circ b$가
극한을 보존하면 $F \times_G H \to E$도 극한을 보존한다.
\end{lemma}

\begin{proof}
$T = \lim_{i \in I} T_i$가
Lemma \ref{spaces-limits-lemma-characterize-relative-limit-preserving}의 특성화 (2)에서와
같다고 하자. 그러면 여과 여극한은 유한 곱과 교환하므로
$$
\colim (F \times_G H)(T_i) =
\colim F(T_i) \times_{\colim G(T_i)} \colim H(T_i)
$$
가 성립한다. 따라서 다음 가환도가 올곱 가환도임을 보이면 된다.
$$
\xymatrix{
\colim F(T_i) \times_{\colim G(T_i)} \colim H(T_i) \ar[r] \ar[d] &
F(T) \times_{G(T)} H(T) \ar[d] \\
\colim_i E(T_i) \ar[r] & E(T)
}
$$
이는 집합의 사상 $E' \to E$, $F \to G$, $H \to G$, $G \to E$가
주어지면
$$
E' \times_E (F \times_G H) =
(E' \times_E F) \times_{(E' \times_E G)} (E' \times_E H)
$$
가 성립한다는 관찰에서 따른다. 몇 가지 세부 사항은 생략한다.
\end{proof}

\begin{lemma}
\label{spaces-limits-lemma-sheafify-finite-presentation}
$S$를 $\Sch_{fppf}$에 속하는 스킴이라 하자.
$F : (\Sch/S)_{fppf}^{opp} \to \textit{Sets}$를 함자라 하자.
$F$가 극한을 보존하면 그 층화 $F^\#$도 극한을 보존한다.
\end{lemma}

\begin{proof}
$F$가 극한을 보존한다고 가정하자.
$F^+$가 극한을 보존함을 보이면 충분하다.
이는 $F^\# = (F^+)^+$이기 때문이다.
Sites, Theorem \ref{sites-theorem-plus}를 보라.
$T$를 $S$ 위의 아핀 스킴이라 하고, $T = \lim T_i$를
아핀 $S$-스킴의 역계의 유향 극한으로 나타내자.
$F^+(T)$는 $\check H^0(\mathcal{V}, F)$의 여극한임을 상기하자.
여기서 극한은 $T$의 모든 $(\Sch/S)_{fppf}$-덮개에 걸쳐 취한다.
% P09-ERR-0052: source says limit after defining a colimit.
아핀 스킴의 모든 fppf 덮개는 표준 fppf 덮개로 세분할 수 있다.
Topologies, Lemma \ref{topologies-lemma-fppf-affine}를 보라.
따라서 다음과 같이 쓸 수 있다.
$$
F^+(T)
=
\colim_{\mathcal{V}\text{ 표준 덮개 }T}
\check H^0(\mathcal{V}, F).
$$
이 여극한에 나타나는 임의의 $\mathcal{V} = \{T_k \to T\}_{k = 1, \ldots, n}$는
$\mathcal{V}_i \times_{T_i} T$의 형태로 쓸 수 있다.
여기서 $i$를 적절히 택하고 표준 fppf 덮개
$\mathcal{V}_i = \{T_{i, k} \to T_i\}_{k = 1, \ldots, n}$를 $T_i$ 위에서 택한다.
$\mathcal{V}_{i'} = \{T_{i', k} \to T_{i'}\}_{k = 1, \ldots, n}$로
$i' \geq i$에 대한 밑변환을 나타내자. 그러면 다음을 얻는다.
% P09-ERR-0053; P09-ERR-0054: both frozen index discrepancies retained.
\begin{align*}
\colim_{i' \geq i} \check H^0(\mathcal{V}_i, F)
& =
\colim_{i' \geq i}
\text{동등자}
\left(
\xymatrix{
\prod F(T_{i', k})
\ar@<1ex>[r] \ar@<-1ex>[r] &
\prod F(T_{i', k} \times_{T_{i'}} T_{i', l})
}
\right)
\\
& =
\text{동등자}
\left(
\xymatrix{
\colim_{i' \geq i}
\prod F(T_{i', k})
\ar@<1ex>[r] \ar@<-1ex>[r] &
\colim_{k' \geq k}
\prod F(T_{i', k} \times_{T_{i'}} T_{i', l})
}
\right)
\\
& =
\text{동등자}
\left(
\xymatrix{
\prod F(T_k)
\ar@<1ex>[r] \ar@<-1ex>[r] &
\prod F(T_k \times_T T_l)
}
\right)
\\
& =
\check H^0(\mathcal{V}, F)
\end{align*}
여기서 두 번째 등식은 여과 여극한이 완전하기 때문에 성립한다.
세 번째 등식은 $F$가 극한을 보존하고,
$\lim_{i' \geq i} T_{i', k} = T_k$ 및
$\lim_{i' \geq i} T_{i', k} \times_{T_{i'}} T_{i', l} = T_k \times_T T_l$가
성립하기 때문에 성립한다.
이는 Limits, Lemma \ref{limits-lemma-scheme-over-limit}에 따른다.
이를 모든 덮개에 동시에 적용하면 다음을 얻는다.
\begin{align*}
F^+(T)
& =
\colim_{\mathcal{V}\text{ 표준 덮개 }T} \check H^0(\mathcal{V}, F) \\
& =
\colim_{i \in I}
\colim_{\mathcal{V}_i\text{ 표준 덮개 }T_i}
\check H^0(T \times_{T_i}\mathcal{V}_i, F) \\
& =
\colim_{i \in I} F^+(T_i)
\end{align*}
여극한을 취하는 순서는 Categories, Lemma
\ref{categories-lemma-colimits-commute}에 의해 바꿀 수 있다.
\end{proof}

\begin{lemma}
\label{spaces-limits-lemma-sheaf-finite-presentation}
$S$를 스킴이라 하자.
$F : (\Sch/S)_{fppf}^{opp} \to \textit{Sets}$를 함자라 하자.
다음을 가정하자.
\begin{enumerate}
\item $F$는 층이다.
\item fppf 덮개 $\{U_j \to S\}_{j \in J}$가 존재하여
$F|_{(\Sch/U_j)_{fppf}}$가 극한을 보존한다.
\end{enumerate}
그러면 $F$는 극한을 보존한다.
\end{lemma}

\begin{proof}
$T$를 $S$ 위의 아핀 스킴이라 하자.
$I$를 유향 집합이라 하고, $T_i$를 $S$ 위의 아핀 스킴의 역계로서
$T = \lim T_i$를 만족하는 것이라 하자.
표준 사상 $\colim F(T_i) \to F(T)$가 전단사임을 보여야 한다.

\medskip\noindent
$0 \in I$를 하나 택하고, 표준 fppf 덮개
$\{V_{0, k} \to T_{0}\}_{k = 1, \ldots, m}$를 택하여
당김 $\{U_j \times_S T_0 \to T_0\}$를 세분하게 하자.
이 당김은 주어진 $S$의 fppf 덮개에서 얻은 것이다.
각 $i \geq 0$에 대해 $V_{i, k} = T_i \times_{T_0} V_{0, k}$로 놓고,
$V_k = T \times_{T_0} V_{0, k}$로 놓자.
$V_k = \lim_{i \geq 0} V_{i, k}$임에 주의하자.
Limits, Lemma \ref{limits-lemma-scheme-over-limit}를 보라.

\medskip\noindent
$x, x' \in \colim F(T_i)$가 $F(T)$의 같은 원소로 간다고 하자.
$x, x'$가 원소 $x_i, x'_i \in F(T_i)$로 주어진다고 하자.
여기서 $i \in I$이다(두 원소에 대해 같은 $i$를 택할 수 있는데,
이는 $I$가 유향이기 때문이다).
가정 (2)와 $x_i, x'_i$가 $F(T)$의 같은 원소로 간다는 사실에 의해
$$
x_i|_{V_{i', k}} = x'_i|_{V_{i', k}}
$$
가 충분히 큰 $i' \in I$에 대해 성립한다.
같은 $i'$를 모든 $k$에 대해 택할 수 있는데,
이는 $k \in \{1, \ldots, m\}$가 유한 집합을 움직이기 때문이다.
$\{V_{i', k} \to T_{i'}\}$가 fppf 덮개이고 $F$가 층이므로,
이로부터 원하는 등식 $x_i|_{T_{i'}} = x'_i|_{T_{i'}}$를 얻는다.
따라서 사상 $\colim F(T_i) \to F(T)$는 단사이다.

\medskip\noindent
전사성도 비슷하게 보인다.
$x \in F(T)$라 하자. 가정 (2)에 의해 각 $k$에 대해
$i$를 택하여 $x|_{V_k}$가 원소 $x_{i, k} \in F(V_{i, k})$에서 오게 할 수 있다.
앞에서와 같이 하나의 $i$를 택하여 모든 $k$에 대해 동시에 성립하게 할 수 있다.
위에서 증명한 단사성에 의해
$$
x_{i, k}|_{V_{i', k} \times_{T_{i'}} V_{i', l}}
=
x_{i, l}|_{V_{i', k} \times_{T_{i'}} V_{i', l}}
$$
가 충분히 큰 $i'$에 대해 성립한다.
따라서 $F$의 층 조건에 의해 원소 $x_{i, k}|_{V_{i', k}}$들을 붙여
원소 $x_{i'} \in F(T_{i'})$를 얻으며, 이것이 원하는 결과이다.
\end{proof}

\begin{lemma}
\label{spaces-limits-lemma-sheafify-finite-presentation-map}
$S$를 $\Sch_{fppf}$에 속하는 스킴이라 하자.
$F, G : (\Sch/S)_{fppf}^{opp} \to \textit{Sets}$를 함자라 하자.
$a : F \to G$가 극한을 보존하는 변환이면,
유도되는 층 사이의 변환 $F^\# \to G^\#$도 극한을 보존한다.
\end{lemma}

\begin{proof}
$T$를 스킴이라 하고 $y \in G^\#(T)$라 하자.
함자 $F^\#_y : (\Sch/T)_{fppf}^{opp} \to \textit{Sets}$가
극한을 보존함을 보여야 한다.
이 함자는 $F^\# \to G^\#$와 $y$로부터
Definition \ref{spaces-limits-definition-locally-finite-presentation}에서와 같이 구성된다.
식 (\ref{spaces-limits-equation-fibre-map-functors})에 의해 $F^\#_y$는 층이다.
fppf 덮개 $\{V_j \to T\}_{j \in J}$를 택하여 $y|_{V_j}$가
원소 $y_j \in F(V_j)$에서 오게 하자.
% P09-ERR-0055; P09-ERR-0056: canonical typing and restriction retained.
$F^\#$를 $(\Sch/V_j)_{fppf}$에 제한하면 바로 $F^\#_{y_j}$가 됨에 주의하자.
$F^\#_{y_j}$가 극한을 보존함을 보일 수 있으면,
Lemma \ref{spaces-limits-lemma-sheaf-finite-presentation}에 의해
$F^\#_y$도 극한을 보존하므로 증명이 끝난다.
이로써 $y \in G(T)$인 경우로 환원된다.

\medskip\noindent
$y \in G(T)$라 하자. 이 경우 $F^\#_y = (F_y)^\#$라고 주장한다.
이는 식 (\ref{spaces-limits-equation-fibre-map-functors})에서 따른다.
따라서 이 경우는 Lemma \ref{spaces-limits-lemma-sheafify-finite-presentation}에서 따른다.
\end{proof}

\begin{proposition}
\label{spaces-limits-proposition-characterize-locally-finite-presentation}
$S$를 스킴이라 하자. $f : X \to Y$를 $S$ 위의 대수공간의 사상이라 하자.
다음 조건들은 동치이다.
\begin{enumerate}
\item 사상 $f$는 국소 유한 표시인 대수공간의 사상이다.
Morphisms of Spaces, Definition
\ref{spaces-morphisms-definition-locally-finite-presentation}를 보라.
\item 사상 $f : X \to Y$는 함자 사이의 변환으로서 극한을 보존한다.
Definition \ref{spaces-limits-definition-locally-finite-presentation}을 보라.
\end{enumerate}
\end{proposition}

\begin{proof}
(1)을 가정하자. $T$를 스킴이라 하고 $y \in Y(T)$라 하자.
$T \times_Y X$가 $T$ 위에서
Definition \ref{spaces-limits-definition-locally-finite-presentation}의 의미로
극한을 보존함을 보여야 한다.
따라서 $X$가 대수공간으로서 $S$ 위에서 국소 유한 표시이면,
함자 $X : (\Sch/S)_{fppf}^{opp} \to \textit{Sets}$로서 극한을 보존함을
보이는 것으로 환원된다. 이를 위해 표시 $X = U/R$를 택하자.
Spaces, Definition \ref{spaces-definition-presentation}을 보라.
Morphisms of Spaces,
Definition \ref{spaces-morphisms-definition-locally-finite-presentation}에 의해
$U$와 $R$은 모두 $S$ 위에서 국소 유한 표시인 스킴이다.
따라서 Limits, Proposition
\ref{limits-proposition-characterize-locally-finite-presentation}에 의해
$$
U(T) = \colim U(T_i), \quad
R(T) = \colim R(T_i)
$$
가 $T = \lim_i T_i$가 $(\Sch/S)_{fppf}$에서 성립할 때마다 성립한다.
따라서 전층
$$
(\Sch/S)_{fppf}^{opp} \longrightarrow \textit{Sets}, \quad
W \longmapsto U(W)/R(W)
$$
는 극한을 보존한다.
그러므로 Lemma \ref{spaces-limits-lemma-sheafify-finite-presentation}에 의해
그 층화 $X = U/R$도 극한을 보존한다.

\medskip\noindent
(2)를 가정하자. 스킴 $V$와 전사 에탈 사상 $V \to Y$를 택하자.
이어서 스킴 $U$와 전사 에탈 사상 $U \to V \times_Y X$를 택하자.
Lemma \ref{spaces-limits-lemma-base-change-locally-finite-presentation}에 의해
함자 사이의 변환 $V \times_Y X \to V$는 극한을 보존한다.
Morphisms of Spaces,
Lemma \ref{spaces-morphisms-lemma-etale-locally-finite-presentation}에 의해
대수공간의 사상 $U \to V \times_Y X$는 국소 유한 표시이다.
따라서 증명의 첫 부분에 의해 함자 사이의 변환으로서 극한을 보존한다.
Lemma \ref{spaces-limits-lemma-composition-locally-finite-presentation}에 의해
합성 $U \to V \times_Y X \to V$는 함자 사이의 변환으로서 극한을 보존한다.
따라서 스킴의 사상 $U \to V$는 Limits, Proposition
\ref{limits-proposition-characterize-locally-finite-presentation}에 의해
국소 유한 표시이다(증명의 마지막 문단에 나오는 집합론적 논점을 전제한다).
이는 정의에 의해 (1)이 성립한다는 뜻이다.

\medskip\noindent
집합론적 논점. $U \to V$를 $(\Sch/S)_{fppf}$의 사상이라 하자.
Limits, Proposition
\ref{limits-proposition-characterize-locally-finite-presentation}에서는
$U \to V$가 국소 유한 표시라는 것을, 아핀 스킴의 {\it 모든} 유향 역계
$(T_i, f_{ii'})$가 $V$ 위에 주어질 때
$U(T) = \colim V(T_i)$가 성립한다는 조건으로 특성화한다.
% P09-ERR-0057: V(T_i) on the right is retained from the frozen source.
하지만 현재 설정에서는 아핀 스킴 $T_i$가 $V$ 위에 있고
$(\Sch/S)_{fppf}$의 어떤 대상과 (동형인) 경우만 고려할 수 있다.
따라서 증명이 성립하도록 $(\Sch/S)_{fppf}$ 안에
충분히 많은 아핀 스킴이 있는지 확인해야 한다.
Limits, Proposition
\ref{limits-proposition-characterize-locally-finite-presentation}의
(2) $\Rightarrow$ (1)의 증명을 살펴보면,
문제는 $U$와 $V$가 아핀인 경우로 환원됨을 알 수 있다.
$U = \Spec(A)$, $V = \Spec(B)$라 하자.
$(\Sch/S)_{fppf}$의 구성에 의해 기수가 $\leq |B|$인 모든 환의 스펙트럼은
$(\Sch/S)_{fppf}$의 어떤 대상과 동형이다.
따라서 Algebra, Lemma \ref{algebra-lemma-characterize-finite-presentation}의
``그때에만'' 방향의 증명에서는 $A$-대수 중 기수가 $\leq |B|$인 것들만
사용된다는 점을 관찰하면 충분하다.
\end{proof}

\begin{remark}
\label{spaces-limits-remark-limit-preserving}
Proposition \ref{spaces-limits-proposition-characterize-locally-finite-presentation}의
중요한 특수한 경우는 다음과 같다.
$S$를 스킴이라 하자. $X$를 $S$ 위의 대수공간이라 하자.
그러면 $X$가 $S$ 위에서 국소 유한 표시일 필요충분조건은
$X$가 함자 $(\Sch/S)^{opp} \to \textit{Sets}$로서 극한을 보존하는 것이다.
Limits, Remark \ref{limits-remark-limit-preserving}와 비교하라.
사실 아래의 Lemma \ref{spaces-limits-lemma-surjection-is-enough}에서,
사상
$$
\colim X(T_i) \longrightarrow X(T)
$$
이 $T = \lim T_i$가 아핀 스킴의 $S$ 위의 유향 극한일 때마다
전사인 것만으로 충분함을 보일 것이다.
\end{remark}

\begin{lemma}
\label{spaces-limits-lemma-surjection-is-enough}
$S$를 스킴이라 하자.
$f : X \to Y$를 $S$ 위의 대수공간의 사상이라 하자.
모든 유향 극한 $T = \lim_{i \in I} T_i$가 $S$ 위의 아핀 스킴으로
이루어질 때 사상
$$
\colim X(T_i) \longrightarrow X(T) \times_{Y(T)} \colim Y(T_i)
$$
이 전사이면, $f$는 국소 유한 표시이다.
달리 말하면, Proposition \ref{spaces-limits-proposition-characterize-locally-finite-presentation}의
(2)에서는 Lemma \ref{spaces-limits-lemma-characterize-relative-limit-preserving}의 판정 조건에서
전사성만 확인하면 충분하다.
\end{lemma}

\begin{proof}
스킴 $V$와 전사 에탈 사상 $g : V \to Y$를 택하자.
이어서 스킴 $U$와 전사 에탈 사상 $h : U \to V \times_Y X$를 택하자.
보조정리에서와 같은 $T = \lim T_i$에 대해 사상
$$
\colim U(T_i) \longrightarrow U(T) \times_{V(T)} \colim V(T_i)
$$
이 전사임을 보이면 충분하다. 그러면 $U \to V$가
Limits, Lemma \ref{limits-lemma-surjection-is-enough}에 의해 국소 유한 표시이기 때문이다
(Proposition \ref{spaces-limits-proposition-characterize-locally-finite-presentation}의 증명에서와
동일한 집합론적 논점을 전제한다).
따라서 $a : T \to U$와 $b_i : T_i \to V$를 택하여
같은 사상 $T \to V$를 정하게 하자. 다음 가환도를 보라.
$$
\xymatrix{
T \ar[d]_a \ar[rr]_{p_i} & & T_i \ar[d]^{b_i} \ar@{..>}[ld] \\
U \ar[r]^-h & X \times_Y V \ar[d] \ar[r] & V \ar[d]^g \\
& X \ar[r]^f & Y
}
$$
보조정리의 가정에 의해 $i$를 크게 한 후
사상 $c_i : T_i \to X$를 찾아
$h \circ a = (b_i, c_i) \circ p_i : T_i \to V \times_Y X$이고
$f \circ c_i = g \circ b_i$가 되게 할 수 있다.
% P09-ERR-0058: the source's T_i domain annotation is retained.
$h$는 대수공간의 에탈 사상이므로(따라서 국소 유한 표시이므로),
사상
$$
\colim U(T_i) \longrightarrow U(T) \times_{(X \times_Y V)(T)}
\colim (X \times_Y V)(T_i)
$$
은 Proposition \ref{spaces-limits-proposition-characterize-locally-finite-presentation}에 의해 전사이다.
따라서 $i$를 다시 크게 하면 원하는 사상 $a_i : T_i \to U$를 찾아
$a = a_i \circ p_i$이고 $b_i = (U \to V) \circ a_i$가 되게 할 수 있다.
\end{proof}



\section{대수공간의 극한}
\label{spaces-limits-section-limits}

\noindent
다음 보조정리는 이 장에서 대수공간의 극한을 어떻게 다루는지 설명한다.
대수공간의 아핀 사상을 밑변환하면 아핀 사상이 된다는 사실을
별도의 언급 없이 사용할 것이다(Morphisms of Spaces,
Lemma \ref{spaces-morphisms-lemma-base-change-affine} 참조).

\begin{lemma}
\label{spaces-limits-lemma-directed-inverse-system-has-limit}
$S$를 스킴이라 하자. $I$를 유향 집합이라 하자.
$(X_i, f_{ii'})$를 $I$를 지표로 하는 $S$ 위의 대수공간의 범주의 역계라 하자.
사상 $f_{ii'} : X_i \to X_{i'}$가 아핀이면,
극한 $X = \lim_i X_i$는 (fppf 층으로서 취했을 때) 대수공간이다.
또한 다음이 성립한다.
\begin{enumerate}
\item 각 사상 $f_i : X \to X_i$는 아핀이다.
\item 임의의 $i \in I$와 대수공간의 임의의 사상 $T \to X_i$에 대해
$$
X \times_{X_i} T = \lim_{i' \geq i} X_{i'} \times_{X_i} T.
$$
가 $S$ 위의 대수공간으로서 성립한다.
\end{enumerate}
\end{lemma}

\begin{proof}
(2)는 극한 $X = \lim X_i$가 $S$ 위의 대수공간으로서 존재한다는 사실에서
형식적으로 따른다.
원소 $0 \in I$를 택하자(유향 집합은 공집합이 아니므로 가능하다).
스킴 $U_0$와 전사 에탈 사상 $U_0 \to X_0$를 택하자.
$R_0 = U_0 \times_{X_0} U_0$로 놓으면 $X_0 = U_0/R_0$이다.
$i \geq 0$에 대해 $U_i = X_i \times_{X_0} U_0$ 및
$R_i = X_i \times_{X_0} R_0 = U_i \times_{X_i} U_i$로 놓자.
Limits, Lemma \ref{limits-lemma-directed-inverse-system-has-limit}에 의해
$U = \lim_{i \geq 0} U_i$와 $R = \lim_{i \geq 0} R_i$는 스킴이다.
또한 두 사상 $s, t : R \to U$는 두 사영 $R_0 \to U_0$를
사상 $U \to U_0$로 밑변환한 것이므로, 특히 에탈이다.
사상 $R \to U \times_S U$는 동치관계를 정의한다.
이는 동치관계의 유향 극한이 동치관계이기 때문이다.
% P09-ERR-0059: source grammar transposition; mathematical assertion unchanged.
따라서 사상 $R \to U \times_S U$는 에탈 동치관계이다.
자연스러운 사상
\begin{equation}
\label{spaces-limits-equation-isomorphism-sheaves}
U/R \longrightarrow \lim X_i
\end{equation}
이 $S$ 위의 스킴의 범주 위에서 fppf 층의 동형이라고 주장한다.
이 주장이 성립하면 $X = \lim X_i$는 Spaces,
Theorem \ref{spaces-theorem-presentation}에 의해 대수공간이다.

\medskip\noindent
$Z$를 스킴이라 하고 $a : Z \to \lim X_i$를 사상이라 하자.
그러면 $a = (a_i)$이며, 여기서 $a_i : Z \to X_i$이다.
$W_0 = Z \times_{a_0, X_0} U_0$로 놓자.
$W_0 = Z \times_{a_i, X_i} U_i$가 모든 $i \geq 0$에 대해 성립함에 주의하자.
이는 위에서 $U_i \to X_i$를 택한 방식에 따른다.
따라서 사상 $W_0 \to \lim_{i \geq 0} U_i = U$를 얻는다.
$W_0 \to Z$가 전사 에탈이므로,
(\ref{spaces-limits-equation-isomorphism-sheaves})는 층 사이의 전사 사상이다.
마지막으로 $Z$를 스킴이라 하고, 두 사상 $a, b : Z \to U/R$가
(\ref{spaces-limits-equation-isomorphism-sheaves})와 합성하면 같아진다고 하자.
$a = b$임을 보여야 한다.
$Z$를 fppf 덮개의 원소들로 대체하면,
사상 $a', b' : Z \to U$가 존재하여 각각 $a$와 $b$를 유도한다고 가정할 수 있다.
$a, b$가 (\ref{spaces-limits-equation-isomorphism-sheaves})와 합성하면 같아진다는 조건은,
각 $i \geq 0$에 대해 합성 $a_i', b_i' : Z \to U \to U_i$가
$U_i/R_i = X_i$로 가는 사상으로서 같다는 뜻이다.
따라서 $(a_i', b_i') : Z \to U_i \times_S U_i$는
$R_i$를 경유하며, 이를 주는 사상을 $c_i : Z \to R_i$라 하자.
$R = \lim_{i \geq 0} R_i$이므로 $c = \lim c_i : Z \to R$는 사상이며,
이 사상으로부터 $a, b$가 $Z$에서 $U/R$로 가는 사상으로서 같음을 알 수 있다.

\medskip\noindent
(1)은 위에서 $U_i \times_{X_i} X = U$임을 보였고,
$U \to U_i$가 구성에 의해 아핀이라는 사실에서 따른다.
\end{proof}

\begin{lemma}
\label{spaces-limits-lemma-space-over-limit}
$S$를 스킴이라 하자. $I$를 유향 집합이라 하자.
$(X_i, f_{ii'})$를 $I$를 지표로 하는 $S$ 위의 대수공간의 역계라 하고,
전이 사상들이 아핀이라고 하자.
$X = \lim_i X_i$라 하자. $0 \in I$라 하자.
$T \to X_0$가 대수공간의 사상이라고 가정하자. 그러면
$$
T \times_{X_0} X = \lim_{i \geq 0} T \times_{X_0} X_i
$$
가 $S$ 위의 대수공간으로서 성립한다.
\end{lemma}

\begin{proof}
극한 $X$는 Lemma \ref{spaces-limits-lemma-directed-inverse-system-has-limit}에 의해 대수공간이다.
등식은 형식적이다.
Categories, Lemma \ref{categories-lemma-colimits-commute}를 보라.
\end{proof}

\begin{lemma}
\label{spaces-limits-lemma-directed-inverse-system-closed-immersions}
$S$를 스킴이라 하자. $I$를 유향 집합이라 하자.
$(X_i, f_{i'i}) \to (Y_i, g_{i'i})$를 $I$를 지표로 하는
$S$ 위의 대수공간의 역계 사이의 사상이라 하자.
다음을 가정하자.
\begin{enumerate}
\item 사상 $f_{i'i} : X_{i'} \to X_i$는 아핀이다.
\item 사상 $g_{i'i} : Y_{i'} \to Y_i$는 아핀이다.
\item 사상 $X_i \to Y_i$는 닫힌 몰입이다.
\end{enumerate}
그러면 $\lim X_i \to \lim Y_i$는 닫힌 몰입이다.
\end{lemma}

\begin{proof}
$\lim X_i$와 $\lim Y_i$가
Lemma \ref{spaces-limits-lemma-directed-inverse-system-has-limit}에 의해 존재함을 관찰하자.
$0 \in I$를 택하고 아핀 스킴 $V_0$와 에탈 사상 $V_0 \to Y_0$를 택하자.
그러면 사상
$V_i = Y_i \times_{Y_0} V_0 \to U_i = X_i \times_{Y_0} V_0$는
아핀 스킴의 닫힌 몰입이다.
% P09-ERR-0060: both reversed source arrows retained.
따라서 사상 $V = Y \times_{Y_0} V_0 \to U = X \times_{Y_0} V_0$는 닫힌 몰입이다.
이는 $V = \lim V_i$, $U = \lim U_i$이고,
아핀 스킴의 닫힌 몰입의 극한이 닫힌 몰입이기 때문이다.
실제로 전사 환 준동형의 여과 여극한은 전사이다.
에탈 사상들 $V \to Y$는 $Y$의 에탈 덮개를 이룬다.
여기서 $V_0 \to Y_0$의 선택을 변화시킨다.
따라서 보조정리가 성립한다.
\end{proof}

\begin{lemma}
\label{spaces-limits-lemma-directed-inverse-system-reduced}
$S$를 스킴이라 하자. $I$를 유향 집합이라 하자.
$(X_i, f_{i'i})$를 $I$를 지표로 하는 $S$ 위의 대수공간의 역계라 하자.
$X_i$가 모든 $i$에 대해 축약이면, $X$는 축약이다.
% P09-ERR-0061: the source's omitted affine-transition hypothesis is not inserted.
\end{lemma}

\begin{proof}
$\lim X_i$가 Lemma \ref{spaces-limits-lemma-directed-inverse-system-has-limit}에 의해
존재함을 관찰하자.
$0 \in I$를 택하고 아핀 스킴 $V_0$와 에탈 사상 $U_0 \to X_0$를 택하자.
% P09-ERR-0062: V_0/U_0 discrepancy retained.
그러면 아핀 스킴 $U_i = X_i \times_{X_0} U_0$는 축약이다.
따라서 $U = X \times_{X_0} U_0$는 축약 아핀 스킴의 극한이므로
축약 아핀 스킴이다. 실제로 축약환의 여과 여극한은 축약이다.
에탈 사상들 $U \to X$는 $X$의 에탈 덮개를 이룬다.
여기서 $U_0 \to X_0$의 선택을 변화시킨다.
따라서 보조정리가 성립한다.
\end{proof}

\begin{lemma}
\label{spaces-limits-lemma-better-characterize-relative-limit-preserving}
$S$를 스킴이라 하자. $X \to Y$를 $S$ 위의 대수공간의 사상이라 하자.
Proposition \ref{spaces-limits-proposition-characterize-locally-finite-presentation}의
서로 동치인 조건 (1), (2)는 다음 조건과도 동치이다.
\begin{enumerate}
\item[(3)] 모든 유향 극한 $T = \lim T_i$가 준콤팩트 준분리 대수공간
$T_i$들의 $S$ 위의 역계에서 얻어지고 전이 사상들이 아핀일 때,
집합의 가환도
$$
\xymatrix{
\colim_i \Mor(T_i, X) \ar[r] \ar[d] & \Mor(T, X) \ar[d] \\
\colim_i \Mor(T_i, Y) \ar[r] & \Mor(T, Y)
}
$$
는 올곱 가환도이다.
\end{enumerate}
\end{lemma}

\begin{proof}
(3)이 (2)를 함의함은 분명하다. (2)를 가정하고 (3)을 증명하자.
증명은 상당히 형식적이므로 독자가 직접 증명을 찾아보기를 권한다.

\medskip\noindent
먼저 $T_i$가 모든 $i$에 대해 분리라는 가정을 추가했을 때 (3)을 증명하자.
$i \in I$를 택하고 전사 에탈 사상 $U_i \to T_i$를 택하되,
$U_i$는 아핀이라고 하자.
Lemma \ref{spaces-limits-lemma-space-over-limit}를 사용하면,
$U = U_i \times_{T_i} T$ 및 $U_{i'} = U_i \times_{T_i} T_{i'}$로 놓았을 때
$U = \lim_{i' \geq i} U_{i'}$임을 알 수 있다.
물론 $U$와 $U_{i'}$는 아핀이다
(Lemma \ref{spaces-limits-lemma-directed-inverse-system-has-limit} 참조).
$T_i$가 분리이므로 올곱 $V_i = U_i \times_{T_i} U_i$도 아핀 스킴이다.
따라서 아핀 스킴 $V = V_i \times_{T_i} T$ 및
$V_{i'} = V_i \times_{T_i} T_{i'}$를 얻고,
$V = \lim_{i' \geq i} V_{i'}$가 성립한다.
$U \to T$와 $U_i \to T_i$가 전사 에탈이며,
$V = U \times_T U$ 및 $V_{i'} = U_{i'} \times_{T_{i'}} U_{i'}$임을 관찰하자.
$\Mor(T, X)$는 두 사상 $\Mor(U, X) \to \Mor(V, X)$의 동등자임에 주의하자.
예를 들어, 이는 $X$가 $(\Sch/S)_{fppf}$ 위의 층으로서
두 사상 $h_V \to h_u$의 쌍대동등자이기 때문에 성립한다.
% P09-ERR-0063; P09-ERR-1455: h_u and the subject X retained.
마찬가지로 $\Mor(T_{i'}, X)$는
두 사상 $\Mor(U_{i'}, X) \to \Mor(V_{i'}, X)$의 동등자이다.
물론 $X$를 $Y$로 바꾸어도 같은 사실이 성립한다.
조건 (2)는 (3)의 가환도들이 $U = \lim U_i$와 $V = \lim V_i$의 경우에
올곱이라는 뜻이다. 이로부터 $T = \lim T_i$에 대해서도 같은 사실이 형식적으로 따른다.

\medskip\noindent
일반적인 경우에는 아핀 스킴 $U$, 원소 $i \in I$,
전사 에탈 사상 $U \to T_i$를 택하자.
앞 문단의 논증을 반복하면 여전히 증명이 성립한다.
스킴 $V_{i'}$, $V$는 더 이상 아핀일 필요는 없지만,
여전히 준콤팩트이고 분리이므로 앞 문단의 결과를 적용할 수 있다.
\end{proof}



\section{성질의 하강}
\label{spaces-limits-section-descent}

\noindent
이 절은 Limits, Section \ref{limits-section-descent}에 대응하는 내용을 다룬다.

\begin{lemma}
\label{spaces-limits-lemma-inverse-limit-sets}
$S$를 스킴이라 하자.
$X = \lim_{i \in I} X_i$를 $S$ 위의 대수공간의 유향 역계의 극한이라 하고,
전이 사상들이 아핀이라고 하자
(Lemma \ref{spaces-limits-lemma-directed-inverse-system-has-limit}).
각 $X_i$가 decent이면(예를 들어 준분리이거나 국소 분리이면),
집합으로서 $|X| = \lim_i |X_i|$이다.
\end{lemma}

\begin{proof}
표준 사상 $|X| \to \lim |X_i|$가 있다. $0 \in I$를 택하자.
$W_0 \subset X_0$가 열린 부분공간이면,
$f_0^{-1}W_0 = \lim_{i \geq 0} f_{i0}^{-1}W_0$이다.
Lemma \ref{spaces-limits-lemma-directed-inverse-system-has-limit}를 보라.
따라서 $X_0$가 준콤팩트인 역계에 대해 보조정리를 증명할 수 있으면,
일반적인 경우에도 보조정리가 따른다.
그러므로 $X_0$가 준콤팩트라고 가정해도 되며, 이제 그렇게 가정한다.

\medskip\noindent
아핀 스킴 $U_0$와 전사 에탈 사상 $U_0 \to X_0$를 택하자.
$U_i = X_i \times_{X_0} U_0$ 및 $U = X \times_{X_0} U_0$로 놓자.
$R_i = U_i \times_{X_i} U_i$ 및 $R = U \times_X U$로 놓자.
$U = \lim U_i$ 및 $R = \lim R_i$임을 상기하자.
Lemma \ref{spaces-limits-lemma-directed-inverse-system-has-limit}의 증명을 보라.
$|X| = |U|/|R|$ 및 $|X_i| = |U_i|/|R_i|$임을 상기하자.
Limits, Lemma \ref{limits-lemma-topology-limit}에 의해
$|U| = \lim |U_i|$ 및 $|R| = \lim |R_i|$이다.

\medskip\noindent
$|X| \to \lim |X_i|$의 전사성.
$(x_i) \in \lim |X_i|$라 하자.
$S_i \subset |U_i|$로 $x_i$의 역상을 나타내자.
이는 decent 공간의 정의에 의해 공집합이 아닌 유한 집합이다
(Decent Spaces, Definition \ref{decent-spaces-definition-very-reasonable}).
따라서 $\lim S_i$는 공집합이 아니다.
Categories, Lemma \ref{categories-lemma-nonempty-limit}를 보라.
$(u_i) \in \lim S_i \subset \lim |U_i|$라 하자.
위의 사실에 의해 이는 점 $u \in |U|$를 정한다.
그 상 $x \in |X|$는 주어진 원소 $(x_i)$로 $\lim |X_i|$ 안에서 간다.

\medskip\noindent
$|X| \to \lim |X_i|$의 단사성.
$x, x' \in |X|$가 $\lim |X_i|$의 같은 점으로 간다고 하자.
올림 $u, u' \in |U|$를 택하고 그 상을 $u_i, u'_i \in |U_i|$라 쓰자.
각 $i$에 대해 $T_i \subset |R_i|$를
$(u_i, u'_i) \in |U_i| \times |U_i|$로 가는 점들의 집합이라 하자.
이는 decent 공간의 정의에 의해 유한 집합이다
(Decent Spaces, Definition \ref{decent-spaces-definition-very-reasonable}).
또한 $T_i$는 공집합이 아니다.
이는 $x$와 $x'$가 $X_i$의 같은 점으로 간다고 가정했기 때문이다.
따라서 $\lim T_i$는 공집합이 아니다.
Categories, Lemma \ref{categories-lemma-nonempty-limit}를 보라.
앞에서와 같이 $r \in |R| = \lim |R_i|$를 $\lim T_i$의 원소에 대응하는 점이라 하자.
그러면 구성에 의해 $r$는 $(u, u')$로 $|U| \times |U|$ 안에서 간다.
따라서 원하는 대로 $x = x'$가 $|X|$ 안에서 성립한다.

\medskip\noindent
괄호 안의 주장: 준분리 대수공간은 decent이다.
Decent Spaces, Section \ref{decent-spaces-section-reasonable-decent}를 보라
(여기서 핵심적인 관찰은 Properties of Spaces,
Lemma \ref{spaces-properties-lemma-finite-fibres-presentation}이다).
국소 분리 대수공간은 Decent Spaces,
Lemma \ref{decent-spaces-lemma-locally-separated-decent}에 의해 decent이다.
\end{proof}

\begin{lemma}
\label{spaces-limits-lemma-topology-limit}
Lemma \ref{spaces-limits-lemma-inverse-limit-sets}와 같은 표기와 가정하에서,
위상공간으로서 $|X| = \lim_i |X_i|$가 성립한다.
\end{lemma}

\begin{proof}
Topology, Lemma \ref{topology-lemma-characterize-limit}의 판정법을 사용할 것이다.
집합으로서 $|X| = \lim_i |X_i|$임은
Lemma \ref{spaces-limits-lemma-inverse-limit-sets}에서 보았다.
사상 $f_i : X \to X_i$는 대수공간의 사상이므로
연속 사상 $|X| \to |X_i|$를 정한다.
따라서 $f_i^{-1}(U_i)$는 각 열린집합 $U_i \subset |X_i|$에 대해 열린집합이다.
마지막으로 $x \in |X|$라 하고 $x \in V \subset |X|$를 열린 근방이라 하자.
$i$와 열린 근방 $W_i \subset |X_i|$를 상 $x$에 대해 찾아서
$f_i^{-1}(W_i) \subset V$가 되게 해야 한다.
$0 \in I$를 택하자. 스킴 $U_0$와 전사 에탈 사상 $U_0 \to X_0$를 택하자.
$U = X \times_{X_0} U_0$ 및 $U_i = X_i \times_{X_0} U_0$로 $i \geq 0$에 대해 놓자.
그러면 $U = \lim_{i \geq 0} U_i$가 스킴의 범주에서 성립한다.
이는 Lemma \ref{spaces-limits-lemma-directed-inverse-system-has-limit}에 따른다.
$u \in U$를 택하여 $x$로 가게 하자.
스킴에 대한 결과(Limits, Lemma \ref{limits-lemma-inverse-limit-top})에 의해
$i \geq 0$와 열린 근방 $E_i \subset U_i$를 $u$의 상에 대해 찾아서,
그 $U$ 안에서의 역상이 $V$의 $U$ 안에서의 역상에 포함되게 할 수 있다.
그러면 $W_i \subset |X_i|$를 $E_i$의 상으로 놓으면 된다.
이렇게 할 수 있는 이유는 $|U_i| \to |X_i|$가 열린 사상이기 때문이다.
\end{proof}

\begin{lemma}
\label{spaces-limits-lemma-limit-nonempty}
$S$를 스킴이라 하자.
$X = \lim_{i \in I} X_i$를 $S$ 위의 대수공간의 유향 역계의 극한이라 하고,
전이 사상들이 아핀이라고 하자
(Lemma \ref{spaces-limits-lemma-directed-inverse-system-has-limit}).
각 $X_i$가 준콤팩트이고 공집합이 아니면 $|X|$는 공집합이 아니다.
\end{lemma}

\begin{proof}
$0 \in I$를 택하자.
아핀 스킴 $U_0$와 전사 에탈 사상 $U_0 \to X_0$를 택하자.
$U_i = X_i \times_{X_0} U_0$ 및 $U = X \times_{X_0} U_0$로 놓자.
그러면 각 $U_i$는 공집합이 아닌 아핀 스킴이다.
따라서 $U = \lim U_i$는 공집합이 아니고
(Limits, Lemma \ref{limits-lemma-limit-nonempty}),
그러므로 $X$도 공집합이 아니다.
\end{proof}

\begin{lemma}
\label{spaces-limits-lemma-inverse-limit-irreducibles}
$S$를 스킴이라 하자.
$X = \lim_{i \in I} X_i$를 $S$ 위의 대수공간의 유향 역계의 극한이라 하고,
전이 사상들이 아핀이라고 하자
(Lemma \ref{spaces-limits-lemma-directed-inverse-system-has-limit}).
$x \in |X|$의 상을 $x_i \in |X_i|$라 하자.
각 $X_i$가 decent이면 $\overline{\{x\}} = \lim_i \overline{\{x_i\}}$가
집합으로서 성립하며, 축약 유도 스킴 구조를 주면 대수공간으로서도 성립한다.
\end{lemma}

\begin{proof}
$Z = \overline{\{x\}} \subset |X|$ 및
$Z_i = \overline{\{x_i\}} \subset |X_i|$로 놓자.
$|X| \to |X_i|$가 연속이므로 $Z$는 $Z_i$ 안으로 각 $i$에 대해 간다.
따라서 단사 사상 $Z \to \lim Z_i$를 얻는다.
이는 집합으로서 $|X| = \lim |X_i|$이기 때문이다
(Lemma \ref{spaces-limits-lemma-inverse-limit-sets}).
$x' \in |X|$가 $Z$에 속하지 않는다고 하자.
그러면 열린 부분집합 $U \subset |X|$가 존재하여
$x' \in U$이고 $x \not \in U$이다.
위상공간으로서 $|X| = \lim |X_i|$이므로
(Lemma \ref{spaces-limits-lemma-topology-limit}),
$U = \bigcup_{j \in J} f_j^{-1}(U_j)$로 쓸 수 있다.
여기서 $J \subset I$는 어떤 부분집합이고 $U_j \subset |X_j|$는 열린집합이다.
Topology, Lemma \ref{topology-lemma-describe-limits}를 보라.
그러면 어떤 $j \in J$에 대해 $f_j(x') \in U_j$이고
$f_j(x) \not \in U_j$임을 알 수 있다.
달리 말하면 $f_j(x') \not \in Z_j$이다.
따라서 집합으로서 $Z = \lim Z_i$이다.

\medskip\noindent
이어서 $Z$와 $Z_i$에 축약 유도 스킴 구조를 주자.
Properties of Spaces, Definition
\ref{spaces-properties-definition-reduced-induced-space}를 보라.
전이 사상 $X_{i'} \to X_i$는 아핀 사상 $Z_{i'} \to Z_i$를 유도하고,
사영 $X \to X_i$는 서로 호환되는 사상 $Z \to Z_i$를 유도한다.
따라서 대수공간의 사상 $Z \to \lim Z_i \to X$를 얻는다.
Lemma \ref{spaces-limits-lemma-directed-inverse-system-closed-immersions}에 의해
$\lim Z_i \to X$는 닫힌 몰입이다.
Lemma \ref{spaces-limits-lemma-directed-inverse-system-reduced}에 의해
대수공간 $\lim Z_i$는 축약이다.
위에서 보인 바와 같이 $Z \to \lim Z_i$는 점들 사이에서 전단사이다.
축약 유도 닫힌 부분스킴 구조의 유일성에 의해
이 사상은 대수공간의 동형이다.
\end{proof}

\begin{situation}
\label{spaces-limits-situation-descent}
$S$를 스킴이라 하자.
$X = \lim_{i \in I} X_i$를 $S$ 위의 대수공간의 유향 역계의 극한이라 하고,
전이 사상들이 아핀이라고 하자
(Lemma \ref{spaces-limits-lemma-directed-inverse-system-has-limit}).
$X_i$가 모든 $i \in I$에 대해 준콤팩트이고 준분리라고 가정한다.
원소 $0 \in I$도 하나 택한다.
\end{situation}

\begin{lemma}
\label{spaces-limits-lemma-descend-section}
표기와 가정은 Situation \ref{spaces-limits-situation-descent}와 같다.
$\mathcal{F}_0$가 $X_0$ 위의 준연접층이라고 가정하자.
$\mathcal{F}_i = f_{0i}^*\mathcal{F}_0$로 $i \geq 0$에 대해 놓고,
$\mathcal{F} = f_0^*\mathcal{F}_0$로 놓자. 그러면
$$
\Gamma(X, \mathcal{F}) = \colim_{i \geq 0} \Gamma(X_i, \mathcal{F}_i)
$$
이다.
\end{lemma}

\begin{proof}
전사 에탈 사상 $U_0 \to X_0$를 택하되 $U_0$는 아핀 스킴이라고 하자
(Properties of Spaces, Lemma
\ref{spaces-properties-lemma-quasi-compact-affine-cover}).
$U_i = X_i \times_{X_0} U_0$로 놓자.
$R_0 = U_0 \times_{X_0} U_0$ 및 $R_i = R_0 \times_{X_0} X_i$로 놓자.
Lemma \ref{spaces-limits-lemma-directed-inverse-system-has-limit}의 증명에서
표시 $X = U/R$가 존재하고 $U = \lim U_i$, $R = \lim R_i$임을 보았다.
$U_i$와 $U$는 아핀이며, $R_i$와 $R$은 준콤팩트이고 분리임에 주의하자
($X_i$가 준분리이기 때문이다).
따라서 Limits, Lemma \ref{limits-lemma-descend-section}에 의해
$$
\mathcal{F}(U) = \colim \mathcal{F}_i(U_i)
\quad\text{및}\quad
\mathcal{F}(R) = \colim \mathcal{F}_i(R_i).
$$
이다. 보조정리는
$\Gamma(X, \mathcal{F}) = \Ker(\mathcal{F}(U) \to \mathcal{F}(R))$이고,
마찬가지로
$\Gamma(X_i, \mathcal{F}_i) =
\Ker(\mathcal{F}_i(U_i) \to \mathcal{F}_i(R_i))$라는 사실에서 따른다.
\end{proof}

\begin{lemma}
\label{spaces-limits-lemma-descend-opens}
표기와 가정은 Situation \ref{spaces-limits-situation-descent}와 같다.
임의의 준콤팩트 열린 부분공간 $U \subset X$에 대해 $i$와
준콤팩트 열린집합 $U_i \subset X_i$가 존재하여,
그 $X$ 안에서의 역상이 $U$가 된다.
\end{lemma}

\begin{proof}
Lemma \ref{spaces-limits-lemma-directed-inverse-system-has-limit}에서의 극한의 구성과
스킴에 대한 대응하는 결과인
Limits, Lemma \ref{limits-lemma-descend-opens}에서 형식적으로 따른다.
\end{proof}

\noindent
다음 보조정리는 뒤에 나오는 더 강한
Lemma \ref{spaces-limits-lemma-descend-isomorphism}으로 대체될 것이다.

\begin{lemma}
\label{spaces-limits-lemma-descend-equality}
표기와 가정은 Situation \ref{spaces-limits-situation-descent}와 같다.
$f_0 : Y_0 \to Z_0$를 $X_0$ 위의 대수공간의 사상이라 하자.
다음을 가정하자. (a) $Y_0 \to X_0$와 $Z_0 \to X_0$는 표현 가능하다.
(b) $Y_0$, $Z_0$는 준콤팩트이고 준분리이다.
(c) $f_0$는 국소 유한 표시이다.
(d) $Y_0 \times_{X_0} X \to Z_0 \times_{X_0} X$는 동형이다.
그러면 $i \geq 0$가 존재하여
$Y_0 \times_{X_0} X_i \to Z_0 \times_{X_0} X_i$가 동형이다.
\end{lemma}

\begin{proof}
아핀 스킴 $U_0$와 전사 에탈 사상 $U_0 \to X_0$를 택하자.
$U_i = U_0 \times_{X_0} X_i$ 및 $U = U_0 \times_{X_0} X$로 놓자.
Limits, Lemma \ref{limits-lemma-descend-isomorphism}을 적용하면
$Y_0 \times_{X_0} U_i \to Z_0 \times_{X_0} U_i$가
어떤 $i \geq 0$에 대해 스킴의 동형임을 알 수 있다(세부 사항은 생략한다).
$U_i \to X_i$가 전사 에탈이므로,
$Y_0 \times_{X_0} X_i \to Z_0 \times_{X_0} X_i$는 동형이다
(세부 사항은 생략한다).
\end{proof}

\begin{lemma}
\label{spaces-limits-lemma-descend-separated}
표기와 가정은 Situation \ref{spaces-limits-situation-descent}와 같다.
$X$가 분리이면 $X_i$는 어떤 $i \in I$에 대해 분리이다.
\end{lemma}

\begin{proof}
아핀 스킴 $U_0$와 전사 에탈 사상 $U_0 \to X_0$를 택하자.
$i \geq 0$에 대해 $U_i = U_0 \times_{X_0} X_i$로 놓고,
$U = U_0 \times_{X_0} X$로 놓자.
$U_i$와 $U$는 전사 에탈 사상 $U_i \to X_i$ 및 $U \to X$를 갖춘
아핀 스킴임에 주의하자.
$R_i = U_i \times_{X_i} U_i$ 및 $R = U \times_X U$로 놓고,
그 사영을 $s_i, t_i : R_i \to U_i$ 및 $s, t : R \to U$라 하자.
$R_i$와 $R$은 준콤팩트 분리 스킴이다
(대수공간 $X_i$와 $X$가 준분리이기 때문이다).
사상 $s_i : R_i \to U_i$ 및 $s : R \to U$는 유한형이다.
정의에 의해 $X_i$가 분리일 필요충분조건은
$(t_i, s_i) : R_i \to U_i \times U_i$가 닫힌 몰입인 것이다.
가정에 의해 $X$가 분리이므로,
사상 $(t, s) : R \to U \times U$는 닫힌 몰입이다.
$R \to U$가 유한형이므로 $i$가 존재하여
사상 $R \to U_i \times U$가 닫힌 몰입이다
(Limits, Lemma \ref{limits-lemma-finite-type-eventually-closed}).
그러한 $i \in I$를 고정하자.
Limits, Lemma \ref{limits-lemma-descend-closed-immersion-finite-presentation}을
사상의 계 $R_{i'} \to U_i \times U_{i'}$에 $i' \geq i$에 대해 적용하자
(실제로 $R_{i'} = R_i \times_{U_i \times U_i} U_i \times U_{i'}$이므로 적용할 수 있다).
그러면 $R_{i'} \to U_i \times U_{i'}$가 충분히 큰 $i'$에 대해 닫힌 몰입이다.
이로부터 $R_{i'} \to U_{i'} \times U_{i'}$도 닫힌 몰입임이 즉시 따라,
보조정리의 증명이 끝난다.
\end{proof}

\begin{lemma}
\label{spaces-limits-lemma-limit-is-affine}
표기와 가정은 Situation \ref{spaces-limits-situation-descent}와 같다.
$X$가 아핀이면 $i$가 존재하여 $X_i$가 아핀이다.
\end{lemma}

\begin{proof}
$0 \in I$를 택하자. 아핀 스킴 $U_0$와 전사 에탈 사상 $U_0 \to X_0$를 택하자.
$U = U_0 \times_{X_0} X$ 및 $U_i = U_0 \times_{X_0} X_i$로
$i \geq 0$에 대해 놓자.
전이 사상들이 아핀이므로 대수공간 $U_i$와 $U$는 아핀이다.
따라서 $U \to X$는 아핀 스킴의 에탈 사상이다.
그러므로 $X = \Spec(A)$, $U = \Spec(B)$ 및
$$
B = A[x_1, \ldots, x_n]/(g_1, \ldots, g_n)
$$
으로 쓰되, $\Delta = \det(\partial g_\lambda/\partial x_\mu)$가
$B$에서 가역이게 할 수 있다.
Algebra, Lemma \ref{algebra-lemma-etale-standard-smooth}를 보라.
$A_i = \mathcal{O}_{X_i}(X_i)$로 놓자.
$A = \colim A_i$가 Lemma \ref{spaces-limits-lemma-descend-section}에 의해 성립한다.
$0$을 크게 한 후, $g_{1, i}, \ldots, g_{n, i} \in A_i[x_1, \ldots, x_n]$가
존재하여 $g_1, \ldots, g_n$으로 간다고 가정할 수 있다.
$$
B_i = A_i[x_1, \ldots, x_n]/(g_{1, i}, \ldots, g_{n, i})
$$
로 모든 $i \geq 0$에 대해 놓자.
필요하면 $0$을 크게 하여
$\Delta_i = \det(\partial g_{\lambda, i}/\partial x_\mu)$가
$B_i$에서 모든 $i \geq 0$에 대해 가역이라고 가정할 수 있다.
따라서 $A_i \to B_i$는 에탈 환 준동형이다.
$0$을 크게 한 후에는 $\Spec(B_i) \to \Spec(A_i)$가 전사라고도 가정할 수 있다.
Limits, Lemma \ref{limits-lemma-descend-surjective}를 보라.
$0$을 다시 크게 하면 원소
$h_{1, i}, \ldots, h_{n, i} \in \mathcal{O}_{U_i}(U_i)$를 택하여
$x_1, \ldots, x_n$의 $B = \mathcal{O}_U(U)$ 안에서의 류로 가고,
$g_{\lambda, i}(h_{\nu, i}) = 0$이 $\mathcal{O}_{U_i}(U_i)$에서 성립하게 할 수 있다.
따라서 다음 가환도를 얻는다.
\begin{equation}
\label{spaces-limits-equation-to-show-cartesian}
\vcenter{
\xymatrix{
X_i \ar[d] & U_i \ar[l] \ar[d] \\
\Spec(A_i) & \Spec(B_i) \ar[l]
}
}
\end{equation}
구성에 의해 $B_i = B_0 \otimes_{A_0} A_i$이고
$B = B_0 \otimes_{A_0} A$이다. 사상
$$
f_0 : U_0 \longrightarrow X_0 \times_{\Spec(A_0)} \Spec(B_0)
$$
를 생각하자. 이는 준콤팩트 준분리 대수공간 사이의 사상이며,
그 대수공간들은 $X_0$ 위에서 표현 가능하고 분리이며 에탈이다.
$f_0$를 $X$로 밑변환하면 선택에 의해 동형이 된다.
따라서 Lemma \ref{spaces-limits-lemma-descend-equality}에 의해
$i$가 존재하여 $f_0$를 $X_i$로 밑변환한 것이 동형이 된다.
달리 말하면 가환도 (\ref{spaces-limits-equation-to-show-cartesian})는 데카르트 가환도이다.
따라서 Descent, Lemma \ref{descent-lemma-descent-data-sheaves}를
fppf 덮개 $\{\Spec(B_i) \to \Spec(A_i)\}$에 적용하고
Descent, Lemma \ref{descent-lemma-affine}와 결합하면,
$X_i \to \Spec(A_i)$는 $\Spec(A_i)$ 위에서 아핀인 스킴으로 표현 가능하다.
이것이 원하는 결과이다.
(물론 여기서 $X_i = \Spec(A_i)$도 따르지만, 이는 필요하지 않다.)
\end{proof}

\begin{lemma}
\label{spaces-limits-lemma-limit-is-scheme}
표기와 가정은 Situation \ref{spaces-limits-situation-descent}와 같다.
$X$가 스킴이면 $i$가 존재하여 $X_i$가 스킴이다.
\end{lemma}

\begin{proof}
유한 아핀 열린 덮개 $X = \bigcup W_j$를 택하자.
Lemma \ref{spaces-limits-lemma-descend-opens}에 의해 $i \in I$와
열린 부분공간 $W_{j, i} \subset X_i$를 찾아,
그 $X$로의 밑변환이 $W_j \to X$가 되게 할 수 있다.
Lemma \ref{spaces-limits-lemma-limit-is-affine}에 의해
각 $W_{j, i}$가 아핀 스킴이라고 가정할 수 있다.
이는 $X_i$가 스킴이라는 뜻이다
(예를 들어 Properties of Spaces, Section
\ref{spaces-properties-section-schematic}을 보라).
\end{proof}

\begin{lemma}
\label{spaces-limits-lemma-finite-type-eventually-closed}
$S$를 스킴이라 하자. $B$를 $S$ 위의 대수공간이라 하자.
$X = \lim X_i$를 대수공간의 $B$ 위의 유향 극한이라 하고,
전이 사상들이 아핀이라고 하자.
$Y \to X$를 $B$ 위의 대수공간의 사상이라 하자.
\begin{enumerate}
\item $Y \to X$가 닫힌 몰입이고, $X_i$가 준콤팩트이며,
$Y \to B$가 국소 유한형이면,
$Y \to X_i$는 충분히 큰 $i$에 대해 닫힌 몰입이다.
\item $Y \to X$가 몰입이고, $X_i$가 준분리이며,
$Y \to B$가 국소 유한형이고 $Y$가 준콤팩트이면,
$Y \to X_i$는 충분히 큰 $i$에 대해 몰입이다.
\item $Y \to X$가 동형이고, $X_i$가 준콤팩트이며,
$X_i \to B$가 국소 유한형이고, 전이 사상 $X_{i'} \to X_i$가 닫힌 몰입이며,
$Y \to B$가 국소 유한 표시이면,
$Y \to X_i$는 충분히 큰 $i$에 대해 동형이다.
\item $Y \to X$가 단사상이고, $X_i$가 준분리이며,
$Y \to B$가 국소 유한형이고 $Y$가 준콤팩트이면,
$Y \to X_i$는 충분히 큰 $i$에 대해 단사상이다.
\end{enumerate}
\end{lemma}

\begin{proof}
(1)의 증명. $0 \in I$를 택하자.
$X_0$가 준콤팩트이므로 아핀 스킴 $W$와 에탈 사상 $W \to B$를 택하여
$|X_0| \to |B|$의 상이 $|W| \to |B|$의 상에 포함되게 할 수 있다.
아핀 스킴 $U_0$와 에탈 사상 $U_0 \to X_0 \times_B W$를 택하되,
$U_0 \to X_0$가 전사이게 하자.
(이는 $W$를 택한 방식과 $X_0$가 준콤팩트라는 사실에 의해 가능하다.
세부 사항은 생략한다.)
$V \to Y$, $U \to X$, $U_i \to X_i$를 각각
$U_0 \to X_0$의 밑변환이라 하자($i \geq 0$에 대해).
$V \to U_i$가 충분히 큰 $i$에 대해 닫힌 몰입임을 보이면 충분하다.
따라서 $V \to U = \lim U_i$에 대한 결과를 $W$ 위에서 증명하는 것으로 환원된다.
이는 스킴의 경우, 즉 Limits, Lemma
\ref{limits-lemma-finite-type-eventually-closed}에서 따른다.

\medskip\noindent
(2)의 증명. $0 \in I$를 택하자.
준콤팩트 열린 부분공간 $X'_0 \subset X_0$를 택하여
$Y \to X_0$가 $X'_0$를 경유하게 하자.
$X_i$를 $X'_0$의 역상으로 $i \geq 0$에 대해 대체하면,
모든 $X_i'$가 준콤팩트이고 준분리라고 가정할 수 있다.
$U \subset X$를 준콤팩트 열린집합으로 택하여
$Y \to X$가 닫힌 몰입 $Y \to U$를 경유하게 하자
($U$는 $Y$가 준콤팩트이므로 존재한다).
Lemma \ref{spaces-limits-lemma-descend-opens}에 의해
$U = \lim U_i$이고 $U_i \subset X_i$가 준콤팩트 열린집합이라고 가정할 수 있다.
(1)에 의해 $Y \to U_i$는 어떤 $i$에 대해 닫힌 몰입이다.
따라서 (2)가 성립한다.

\medskip\noindent
(3)의 증명. $0 \in I$를 택하자.
아핀 스킴 $U_0$와 전사 에탈 사상 $U_0 \to X_0$를 택하자.
$U_i = X_i \times_{X_0} U_0$,
$U = X \times_{X_0} U_0 = Y \times_{X_0} U_0$로 놓자.
그러면 $U = \lim U_i$는 아핀 스킴의 극한이고,
이 계의 전이 사상들은 닫힌 몰입이며, $U \to U_0$는 유한 표시이다
($U \to B$는 국소 유한 표시이고 $U_0 \to B$는 국소 유한형이므로
Morphisms of Spaces,
Lemma \ref{spaces-morphisms-lemma-finite-presentation-permanence}를 적용한다).
이로써 다음 대수적 사실로 환원된다.
$A = \lim A_i$가 $R$-대수의 유향 여극한이고 전이 준동형들이 전사이며,
$A$가 $A_0$ 위에서 유한 표시이면, $A = A_i$가 어떤 $i$에 대해 성립한다.
실제로 $A = A_0/(f_1, \ldots, f_n)$으로 쓰자.
$i$를 택하여 $f_1, \ldots, f_n$이 전사 준동형 $A_0 \to A_i$ 아래에서
영으로 가게 하면 된다.

\medskip\noindent
(4)의 증명. $Z_i = Y \times_{X_i} Y$로 놓자.
전이 사상 $X_{i'} \to X_i$는 아핀이므로 분리이다.
따라서 전이 사상 $Z_{i'} \to Z_i$는 닫힌 몰입이다.
Morphisms of Spaces, Lemma
\ref{spaces-morphisms-lemma-fibre-product-after-map}을 보라.
$\lim Z_i = Y \times_X Y = Y$이다.
이는 $Y \to X$가 단사상이기 때문이다.
$0 \in I$를 택하자. $Y \to X_0$가 국소 유한형이므로
(Morphisms of Spaces, Lemma \ref{spaces-morphisms-lemma-permanence-finite-type}),
사상 $Y \to Z_0$는 국소 유한 표시이다
(Morphisms of Spaces, Lemma
\ref{spaces-morphisms-lemma-diagonal-morphism-finite-type}).
사상 $Z_i \to Z_0$는 국소 유한형이다(닫힌 몰입이기 때문이다).
마지막으로 $Z_i = Y \times_{X_i} Y$는 준콤팩트이다.
이는 $X_i$가 준분리이고 $Y$가 준콤팩트이기 때문이다.
따라서 (3)을 $Y = \lim_{i \geq 0} Z_i$에 $Z_0$ 위에서 적용할 수 있고,
$Y = Z_i$가 어떤 $i$에 대해 성립함을 얻는다.
이로써 (4)와 보조정리가 증명된다.
\end{proof}

\begin{lemma}
\label{spaces-limits-lemma-eventually-separated}
$S$를 스킴이라 하자. $Y$를 $S$ 위의 대수공간이라 하자.
$X = \lim X_i$를 대수공간의 $Y$ 위의 유향 극한이라 하고,
전이 사상들이 아핀이라고 하자. 다음을 가정하자.
\begin{enumerate}
\item $Y$는 준분리이다.
\item $X_i$는 준콤팩트이고 준분리이다.
\item 사상 $X \to Y$는 분리이다.
\end{enumerate}
그러면 $X_i \to Y$는 충분히 큰 모든 $i$에 대해 분리이다.
\end{lemma}

\begin{proof}
$0 \in I$라 하자. 아핀 스킴 $W$와 에탈 사상 $W \to Y$를 택하여
$|W| \to |Y|$의 상이 $|X_0| \to |Y|$의 상을 포함하게 하자.
이는 $X_0$가 준콤팩트이므로 가능하다.
$W \times_Y X_i \to W$가 어떤 $i \geq 0$에 대해 분리임을 확인하면 충분하다.
이는 $W \times_Y X_i$의 $W$ 위의 대각 사상이
$X_i \to X_i \times_Y X_i$를 전사 에탈 사상
$(X_i \times_Y X_i) \times_Y W \to X_i \times_Y X_i$로 밑변환한 것이기 때문이다.
$Y$가 준분리이므로 대수공간 $W \times_Y X_i$는 준콤팩트이다
(또한 준분리이다).
따라서 $W$로 밑변환하여 $Y$가 아핀 스킴이라고 가정할 수 있다.
$Y$가 아핀 스킴이면, $X_i$가 충분히 큰 $i$에 대해 분리 대수공간임을
보여야 하며, $X$가 분리 대수공간이라는 가정이 주어져 있다.
따라서 이 경우는 Lemma \ref{spaces-limits-lemma-descend-separated}에서 따른다.
\end{proof}

\begin{lemma}
\label{spaces-limits-lemma-eventually-affine}
$S$를 스킴이라 하자. $Y$를 $S$ 위의 대수공간이라 하자.
$X = \lim X_i$를 대수공간의 $Y$ 위의 유향 극한이라 하고,
전이 사상들이 아핀이라고 하자. 다음을 가정하자.
\begin{enumerate}
\item $Y$는 준콤팩트이고 준분리이다.
\item $X_i$는 준콤팩트이고 준분리이다.
\item $X \to Y$는 아핀이다.
\end{enumerate}
그러면 $X_i \to Y$는 충분히 큰 $i$에 대해 아핀이다.
\end{lemma}

\begin{proof}
아핀 스킴 $W$와 전사 에탈 사상 $W \to Y$를 택하자.
그러면 $X \times_Y W$는 아핀이고,
$X_i \times_Y W$가 어떤 $i$에 대해 아핀임을 확인하면 충분하다
(Morphisms of Spaces, Lemma \ref{spaces-morphisms-lemma-affine-local}).
이는 Lemma \ref{spaces-limits-lemma-limit-is-affine}에서 따른다.
\end{proof}

\begin{lemma}
\label{spaces-limits-lemma-eventually-finite}
$S$를 스킴이라 하자. $Y$를 $S$ 위의 대수공간이라 하자.
$X = \lim X_i$를 대수공간의 $Y$ 위의 유향 극한이라 하고,
전이 사상들이 아핀이라고 하자. 다음을 가정하자.
\begin{enumerate}
\item $Y$는 준콤팩트이고 준분리이다.
\item $X_i$는 준콤팩트이고 준분리이다.
\item 전이 사상 $X_{i'} \to X_i$는 유한이다.
\item $X_i \to Y$는 국소 유한형이다.
\item $X \to Y$는 정수적이다.
\end{enumerate}
그러면 $X_i \to Y$는 충분히 큰 $i$에 대해 유한이다.
\end{lemma}

\begin{proof}
아핀 스킴 $W$와 전사 에탈 사상 $W \to Y$를 택하자.
그러면 $X \times_Y W$는 $W$ 위에서 유한이고,
% P09-ERR-1456: the canonical proof says finite, not integral, here.
$X_i \times_Y W$가 $W$ 위에서 어떤 $i$에 대해 유한임을 확인하면 충분하다
(Morphisms of Spaces, Lemma \ref{spaces-morphisms-lemma-integral-local}).
Lemma \ref{spaces-limits-lemma-limit-is-scheme}에 의해 스킴의 경우로 환원된다.
스킴의 경우에는 Limits, Lemma \ref{limits-lemma-eventually-finite}에서 따른다.
\end{proof}

\begin{lemma}
\label{spaces-limits-lemma-eventually-closed-immersion}
$S$를 스킴이라 하자. $Y$를 $S$ 위의 대수공간이라 하자.
$X = \lim X_i$를 대수공간의 $Y$ 위의 유향 극한이라 하고,
전이 사상들이 아핀이라고 하자. 다음을 가정하자.
\begin{enumerate}
\item $Y$는 준콤팩트이고 준분리이다.
\item $X_i$는 준콤팩트이고 준분리이다.
\item 전이 사상 $X_{i'} \to X_i$는 닫힌 몰입이다.
\item $X_i \to Y$는 국소 유한형이다.
\item $X \to Y$는 닫힌 몰입이다.
\end{enumerate}
그러면 $X_i \to Y$는 충분히 큰 $i$에 대해 닫힌 몰입이다.
\end{lemma}

\begin{proof}
아핀 스킴 $W$와 전사 에탈 사상 $W \to Y$를 택하자.
그러면 $X \times_Y W$는 $W$의 닫힌 부분공간이고,
$X_i \times_Y W$가 $W$의 닫힌 부분공간이 어떤 $i$에 대해 됨을 확인하면 충분하다
(Morphisms of Spaces, Lemma \ref{spaces-morphisms-lemma-closed-immersion-local}).
Lemma \ref{spaces-limits-lemma-limit-is-scheme}에 의해 스킴의 경우로 환원된다.
스킴의 경우에는 Limits, Lemma \ref{limits-lemma-eventually-closed-immersion}에서 따른다.
\end{proof}



\section{사상의 성질의 하강}
\label{spaces-limits-section-descent-of-properties}

\noindent
이 절에서는 Section \ref{spaces-limits-section-descent}에 대응하는 결과를
사상의 성질에 대해 다룬다. 다음 상황에서 논의할 것이다.

\begin{situation}
\label{spaces-limits-situation-descent-property}
$S$를 스킴이라 하자. $B = \lim B_i$를 유향 역계의 극한이라 하고,
그 역계는 $S$ 위의 대수공간들로 이루어지며 전이 사상들이 아핀이라고 하자
(Lemma \ref{spaces-limits-lemma-directed-inverse-system-has-limit}).
$0 \in I$라 하고 $f_0 : X_0 \to Y_0$를 $B_0$ 위의 대수공간의 사상이라 하자.
$B_0$, $X_0$, $Y_0$가 준콤팩트이고 준분리라고 가정하자.
$f_i : X_i \to Y_i$를 $f_0$의 $B_i$로의 밑변환이라 하고,
$f : X \to Y$를 $f_0$의 $B$로의 밑변환이라 하자.
\end{situation}

\begin{lemma}
\label{spaces-limits-lemma-descend-etale}
표기와 가정은 Situation \ref{spaces-limits-situation-descent-property}와 같다.
다음을 가정하자.
\begin{enumerate}
\item $f$는 에탈이다.
\item $f_0$는 국소 유한 표시이다.
\end{enumerate}
그러면 $f_i$는 어떤 $i \geq 0$에 대해 에탈이다.
\end{lemma}

\begin{proof}
아핀 스킴 $V_0$와 전사 에탈 사상 $V_0 \to Y_0$를 택하자.
아핀 스킴 $U_0$와 전사 에탈 사상 $U_0 \to V_0 \times_{Y_0} X_0$를 택하자.
다음 가환도를 보라.
$$
\xymatrix{
U_0 \ar[d] \ar[r] & V_0 \ar[d] \\
X_0 \ar[r] & Y_0
}
$$
수직 화살표들은 구성에 의해 전사 에탈이다.
이 가환도를 $B_i$ 또는 $B$로 밑변환하면 다음을 얻는다.
$$
\vcenter{
\xymatrix{
U_i \ar[d] \ar[r] & V_i \ar[d] \\
X_i \ar[r] & Y_i
}
}
\quad\text{및}\quad
\vcenter{
\xymatrix{
U \ar[d] \ar[r] & V \ar[d] \\
X \ar[r] & Y
}
}
$$
$U_i, V_i, U, V$는 아핀 스킴이고 수직 사상들은 전사 에탈이며,
사상 $U_i \to V_i$의 극한은 $U \to V$임에 주의하자.
$X_i \to Y_i$가 에탈일 필요충분조건은 $U_i \to V_i$가 에탈인 것이며,
마찬가지로 $X \to Y$가 에탈일 필요충분조건은 $U \to V$가 에탈인 것이다
(Morphisms of Spaces, Lemma \ref{spaces-morphisms-lemma-etale-local}).
$f_0$가 국소 유한 표시이므로 사상 $U_0 \to V_0$도 국소 유한 표시이다.
따라서 보조정리는 Limits, Lemma \ref{limits-lemma-descend-etale}에서 따른다.
\end{proof}

\begin{lemma}
\label{spaces-limits-lemma-descend-smooth}
표기와 가정은 Situation \ref{spaces-limits-situation-descent-property}와 같다.
다음을 가정하자.
\begin{enumerate}
\item $f$는 매끄럽다.
\item $f_0$는 국소 유한 표시이다.
\end{enumerate}
그러면 $f_i$는 어떤 $i \geq 0$에 대해 매끄럽다.
\end{lemma}

\begin{proof}
아핀 스킴 $V_0$와 전사 에탈 사상 $V_0 \to Y_0$를 택하자.
아핀 스킴 $U_0$와 전사 에탈 사상 $U_0 \to V_0 \times_{Y_0} X_0$를 택하자.
다음 가환도를 보라.
$$
\xymatrix{
U_0 \ar[d] \ar[r] & V_0 \ar[d] \\
X_0 \ar[r] & Y_0
}
$$
수직 화살표들은 구성에 의해 전사 에탈이다.
이 가환도를 $B_i$ 또는 $B$로 밑변환하면 다음을 얻는다.
$$
\vcenter{
\xymatrix{
U_i \ar[d] \ar[r] & V_i \ar[d] \\
X_i \ar[r] & Y_i
}
}
\quad\text{및}\quad
\vcenter{
\xymatrix{
U \ar[d] \ar[r] & V \ar[d] \\
X \ar[r] & Y
}
}
$$
$U_i, V_i, U, V$는 아핀 스킴이고 수직 사상들은 전사 에탈이며,
사상 $U_i \to V_i$의 극한은 $U \to V$임에 주의하자.
$X_i \to Y_i$가 매끄러울 필요충분조건은 $U_i \to V_i$가 매끄러운 것이며,
마찬가지로 $X \to Y$가 매끄러울 필요충분조건은 $U \to V$가 매끄러운 것이다
(Morphisms of Spaces, Definition \ref{spaces-morphisms-definition-smooth}).
$f_0$가 국소 유한 표시이므로 사상 $U_0 \to V_0$도 국소 유한 표시이다.
따라서 보조정리는 Limits, Lemma \ref{limits-lemma-descend-smooth}에서 따른다.
\end{proof}

\begin{lemma}
\label{spaces-limits-lemma-descend-surjective}
표기와 가정은 Situation \ref{spaces-limits-situation-descent-property}와 같다.
다음을 가정하자.
\begin{enumerate}
\item $f$는 전사이다.
\item $f_0$는 국소 유한 표시이다.
\end{enumerate}
그러면 $f_i$는 어떤 $i \geq 0$에 대해 전사이다.
\end{lemma}

\begin{proof}
아핀 스킴 $V_0$와 전사 에탈 사상 $V_0 \to Y_0$를 택하자.
아핀 스킴 $U_0$와 전사 에탈 사상 $U_0 \to V_0 \times_{Y_0} X_0$를 택하자.
다음 가환도를 보라.
$$
\xymatrix{
U_0 \ar[d] \ar[r] & V_0 \ar[d] \\
X_0 \ar[r] & Y_0
}
$$
수직 화살표들은 구성에 의해 전사 에탈이다.
이 가환도를 $B_i$ 또는 $B$로 밑변환하면 다음을 얻는다.
$$
\vcenter{
\xymatrix{
U_i \ar[d] \ar[r] & V_i \ar[d] \\
X_i \ar[r] & Y_i
}
}
\quad\text{및}\quad
\vcenter{
\xymatrix{
U \ar[d] \ar[r] & V \ar[d] \\
X \ar[r] & Y
}
}
$$
$U_i, V_i, U, V$는 아핀 스킴이고 수직 사상들은 전사 에탈이며,
사상 $U_i \to V_i$의 극한은 $U \to V$이다.
또한 사상 $U_i \to X_i \times_{Y_i} V_i$ 및 $U \to X \times_Y V$는 전사이다
($U_0 \to X_0 \times_{Y_0} V_0$의 밑변환이기 때문이다).
특히 $X_i \to Y_i$가 전사일 필요충분조건은 $U_i \to V_i$가 전사인 것이며,
마찬가지로 $X \to Y$가 전사일 필요충분조건은 $U \to V$가 전사인 것이다.
$f_0$가 국소 유한 표시이므로 사상 $U_0 \to V_0$도 국소 유한 표시이다.
따라서 보조정리는 스킴의 경우
(Limits, Lemma \ref{limits-lemma-descend-surjective})에서 따른다.
\end{proof}

\begin{lemma}
\label{spaces-limits-lemma-descend-universally-injective}
표기와 가정은 Situation \ref{spaces-limits-situation-descent-property}와 같다.
다음을 가정하자.
\begin{enumerate}
\item $f$는 보편 단사이다.
\item $f_0$는 국소 유한형이다.
\end{enumerate}
그러면 $f_i$는 어떤 $i \geq 0$에 대해 보편 단사이다.
\end{lemma}

\begin{proof}
사상 $X \to Y$가 보편 단사일 필요충분조건은
대각 사상 $X \to X \times_Y X$가 전사인 것임을 상기하자
(Morphisms of Spaces, Definition
\ref{spaces-morphisms-definition-universally-injective} 및
Lemma \ref{spaces-morphisms-lemma-universally-injective}).
$X_0 \to X_0 \times_{Y_0} X_0$는 국소 유한 표시임을 관찰하자
(Morphisms of Spaces, Lemma
\ref{spaces-morphisms-lemma-diagonal-morphism-finite-type}).
따라서 사상 $X_0 \to X_0 \times_{Y_0} X_0$를 생각하고
Lemma \ref{spaces-limits-lemma-descend-surjective}를 적용하면 보조정리가 따른다.
\end{proof}

\begin{lemma}
\label{spaces-limits-lemma-descend-affine}
표기와 가정은 Situation \ref{spaces-limits-situation-descent-property}와 같다.
$f$가 아핀이면 $f_i$는 어떤 $i \geq 0$에 대해 아핀이다.
\end{lemma}

\begin{proof}
아핀 스킴 $V_0$와 전사 에탈 사상 $V_0 \to Y_0$를 택하자.
$V_i = V_0 \times_{Y_0} Y_i$ 및 $V = V_0 \times_{Y_0} Y$로 놓자.
$f$가 아핀이므로 $V \times_Y X = \lim V_i \times_{Y_i} X_i$는 아핀이다.
Lemma \ref{spaces-limits-lemma-limit-is-affine}에 의해
$V_i \times_{Y_i} X_i$는 어떤 $i \geq 0$에 대해 아핀이다.
이 $i$에 대해 사상 $f_i$는 아핀이다
(Morphisms of Spaces, Lemma \ref{spaces-morphisms-lemma-affine-local}).
\end{proof}

\begin{lemma}
\label{spaces-limits-lemma-descend-finite}
표기와 가정은 Situation \ref{spaces-limits-situation-descent-property}와 같다.
다음을 가정하자.
\begin{enumerate}
\item $f$는 유한이다.
\item $f_0$는 국소 유한형이다.
\end{enumerate}
그러면 $f_i$는 어떤 $i \geq 0$에 대해 유한이다.
\end{lemma}

\begin{proof}
아핀 스킴 $V_0$와 전사 에탈 사상 $V_0 \to Y_0$를 택하자.
$V_i = V_0 \times_{Y_0} Y_i$ 및 $V = V_0 \times_{Y_0} Y$로 놓자.
$f$가 유한이므로 $V \times_Y X = \lim V_i \times_{Y_i} X_i$는
$V$ 위에서 유한인 스킴이다.
Lemma \ref{spaces-limits-lemma-limit-is-affine}에 의해
$V_i \times_{Y_i} X_i$는 어떤 $i \geq 0$에 대해 아핀이다.
필요하면 $i$를 크게 하여 $V_i \times_{Y_i} X_i \to V_i$가 유한임을
Limits, Lemma \ref{limits-lemma-descend-finite-finite-presentation}에서 얻는다.
이 $i$에 대해 사상 $f_i$는 유한이다
(Morphisms of Spaces, Lemma \ref{spaces-morphisms-lemma-integral-local}).
\end{proof}

\begin{lemma}
\label{spaces-limits-lemma-descend-closed-immersion}
표기와 가정은 Situation \ref{spaces-limits-situation-descent-property}와 같다.
다음을 가정하자.
\begin{enumerate}
\item $f$는 닫힌 몰입이다.
\item $f_0$는 국소 유한형이다.
\end{enumerate}
그러면 $f_i$는 어떤 $i \geq 0$에 대해 닫힌 몰입이다.
\end{lemma}

\begin{proof}
아핀 스킴 $V_0$와 전사 에탈 사상 $V_0 \to Y_0$를 택하자.
$V_i = V_0 \times_{Y_0} Y_i$ 및 $V = V_0 \times_{Y_0} Y$로 놓자.
$f$가 닫힌 몰입이므로 $V \times_Y X = \lim V_i \times_{Y_i} X_i$는
아핀 스킴 $V$의 닫힌 부분스킴이다.
Lemma \ref{spaces-limits-lemma-limit-is-affine}에 의해
$V_i \times_{Y_i} X_i$는 어떤 $i \geq 0$에 대해 아핀이다.
필요하면 $i$를 크게 하여 $V_i \times_{Y_i} X_i \to V_i$가 닫힌 몰입임을
Limits, Lemma \ref{limits-lemma-descend-closed-immersion-finite-presentation}에서 얻는다.
이 $i$에 대해 사상 $f_i$는 닫힌 몰입이다
(Morphisms of Spaces, Lemma \ref{spaces-morphisms-lemma-integral-local}).
\end{proof}

\begin{lemma}
\label{spaces-limits-lemma-descend-separated-morphism}
표기와 가정은 Situation \ref{spaces-limits-situation-descent-property}와 같다.
$f$가 분리이면 $f_i$는 어떤 $i \geq 0$에 대해 분리이다.
\end{lemma}

\begin{proof}
Lemma \ref{spaces-limits-lemma-descend-closed-immersion}을
대각 사상 $\Delta_{X_0/Y_0} : X_0 \to X_0 \times_{Y_0} X_0$에 적용하자.
(대각 사상은 국소 유한형이고,
올곱 $X_0 \times_{Y_0} X_0$는 준콤팩트이고 준분리이다.
몇 가지 세부 사항은 생략한다.)
\end{proof}

\begin{lemma}
\label{spaces-limits-lemma-descend-isomorphism}
표기와 가정은 Situation \ref{spaces-limits-situation-descent-property}와 같다.
다음을 가정하자.
\begin{enumerate}
\item $f$는 동형이다.
\item $f_0$는 국소 유한 표시이다.
\end{enumerate}
그러면 $f_i$는 어떤 $i \geq 0$에 대해 동형이다.
\end{lemma}

\begin{proof}
동형이라는 것은 에탈이고 보편 단사이며 전사라는 것과 동치이다.
Morphisms of Spaces, Lemma
\ref{spaces-morphisms-lemma-etale-universally-injective-open}을 보라.
따라서 보조정리는 Lemmas \ref{spaces-limits-lemma-descend-etale},
\ref{spaces-limits-lemma-descend-surjective}, \ref{spaces-limits-lemma-descend-universally-injective}에서 따른다.
\end{proof}

\begin{lemma}
\label{spaces-limits-lemma-descend-monomorphism}
표기와 가정은 Situation \ref{spaces-limits-situation-descent-property}와 같다.
다음을 가정하자.
\begin{enumerate}
\item $f$는 단사상이다.
\item $f_0$는 국소 유한형이다.
\end{enumerate}
그러면 $f_i$는 어떤 $i \geq 0$에 대해 단사상이다.
\end{lemma}

\begin{proof}
사상이 단사상일 필요충분조건은 그 대각 사상이 동형인 것임을 상기하자.
사상 $X_0 \to X_0 \times_{Y_0} X_0$는
Morphisms of Spaces, Lemma
\ref{spaces-morphisms-lemma-diagonal-morphism-finite-type}에 의해 국소 유한 표시이다.
$X_0 \times_{Y_0} X_0$가 준콤팩트이고 준분리이므로,
Lemma \ref{spaces-limits-lemma-descend-isomorphism}에 의해
$\Delta_i : X_i \to X_i \times_{Y_i} X_i$는 어떤 $i \geq 0$에 대해 동형이다.
이 $i$에 대해 사상 $f_i$는 단사상이다.
\end{proof}

\begin{lemma}
\label{spaces-limits-lemma-descend-flat}
표기와 가정은 Situation \ref{spaces-limits-situation-descent-property}와 같다.
$\mathcal{F}_0$를 준연접 $\mathcal{O}_{X_0}$-가군이라 하고,
$\mathcal{F}_i$로 그 $X_i$로의 당김을,
$\mathcal{F}$로 그 $X$로의 당김을 나타내자.
다음을 가정하자.
\begin{enumerate}
\item $\mathcal{F}$는 $Y$ 위에서 평탄하다.
\item $\mathcal{F}_0$는 유한 표시이다.
\item $f_0$는 국소 유한 표시이다.
\end{enumerate}
그러면 $\mathcal{F}_i$는 $Y_i$ 위에서 어떤 $i \geq 0$에 대해 평탄하다.
특히 $f_0$가 국소 유한 표시이고 $f$가 평탄하면,
$f_i$는 어떤 $i \geq 0$에 대해 평탄하다.
\end{lemma}

\begin{proof}
아핀 스킴 $V_0$와 전사 에탈 사상 $V_0 \to Y_0$를 택하자.
아핀 스킴 $U_0$와 전사 에탈 사상 $U_0 \to V_0 \times_{Y_0} X_0$를 택하자.
다음 가환도를 보라.
$$
\xymatrix{
U_0 \ar[d] \ar[r] & V_0 \ar[d] \\
X_0 \ar[r] & Y_0
}
$$
수직 화살표들은 구성에 의해 전사 에탈이다.
이 가환도를 $B_i$ 또는 $B$로 밑변환하면 다음을 얻는다.
$$
\vcenter{
\xymatrix{
U_i \ar[d] \ar[r] & V_i \ar[d] \\
X_i \ar[r] & Y_i
}
}
\quad\text{및}\quad
\vcenter{
\xymatrix{
U \ar[d] \ar[r] & V \ar[d] \\
X \ar[r] & Y
}
}
$$
$U_i, V_i, U, V$는 아핀 스킴이고 수직 사상들은 전사 에탈이며,
사상 $U_i \to V_i$의 극한은 $U \to V$임에 주의하자.
$\mathcal{F}_i$가 $Y_i$ 위에서 평탄할 필요충분조건은
$\mathcal{F}_i|_{U_i}$가 $V_i$ 위에서 평탄한 것이며,
마찬가지로 $\mathcal{F}$가 $Y$ 위에서 평탄할 필요충분조건은
$\mathcal{F}|_U$가 $V$ 위에서 평탄한 것이다
(Morphisms of Spaces, Definition \ref{spaces-morphisms-definition-flat}).
$f_0$가 국소 유한 표시이므로 사상 $U_0 \to V_0$도 국소 유한 표시이다.
따라서 보조정리는 Limits, Lemma
\ref{limits-lemma-descend-module-flat-finite-presentation}에서 따른다.
\end{proof}

\begin{lemma}
\label{spaces-limits-lemma-eventually-proper}
가정과 표기는 Situation \ref{spaces-limits-situation-descent-property}와 같다.
다음을 가정하자.
\begin{enumerate}
\item $f$는 고유이다.
\item $f_0$는 국소 유한형이다.
\end{enumerate}
그러면 $i$가 존재하여 $f_i$가 고유이다.
\end{lemma}

\begin{proof}
아핀 스킴 $V_0$와 전사 에탈 사상 $V_0 \to Y_0$를 택하자.
$V_i = Y_i \times_{Y_0} V_0$ 및 $V = Y \times_{Y_0} V_0$로 놓자.
$f_i$를 $V_i$로 밑변환한 것이 고유임을 보이면 충분하다.
Morphisms of Spaces, Lemma \ref{spaces-morphisms-lemma-proper-local}을 보라.
따라서 $Y_0$가 아핀이라고 가정할 수 있다.

\medskip\noindent
Lemma \ref{spaces-limits-lemma-descend-separated-morphism}에 의해
$f_i$는 어떤 $i \geq 0$에 대해 분리이다.
$0$을 $i$로 대체하여 $f_0$가 분리라고 가정할 수 있다.
$f_0$가 준콤팩트임을 관찰하자. 따라서 $f_0$는 분리이고 유한형이다.
Cohomology of Spaces, Lemma \ref{spaces-cohomology-lemma-weak-chow}에 의해
다음 가환도를 택할 수 있다.
$$
\xymatrix{
X_0 \ar[rd] & X_0' \ar[d] \ar[l]^\pi \ar[r] & \mathbf{P}^n_{Y_0} \ar[dl] \\
& Y_0 &
}
$$
여기서 $X_0' \to \mathbf{P}^n_{Y_0}$는 몰입이고,
$\pi : X_0' \to X_0$는 고유이며 전사이다.
$X' = X_0' \times_{Y_0} Y$ 및 $X_i' = X_0' \times_{Y_0} Y_i$를 도입하자.
Morphisms of Spaces, Lemmas \ref{spaces-morphisms-lemma-composition-proper}와
\ref{spaces-morphisms-lemma-base-change-proper}에 의해
$X' \to Y$는 고유이다.
따라서 $X' \to \mathbf{P}^n_Y$는 닫힌 몰입이다
(Morphisms of Spaces, Lemma
\ref{spaces-morphisms-lemma-universally-closed-permanence}).
Morphisms of Spaces, Lemma \ref{spaces-morphisms-lemma-image-proper-is-proper}에 의해
$X'_i \to Y_i$가 어떤 $i$에 대해 고유임을 보이면 충분하다.
Lemma \ref{spaces-limits-lemma-descend-closed-immersion}에 의해
$X'_i \to \mathbf{P}^n_{Y_i}$는 충분히 큰 $i$에 대해 닫힌 몰입이다.
그러면 $X'_i \to Y_i$는 고유이므로 증명이 끝난다.
\end{proof}

\begin{lemma}
\label{spaces-limits-lemma-eventually-relative-dimension}
가정과 표기는 Situation \ref{spaces-limits-situation-descent-property}와 같다.
$d \geq 0$라 하자. 다음을 가정하자.
\begin{enumerate}
\item $f$의 상대 차원은 $\leq d$이다
(Morphisms of Spaces, Definition \ref{spaces-morphisms-definition-relative-dimension}).
\item $f_0$는 국소 유한형이다.
\end{enumerate}
그러면 $i$가 존재하여 $f_i$의 상대 차원이 $\leq d$이다.
\end{lemma}

\begin{proof}
아핀 스킴 $V_0$와 전사 에탈 사상 $V_0 \to Y_0$를 택하자.
아핀 스킴 $U_0$와 전사 에탈 사상 $U_0 \to V_0 \times_{Y_0} X_0$를 택하자.
다음 가환도를 보라.
$$
\xymatrix{
U_0 \ar[d] \ar[r] & V_0 \ar[d] \\
X_0 \ar[r] & Y_0
}
$$
수직 화살표들은 구성에 의해 전사 에탈이다.
이 가환도를 $B_i$ 또는 $B$로 밑변환하면 다음을 얻는다.
$$
\vcenter{
\xymatrix{
U_i \ar[d] \ar[r] & V_i \ar[d] \\
X_i \ar[r] & Y_i
}
}
\quad\text{및}\quad
\vcenter{
\xymatrix{
U \ar[d] \ar[r] & V \ar[d] \\
X \ar[r] & Y
}
}
$$
$U_i, V_i, U, V$는 아핀 스킴이고 수직 사상들은 전사 에탈이며,
사상 $U_i \to V_i$의 극한은 $U \to V$임에 주의하자.
이 상황에서 $X_i \to Y_i$의 상대 차원이 $\leq d$일 필요충분조건은
$U_i \to V_i$의 상대 차원이 $\leq d$인 것이다
(Morphisms, Definition \ref{morphisms-definition-relative-dimension-d}의 정의를 따른다).
이 동치를 보려면 대수공간의 사상에 대한 정의가
Morphisms of Spaces, Definition \ref{spaces-morphisms-definition-dimension-fibre}를
사용하며, 그 정의는 에탈 국소화를 사용한다는 사실을 이용하면 된다.
$X \to Y$와 $U \to V$에 대해서도 같은 사실이 성립한다.
$f_0$가 국소 유한형이므로 사상 $U_0 \to V_0$도 국소 유한형이다.
따라서 보조정리는 더 일반적인
Limits, Lemma \ref{limits-lemma-limit-dimension}에서 따른다.
\end{proof}



\section{상대적 대상의 하강}
\label{spaces-limits-section-descending-relative}

\noindent
다음 보조정리는 이 절에서 다루는 결과의 전형이다.

\begin{lemma}
\label{spaces-limits-lemma-descend-finite-presentation}
$S$를 스킴이라 하자. $I$를 유향 집합이라 하자.
$(X_i, f_{ii'})$를 $I$를 지표로 하는 $S$ 위의 대수공간의 역계라 하자.
다음을 가정하자.
\begin{enumerate}
\item 사상 $f_{ii'} : X_i \to X_{i'}$는 아핀이다.
\item 공간 $X_i$는 준콤팩트이고 준분리이다.
\end{enumerate}
$X = \lim_i X_i$라 하자.
그러면 $X$ 위에서 유한 표시인 대수공간의 범주는
$I$에 걸친 여극한으로 주어지며, 그 계는
$X_i$ 위에서 유한 표시인 대수공간의 범주들로 이루어진다.
\end{lemma}

\begin{proof}
$0 \in I$를 택하자. 전사 에탈 사상 $U_0 \to X_0$를 택하되,
$U_0$는 아핀 스킴이라고 하자
(Properties of Spaces, Lemma \ref{spaces-properties-lemma-quasi-compact-affine-cover}).
$U_i = X_i \times_{X_0} U_0$로 놓자.
$R_0 = U_0 \times_{X_0} U_0$ 및 $R_i = R_0 \times_{X_0} X_i$로 놓자.
두 사영을 $s_i, t_i : R_i \to U_i$ 및 $s, t : R \to U$라 쓰자.
Lemma \ref{spaces-limits-lemma-directed-inverse-system-has-limit}의 증명에서
표시 $X = U/R$가 존재하여 $U = \lim U_i$ 및 $R = \lim R_i$임을 보았다.
$U_i$와 $U$는 아핀이며, $R_i$와 $R$은 준콤팩트이고 분리임에 주의하자
($X_i$가 준분리이기 때문이다).
$Y$를 $S$ 위의 대수공간이라 하고 $Y \to X$를 유한 표시 사상이라 하자.
$V = U \times_X Y$로 놓자. 이는 $U$ 위에서 유한 표시인 대수공간이다.
아핀 스킴 $W$와 전사 에탈 사상 $W \to V$를 택하자.
그러면 $W \to Y$도 전사 에탈이다.
$R' = W \times_Y W$로 놓으면 $Y = W/R'$이다
(Spaces, Section \ref{spaces-section-presentations} 참조).
$W$는 $U$ 위에서 유한 표시인 스킴이고,
$R'$은 $R$ 위에서 유한 표시인 스킴임에 주의하자(세부 사항은 생략한다).
Limits, Lemma \ref{limits-lemma-descend-finite-presentation}에 의해
지표 $i$와 유한 표시인 스킴의 사상 $W_i \to U_i$를 찾아,
그 $U$로의 밑변환이 $W \to U$가 되게 할 수 있다.
마찬가지로 필요하면 $i$를 크게 하여 스킴 $R'_i$를 찾을 수 있다.
이는 $R_i$ 위에서 유한 표시이고, 그 $R$로의 밑변환은 $R'$이다.
사영 사상 $s', t' : R' \to W$는
사영 사상 $s, t : R \to U$ 위의 사상이다.
따라서 $s'$, $t'$를 각각 $U$ 위에서 유한 표시인 스킴 사이의 사상으로 볼 수 있다
(구조 사상 $R' \to U$는 $R' \to R$ 뒤에 각각 $s$, $t$를 합성하여 주어진다).
그러므로 Limits, Lemma \ref{limits-lemma-descend-finite-presentation}을
다시 적용하면, 필요에 따라 $i$를 크게 한 후
사상 $s'_i, t'_i : R'_i \to W_i$가 존재하여
그 $U$로의 밑변환이 $S', t'$가 됨을 알 수 있다.
Limits, Lemmas \ref{limits-lemma-descend-etale} 및
\ref{limits-lemma-descend-monomorphism}에 의해
$s'_i, t'_i$가 에탈이고,
$j'_i : R'_i \to W_i \times_{X_i} W_i$가 단사상이라고 가정할 수 있다
(여기서 $j'_i$는 $U_i$ 위에서 유한 표시인 스킴 사이의 사상으로 본다.
구조 사상에는 어느 사영을 사용해도 상관없다).
$Y_i = W_i/R'_i$로 놓으면
(Spaces, Theorem \ref{spaces-theorem-presentation} 참조),
$X_i$ 위에서 유한 표시인 대수공간을 얻고,
그 $X$로의 밑변환은 $Y$와 동형이다.

\medskip\noindent
이로써 $X$ 위에서 유한 표시인 모든 대수공간이
어떤 $X_i$ 위에서 유한 표시인 대수공간에서 옴을 보였다.
즉 보조정리의 함자가 본질적 전사임을 보였다.
이 함자가 충만 충실함을 보이기 위해, 지표 $0 \in I$와
두 대수공간 $Y_0, Z_0$를 $X_0$ 위에서 유한 표시인 것으로 택하자.
$Y_i = X_i \times_{X_0} Y_0$, $Y = X \times_{X_0} Y_0$,
$Z_i = X_i \times_{X_0} Z_0$, $Z = X \times_{X_0} Z_0$로 놓자.
$\alpha : Y \to Z$를 $X$ 위의 대수공간의 사상이라 하자.
전사 에탈 사상 $V_0 \to Y_0$를 택하되 $V_0$는 아핀 스킴이라고 하자.
$V_i = V_0 \times_{Y_0} Y_i$ 및 $V = V_0 \times_{Y_0} Y$로 놓자.
이들은 각각 $Y_i$와 $Y$로 가는 전사 에탈 사상을 갖춘 아핀 스킴이다.
합성 $V \to Y \to Z \to Z_0$는 본질적으로 유일한 사상
$V_i \to Z_0$에서 어떤 $i \geq 0$에 대해 온다.
이는 Proposition \ref{spaces-limits-proposition-characterize-locally-finite-presentation}에 따른다
($Z_0 \to X_0$에 적용하며, 이 사상은 가정에 의해 유한 표시이다).
$i$를 크게 한 후에는 두 합성
$$
V_i \times_{Y_i} V_i \to V_i \to Z_0
$$
이 같다. 극한에서 두 합성이 같기 때문이다.
따라서 본질적으로 유일한 사상 $Y_i \to Z_0$를 얻는다.
이는 $X_0$ 위의 사상이므로,
원하는 대로 $Z_i = Z_0 \times_{X_0} X_i$로 가는 사상을 유도한다.
\end{proof}

\begin{lemma}
\label{spaces-limits-lemma-descend-modules-finite-presentation}
표기와 가정은 Lemma \ref{spaces-limits-lemma-descend-finite-presentation}과 같다.
유한 표시인 $\mathcal{O}_X$-가군의 범주는 $I$에 걸친 여극한으로 주어지며,
그 계는 유한 표시인 $\mathcal{O}_{X_i}$-가군의 범주들로 이루어진다.
\end{lemma}

\begin{proof}
$0 \in I$를 택하자. 아핀 스킴 $U_0$와 전사 에탈 사상 $U_0 \to X_0$를 택하자.
$U_i = X_i \times_{X_0} U_0$로 놓자.
$R_0 = U_0 \times_{X_0} U_0$ 및 $R_i = R_0 \times_{X_0} X_i$로 놓자.
두 사영을 $s_i, t_i : R_i \to U_i$ 및 $s, t : R \to U$라 쓰자.
Lemma \ref{spaces-limits-lemma-directed-inverse-system-has-limit}의 증명에서
표시 $X = U/R$가 존재하여 $U = \lim U_i$ 및 $R = \lim R_i$임을 보았다.
$U_i$와 $U$는 아핀이며, $R_i$와 $R$은 준콤팩트이고 분리임에 주의하자
($X_i$가 준분리이기 때문이다). 또한
$R \times_{s, U, t} R = \colim R_i \times_{s_i, U_i, t_i} R_i$도 성립한다.
따라서
$\QCoh(\mathcal{O}_U) = \colim \QCoh(\mathcal{O}_{U_i})$,
$\QCoh(\mathcal{O}_R) = \colim \QCoh(\mathcal{O}_{R_i})$ 및
$\QCoh(\mathcal{O}_{R \times_{s, U, t} R}) =
\colim \QCoh(\mathcal{O}_{R_i \times_{s_i, U_i, t_i} R_i})$임을
Limits, Lemma \ref{limits-lemma-descend-modules-finite-presentation}에 의해 알 수 있다.
$\QCoh(\mathcal{O}_X) = \QCoh(U, R, s, t, c)$ 및
$\QCoh(\mathcal{O}_{X_i}) = \QCoh(U_i, R_i, s_i, t_i, c_i)$가 성립한다.
Properties of Spaces, Proposition \ref{spaces-properties-proposition-quasi-coherent}를 보라.
따라서 결과는 형식적으로 따른다.
\end{proof}

\begin{lemma}
\label{spaces-limits-lemma-descend-invertible-modules}
표기와 가정은 Lemma \ref{spaces-limits-lemma-descend-finite-presentation}과 같다.
그러면 다음이 성립한다.
\begin{enumerate}
\item 모든 유한 국소 자유 $\mathcal{O}_X$-가군은
유한 국소 자유 $\mathcal{O}_{X_i}$-가군의 당김으로 어떤 $i$에 대해 주어진다.
\item 모든 가역 $\mathcal{O}_X$-가군은
가역 $\mathcal{O}_{X_i}$-가군의 당김으로 어떤 $i$에 대해 주어진다.
\end{enumerate}
\end{lemma}

\begin{proof}
(2)의 증명.
$\mathcal{L}$을 가역 $\mathcal{O}_X$-가군이라 하자.
가역 가군은 유한 표시이므로, $i$와 유한 표시 가군
$\mathcal{L}_i$, $\mathcal{N}_i$를 $X_i$ 위에서 찾아
$f_i^*\mathcal{L}_i \cong \mathcal{L}$ 및
$f_i^*\mathcal{N}_i \cong \mathcal{L}^{\otimes -1}$이 되게 할 수 있다.
Lemma \ref{spaces-limits-lemma-descend-modules-finite-presentation}을 보라.
당김은 텐서곱과 교환하므로
$f_i^*(\mathcal{L}_i \otimes_{\mathcal{O}_{X_i}} \mathcal{N}_i)$는
$\mathcal{O}_X$와 동형이다.
유한 표시 가군의 텐서곱은 유한 표시이므로, 같은 보조정리에 의해
$f_{i'i}^*\mathcal{L}_i
\otimes_{\mathcal{O}_{X_{i'}}} f_{i'i}^*\mathcal{N}_i$는
$\mathcal{O}_{X_{i'}}$와 어떤 $i' \geq i$에 대해 동형이다.
따라서 $f_{i'i}^*\mathcal{L}_i$는 가역이고
(Modules on Sites, Lemma \ref{sites-modules-lemma-invertible}),
증명이 끝난다.

\medskip\noindent
(1)의 증명. 생략한다. 힌트: (2)의 증명과 같이 논증하되,
(국소환 달린 사이트 위의) 가군이 유한 국소 자유일 필요충분조건이
쌍대를 갖는 것이라는 사실을 사용하라.
Modules on Sites, Section \ref{sites-modules-section-duals}를 보라.
또는 스킴에 대한 증명과 같이 논증하라.
Limits, Lemma \ref{limits-lemma-descend-invertible-modules}를 보라.
\end{proof}



\section{절대 Noether 근사}
\label{spaces-limits-section-approximation}

\noindent
다음 결과는 \cite[Theorem 1.2.2]{CLO}이다.
증명의 핵심적인 도구는 Decent Spaces, Lemma
\ref{decent-spaces-lemma-filter-quasi-compact-quasi-separated}이다.

\begin{proposition}
\label{spaces-limits-proposition-approximate}
\begin{reference}
우리의 증명은 \cite[Theorem 1.2.2]{CLO}에 제시된 증명을 밀접하게 따른다.
\end{reference}
$X$를 준콤팩트 준분리 대수공간으로서
$\Spec(\mathbf{Z})$ 위에 주어진 것이라 하자.
유향 집합 $I$와 대수공간의 역계 $(X_i, f_{ii'})$가 $I$ 위에 존재하여
다음을 만족한다.
\begin{enumerate}
\item 전이 사상 $f_{ii'}$는 아핀이다.
\item 각 $X_i$는 준분리이고 $\mathbf{Z}$ 위에서 유한형이다.
\item $X = \lim X_i$이다.
\end{enumerate}
\end{proposition}

\begin{proof}
Decent Spaces, Lemma
\ref{decent-spaces-lemma-filter-quasi-compact-quasi-separated}를 적용하여
열린 부분공간 $U_p \subset X$, 스킴 $V_p$, 사상 $f_p : V_p \to U_p$를
명시된 성질들을 만족하도록 얻는다.
$f_n : V_n \to U_n$는 대수공간의 에탈 사상이며,
$T_n = (V_n)_{red}$의 역상에 제한하면 동형이 됨에 주의하자.
따라서 $f_n$는 동형이다. 예를 들어 Morphisms of Spaces,
Lemma \ref{spaces-morphisms-lemma-etale-universally-injective-open}에 따른다.
특히 $U_n$는 준콤팩트 분리 스킴이다.
따라서 $U_n = \lim U_{n, i}$를 $\mathbf{Z}$ 위의 유한형 스킴의 유향 극한으로
쓰되, 전이 사상들이 아핀이게 할 수 있다.
Limits, Proposition \ref{limits-proposition-approximate}를 보라.
그러므로 $p$에 대한 내림 귀납법을 적용하면,
다음 문단에 제시된 문제로 환원됨을 알 수 있다.

\medskip\noindent
여기서는 $U \subset X$, $U = \lim U_i$, $Z \subset X$ 및
$f : V \to X$가 주어지고 다음 성질들을 만족한다.
\begin{enumerate}
\item $X$는 준콤팩트 준분리 대수공간이다.
\item $V$는 준콤팩트 분리 스킴이다.
\item $U \subset X$는 준콤팩트 열린 부분공간이다.
\item $(U_i, g_{ii'})$는 준분리 대수공간의 유향 역계이며,
각 대수공간은 $\mathbf{Z}$ 위에서 유한형이고, 전이 사상들이 아핀이며,
그 극한은 $U$이다.
\item $Z \subset X$는 닫힌 부분공간으로서 $|X| = |U| \amalg |Z|$를 만족한다.
\item $f : V \to X$는 전사 에탈 사상으로서
$f^{-1}(Z) \to Z$가 동형이다.
\end{enumerate}
문제: 명제의 결론이 $X$에 대해 성립함을 보여라.

\medskip\noindent
$W = f^{-1}(U) \subset V$는 준콤팩트 열린 부분스킴이며
$U$ 위에서 에탈임에 주의하자.
따라서 Lemmas \ref{spaces-limits-lemma-descend-finite-presentation} 및 \ref{spaces-limits-lemma-descend-etale}을
적용하여 지표 $0 \in I$와 유한 표시인 에탈 사상 $W_0 \to U_0$를 찾아,
그 $U$로의 밑변환이 $W$가 되게 할 수 있다.
$W_i = W_0 \times_{U_0} U_i$로 놓으면 $W = \lim_{i \geq 0} W_i$이다.
$0$을 크게 한 후에는 $W_i$들이 스킴이라고 가정할 수 있다.
Lemma \ref{spaces-limits-lemma-limit-is-scheme}를 보라.
또한 $W_i$는 $\mathbf{Z}$ 위에서 유한형이다.

\medskip\noindent
Limits, Lemma \ref{limits-lemma-approximate}를
$W = \lim_{i \geq 0} W_i$와 포함 사상 $W \subset V$에 적용하자.
$I$를 그 보조정리에서 얻은 유향 집합 $J$로 대체하자.
그러면 $V$를 유향 극한 $V = \lim V_i$로 쓸 수 있다.
여기서 각 스킴은 $\mathbf{Z}$ 위에서 유한형이고 전이 사상들은 아핀이며,
각 $V_i$는 $W_i$를 열린 부분스킴으로 포함한다
(전이 사상들과 호환되게 포함한다).
각 $i$에 대해 스킴의 범주에서 다음 푸시아웃을 만들 수 있다.
$$
\xymatrix{
W_i \ar[r] \ar[d]_\Delta & V_i \ar[d] \\
W_i \times_{U_i} W_i \ar[r] & R_i
}
$$
실제로 왼쪽 수직 화살표와 위쪽 수평 화살표는 스킴의 열린 몰입이다.
달리 말하면 $R_i$를 $V_i$와 $W_i \times_{U_i} W_i$를
공통 열린집합 $W_i$를 따라 붙여서 구성할 수 있다
(Schemes, Section \ref{schemes-section-glueing-schemes} 참조).
에탈 사영 사상 $W_i \times_{U_i} W_i \to W_i$는
에탈 사상 $s_i, t_i : R_i \to V_i$로 확장됨에 주의하자.
사상 $j_i = (t_i, s_i) : R_i \to V_i \times V_i$는
$V_i$ 위의 에탈 동치관계임이 분명하다.
$W_i \times_{U_i} W_i$는 준콤팩트이고
($U_i$가 준분리이고 $W_i$가 준콤팩트이기 때문이다),
$V_i$도 준콤팩트이므로 $R_i$는 준콤팩트이다.
$i \geq i'$에 대해 가환도
\begin{equation}
\label{spaces-limits-equation-cartesian}
\vcenter{
\xymatrix{
R_i \ar[r] \ar[d]_{s_i} & R_{i'} \ar[d]^{s_{i'}} \\
V_i \ar[r] & V_{i'}
}
}
\end{equation}
는 데카르트 가환도이다. 이는
$$
(W_{i'} \times_{U_{i'}} W_{i'}) \times_{U_{i'}} U_i =
W_{i'} \times_{U_{i'}} U_i \times_{U_i} U_i \times_{U_{i'}} W_{i'} =
W_i \times_{U_i} W_i.
$$
이기 때문이다.
대수공간 $X_i = V_i/R_i$를 생각하자
(Spaces, Theorem \ref{spaces-theorem-presentation} 참조).
$V_i$가 $\mathbf{Z}$ 위에서 유한형이고 $R_i$가 준콤팩트이므로,
$X_i$는 준분리이고 $\mathbf{Z}$ 위에서 유한형이다
(Properties of Spaces, Lemma \ref{spaces-properties-lemma-quasi-separated}와
Morphisms of Spaces, Lemmas \ref{spaces-morphisms-lemma-surjection-from-quasi-compact} 및
\ref{spaces-morphisms-lemma-finite-type-local}을 보라).
위의 $R_i$의 구성은 전이 사상들과 호환되므로,
대수공간의 사상 $X_i \to X_{i'}$를 $i \geq i'$에 대해 얻는다.
가환도
$$
\xymatrix{
V_i \ar[r] \ar[d] & V_{i'} \ar[d] \\
X_i \ar[r] & X_{i'}
}
$$
들은 (\ref{spaces-limits-equation-cartesian})가 데카르트 가환도이므로 데카르트 가환도이다.
Groupoids, Lemma \ref{groupoids-lemma-criterion-fibre-product}를 보라.
$V_i \to V_{i'}$가 아핀이므로 $X_i \to X_{i'}$도 아핀이다.
Morphisms of Spaces, Lemma \ref{spaces-morphisms-lemma-affine-local}을 보라.
따라서 극한 $X' = \lim X_i$를
Lemma \ref{spaces-limits-lemma-directed-inverse-system-has-limit}에 의해 만들 수 있다.
$X \cong X'$라고 주장한다. 이로써 명제의 증명이 끝난다.

\medskip\noindent
주장의 증명. $R = \lim R_i$로 놓자.
구성에 의해 대수공간 $X'$에는 전사 에탈 사상 $V \to X'$가 주어지며,
$$
V \times_{X'} V \cong R
$$
가 성립한다(Lemma \ref{spaces-limits-lemma-directed-inverse-system-has-limit}를 사용하라).
구성에 의해 $\lim W_i \times_{U_i} W_i = W \times_U W$이고 $V = \lim V_i$이므로,
$R$은 $W \times_U W$와 $V$를 $W$를 따라 붙인 합집합이다.
성질 (6)에 의해 사영 $V \times_X V \to V$는
$f^{-1}(Z) \subset V$ 위에서 동형이다.
따라서 스킴 $V \times_X V$는 열린집합 $\Delta_{V/X}(V)$와
$W \times_U W$의 합집합이며, 이들의 교집합은 $\Delta_{W/X}(W)$이다.
그러므로 동형 $R \cong V \times_X V$가 유일하게 존재하여
$V$로의 사영들과 호환된다.
$V \to X$와 $V \to X'$가 전사 에탈이므로
$$
X = V/ V \times_X V = V/R = V/V \times_{X'} V = X'
$$
임을 Spaces, Lemma \ref{spaces-lemma-space-presentation}에 의해 알 수 있고,
증명이 끝난다.
\end{proof}



\section{응용}
\label{spaces-limits-section-applications}

\noindent
다음 보조정리는 절대 Noether 근사를 거치지 않고도
Decent Spaces, Lemma
\ref{decent-spaces-lemma-filter-quasi-compact-quasi-separated}에서 직접 유도할 수 있다.

\begin{lemma}
\label{spaces-limits-lemma-colimit-finitely-presented}
$S$를 스킴이라 하자. $X$를 $S$ 위의 준콤팩트 준분리 대수공간이라 하자.
모든 준연접 $\mathcal{O}_X$-가군은
유한 표시 $\mathcal{O}_X$-가군들의 여과 여극한이다.
\end{lemma}

\begin{proof}
$X$를 $\Spec(\mathbf{Z})$ 위의 대수공간으로 볼 수 있다.
Spaces, Definition \ref{spaces-definition-base-change} 및
Properties of Spaces, Definition \ref{spaces-properties-definition-separated}를 보라.
따라서 Proposition \ref{spaces-limits-proposition-approximate}를 적용하여
$X = \lim X_i$로 쓰되 $X_i$가 $\mathbf{Z}$ 위에서 유한 표시이게 할 수 있다.
그러면 $X_i$는 Noether 대수공간이다.
Morphisms of Spaces, Lemma
\ref{spaces-morphisms-lemma-finite-presentation-noetherian}을 보라.
사상 $X \to X_i$는 아핀이다.
Lemma \ref{spaces-limits-lemma-directed-inverse-system-has-limit}를 보라.
결론은 Cohomology of Spaces, Lemma
\ref{spaces-cohomology-lemma-direct-colimit-finite-presentation}에서 따른다.
\end{proof}

\noindent
이 절의 나머지 내용은 Lemma \ref{spaces-limits-lemma-colimit-finitely-presented}의
직접적인 응용들이다.

\begin{lemma}
\label{spaces-limits-lemma-directed-colimit-finite-type}
$S$를 스킴이라 하자. $X$를 $S$ 위의 준콤팩트 준분리 대수공간이라 하자.
$\mathcal{F}$를 준연접 $\mathcal{O}_X$-가군이라 하자.
그러면 $\mathcal{F}$는 자신의 유한형 준연접 부분가군들의 유향 여극한이다.
\end{lemma}

\begin{proof}
$\mathcal{G}, \mathcal{H} \subset \mathcal{F}$가 유한형 준연접
$\mathcal{O}_X$-부분가군이면,
$\mathcal{G} \oplus \mathcal{H} \to \mathcal{F}$의 상은
이 둘을 모두 포함하는 또 다른 유한형 준연접 $\mathcal{O}_X$-부분가군이다.
이로써 이 계가 유향임을 알 수 있다.
$\mathcal{F}$가 이 계의 여극한임을 보이기 위해
$\mathcal{F} = \colim_i \mathcal{F}_i$를 유한 표시 준연접층의 유향 여극한으로 쓰자.
Lemma \ref{spaces-limits-lemma-colimit-finitely-presented}에서와 같이 쓴다.
그러면 상 $\mathcal{G}_i = \Im(\mathcal{F}_i \to \mathcal{F})$는
$\mathcal{F}$의 유한형 준연접 부분층이다.
$\mathcal{F}$가 이들의 여극한이므로 결과가 따른다.
\end{proof}

\begin{lemma}
\label{spaces-limits-lemma-finite-directed-colimit-surjective-maps}
$S$를 스킴이라 하자. $X$를 $S$ 위의 준콤팩트 준분리 대수공간이라 하자.
$\mathcal{F}$를 유한형 준연접 $\mathcal{O}_X$-가군이라 하자.
그러면 $\mathcal{F} = \lim \mathcal{F}_i$로 쓸 수 있다.
여기서 각 $\mathcal{F}_i$는 유한 표시 $\mathcal{O}_X$-가군이고,
모든 전이 사상 $\mathcal{F}_i \to \mathcal{F}_{i'}$는 전사이다.
\end{lemma}

\begin{proof}
$\mathcal{F} = \colim \mathcal{G}_i$를
유한 표시 $\mathcal{O}_X$-가군들의 여과 여극한으로 쓰자
(Lemma \ref{spaces-limits-lemma-colimit-finitely-presented}).
$\mathcal{G}_i \to \mathcal{F}$가 어떤 $i$에 대해 전사라고 주장한다.
실제로 에탈 전사 사상 $U \to X$를 택하되 $U$는 아핀 스킴이라고 하자.
유한 개의 단면 $s_k \in \mathcal{F}(U)$를 택하여
$\mathcal{F}|_U$를 생성하게 하자.
$U$가 아핀이므로 $s_k$는 $\mathcal{G}_i \to \mathcal{F}$의 상에
충분히 큰 $i$에 대해 속한다.
따라서 $\mathcal{G}_i \to \mathcal{F}$는 충분히 큰 $i$에 대해 전사이다.
그러한 $i$를 택하고 $\mathcal{K} \subset \mathcal{G}_i$를
사상 $\mathcal{G}_i \to \mathcal{F}$의 핵이라 하자.
$\mathcal{K} = \colim \mathcal{K}_a$를
그 유한형 준연접 부분가군들의 여과 여극한으로 쓰자
(Lemma \ref{spaces-limits-lemma-directed-colimit-finite-type}).
그러면 $\mathcal{F} = \colim \mathcal{G}_i/\mathcal{K}_a$가
보조정리에서 제시한 문제의 해이다.
\end{proof}

\noindent
$X$를 대수공간이라 하자. 다음 보조정리에서는
{\it 유한 표시 준연접 $\mathcal{O}_X$-대수 $\mathcal{A}$}라는 개념을 사용한다.
이는 모든 아핀 스킴 $U = \Spec(R)$가 $X$ 위에서 에탈일 때,
$\mathcal{A}|_U = \widetilde{A}$이고,
여기서 $A$는 (가환) $R$-대수로서 $R$-대수로서 유한 표시라는 뜻이다.

\begin{lemma}
\label{spaces-limits-lemma-algebra-directed-colimit-finite-presentation}
$S$를 스킴이라 하자. $X$를 $S$ 위의 준콤팩트 준분리 대수공간이라 하자.
$\mathcal{A}$를 준연접 $\mathcal{O}_X$-대수라 하자.
그러면 $\mathcal{A}$는 유한 표시 준연접 $\mathcal{O}_X$-대수들의 유향 여극한이다.
\end{lemma}

\begin{proof}
먼저 $\mathcal{A} = \colim_i \mathcal{F}_i$를
Lemma \ref{spaces-limits-lemma-colimit-finitely-presented}에서와 같이
유한 표시 준연접층의 유향 여극한으로 쓰자.
각 $i$에 대해 $\mathcal{B}_i = \text{Sym}(\mathcal{F}_i)$를
$\mathcal{F}_i$의 $\mathcal{O}_X$ 위의 대칭대수라 하자.
$\mathcal{I}_i = \Ker(\mathcal{B}_i \to \mathcal{A})$로 쓰자.
$\mathcal{I}_i = \colim_j \mathcal{F}_{i, j}$로 쓰되,
$\mathcal{F}_{i, j}$는 $\mathcal{I}_i$의 유한형 준연접 부분가군이라고 하자.
Lemma \ref{spaces-limits-lemma-directed-colimit-finite-type}를 보라.
$\mathcal{I}_{i, j} \subset \mathcal{I}_i$를
$\mathcal{B}_i$의 아이디얼 중 $\mathcal{F}_{i, j}$가 생성하는 것으로 놓자.
$\mathcal{A}_{i, j} = \mathcal{B}_i/\mathcal{I}_{i, j}$로 놓자.
그러면 $\mathcal{A}_{i, j}$는 준연접 유한 표시 $\mathcal{O}_X$-대수이다.
관계 $(i, j) \leq (i', j')$를 다음 조건으로 정의한다.
$i \leq i'$이고, 사상 $\mathcal{B}_i \to \mathcal{B}_{i'}$가
아이디얼 $\mathcal{I}_{i, j}$를 아이디얼 $\mathcal{I}_{i', j'}$ 안으로 보낸다.
그러면 $\mathcal{A} = \colim_{i, j} \mathcal{A}_{i, j}$임이 분명하다.
\end{proof}

\noindent
$X$를 대수공간이라 하자. 다음 보조정리에서는
{\it 준연접 $\mathcal{O}_X$-대수 $\mathcal{A}$가 유한형}이라는 개념을 사용한다.
이는 모든 아핀 스킴 $U = \Spec(R)$가 $X$ 위에서 에탈일 때,
$\mathcal{A}|_U = \widetilde{A}$이고,
여기서 $A$는 (가환) $R$-대수로서 $R$-대수로서 유한형이라는 뜻이다.

\begin{lemma}
\label{spaces-limits-lemma-algebra-directed-colimit-finite-type}
$S$를 스킴이라 하자. $X$를 $S$ 위의 준콤팩트 준분리 대수공간이라 하자.
$\mathcal{A}$를 준연접 $\mathcal{O}_X$-대수라 하자.
그러면 $\mathcal{A}$는 자신의 유한형 준연접
$\mathcal{O}_X$-부분대수들의 유향 여극한이다.
\end{lemma}

\begin{proof}
생략한다. 힌트: Lemma \ref{spaces-limits-lemma-directed-colimit-finite-type}의 증명과 비교하라.
\end{proof}

\noindent
$X$를 대수공간이라 하자. 다음 보조정리에서는
{\it 유한(각각 정수적) 준연접 $\mathcal{O}_X$-대수 $\mathcal{A}$}라는 개념을 사용한다.
이는 모든 아핀 스킴 $U = \Spec(R)$가 $X$ 위에서 에탈일 때,
$\mathcal{A}|_U = \widetilde{A}$이고,
여기서 $A$는 (가환) $R$-대수로서 $R$-대수로서 유한(각각 정수적)이라는 뜻이다.

\begin{lemma}
\label{spaces-limits-lemma-finite-algebra-directed-colimit-finite-finitely-presented}
$S$를 스킴이라 하자. $X$를 $S$ 위의 준콤팩트 준분리 대수공간이라 하자.
$\mathcal{A}$를 유한 준연접 $\mathcal{O}_X$-대수라 하자.
그러면 $\mathcal{A} = \colim \mathcal{A}_i$는
유한이면서 유한 표시인 준연접 $\mathcal{O}_X$-대수들의 유향 여극한으로,
전이 사상들은 전사이다.
\end{lemma}

\begin{proof}
Lemma \ref{spaces-limits-lemma-finite-directed-colimit-surjective-maps}에 의해
유한 표시 $\mathcal{O}_X$-가군 $\mathcal{F}$와 전사 사상
$\mathcal{F} \to \mathcal{A}$가 존재한다.
대수 구조를 사용하면 전사 사상
$$
\text{Sym}^*_{\mathcal{O}_X}(\mathcal{F}) \longrightarrow \mathcal{A}
$$
를 얻는다. 그 핵을 $\mathcal{J}$라 쓰자.
$\mathcal{J} = \colim \mathcal{E}_i$를
유한형 $\mathcal{O}_X$-부분가군 $\mathcal{E}_i$들의 여과 여극한으로 쓰자
(Lemma \ref{spaces-limits-lemma-directed-colimit-finite-type}).
$$
\mathcal{A}_i = \text{Sym}^*_{\mathcal{O}_X}(\mathcal{F})/(\mathcal{E}_i)
$$
로 놓자. 여기서 $(\mathcal{E}_i)$는
$\mathcal{E}_i \to \text{Sym}^*_{\mathcal{O}_X}(\mathcal{F})$의 상이 생성하는
아이디얼 층을 나타낸다.
그러면 각 $\mathcal{A}_i$는 유한 표시 $\mathcal{O}_X$-대수이고,
전이 사상들은 전사이며, $\mathcal{A} = \colim \mathcal{A}_i$이다.
증명을 끝내려면 $\mathcal{A}_i$가 유한 $\mathcal{O}_X$-대수임을
충분히 큰 $i$에 대해 보여야 한다.
이를 위해 에탈 전사 사상 $U \to X$를 택하되 $U$는 아핀 스킴이라고 하자.
생성원 $f_1, \ldots, f_m \in \Gamma(U, \mathcal{F})$를 택하자.
$\mathcal{A}(U)$는 유한 $\mathcal{O}_X(U)$-대수이므로,
각 $j$에 대해 일계수 다항식 $P_j \in \mathcal{O}(U)[T]$가 존재하여
$P_j(f_j)$가 $\mathcal{A}(U)$에서 영이 됨을 알 수 있다.
구성에 의해 $\mathcal{A} = \colim \mathcal{A}_i$이므로,
$P_j(f_j) = 0$이 $\mathcal{A}_i(U)$에서 충분히 큰 모든 $i$에 대해 성립한다.
그러한 $i$에 대해 대수 $\mathcal{A}_i$는 유한이다.
\end{proof}

\begin{lemma}
\label{spaces-limits-lemma-integral-algebra-directed-colimit-finite}
$S$를 스킴이라 하자. $X$를 $S$ 위의 준콤팩트 준분리 대수공간이라 하자.
$\mathcal{A}$를 정수적 준연접 $\mathcal{O}_X$-대수라 하자.
그러면 다음이 성립한다.
\begin{enumerate}
\item $\mathcal{A}$는 자신의 유한 준연접
$\mathcal{O}_X$-부분대수들의 유향 여극한이다.
\item $\mathcal{A}$는 유한이면서 유한 표시인
$\mathcal{O}_X$-대수들의 유향 여극한이다.
\end{enumerate}
\end{lemma}

\begin{proof}
Lemma \ref{spaces-limits-lemma-algebra-directed-colimit-finite-type}에 의해
$\mathcal{A} = \colim \mathcal{A}_i$이다.
여기서 $\mathcal{A}_i \subset \mathcal{A}$는
유한형 준연접 $\mathcal{O}_X$-부분대수들을 움직인다.
모든 유한형 준연접 $\mathcal{O}_X$-부분대수는
$\mathcal{A}$ 안에서 유한이다(Algebra, Lemma
\ref{algebra-lemma-characterize-finite-in-terms-of-integral}을
$X$ 위에서 에탈인 아핀 스킴들에 적용하라).
이로써 (1)이 증명된다.

\medskip\noindent
(2)를 증명하기 위해 $\mathcal{A} = \colim \mathcal{F}_i$를
유한 표시 $\mathcal{O}_X$-가군들의 여극한으로 쓰자.
Lemma \ref{spaces-limits-lemma-colimit-finitely-presented}를 사용한다.
각 $i$에 대해 $\mathcal{J}_i$를 사상
$$
\text{Sym}^*_{\mathcal{O}_X}(\mathcal{F}_i) \longrightarrow \mathcal{A}
$$
의 핵이라 하자.
$i' \geq i$에 대해 유도되는 사상 $\mathcal{J}_i \to \mathcal{J}_{i'}$가 있으며,
$\mathcal{A} =
\colim \text{Sym}^*_{\mathcal{O}_X}(\mathcal{F}_i)/\mathcal{J}_i$이다.
또한 준연접 $\mathcal{O}_X$-대수
$\text{Sym}^*_{\mathcal{O}_X}(\mathcal{F}_i)/\mathcal{J}_i$는 유한이다(위의 논의를 보라).
$\mathcal{J}_i = \colim \mathcal{E}_{ik}$를
유한 표시 $\mathcal{O}_X$-가군들의 여극한으로 쓰자.
$i' \geq i$와 $k$가 주어지면 $k'$가 존재하여
사상 $\mathcal{E}_{ik} \to \mathcal{E}_{i'k'}$를 다음 가환도가 가환하게 택할 수 있다.
$$
\xymatrix{
\mathcal{J}_i \ar[r] & \mathcal{J}_{i'} \\
\mathcal{E}_{ik} \ar[u] \ar[r] & \mathcal{E}_{i'k'} \ar[u]
}
$$
이는 Cohomology of Spaces, Lemma
\ref{spaces-cohomology-lemma-finite-presentation-quasi-compact-colimit}에서 따른다.
이 사상은 다음 사상을 유도한다.
$$
\mathcal{A}_{ik} =
\text{Sym}^*_{\mathcal{O}_X}(\mathcal{F}_i)/(\mathcal{E}_{ik})
\longrightarrow
\text{Sym}^*_{\mathcal{O}_X}(\mathcal{F}_{i'})/(\mathcal{E}_{i'k'}) =
\mathcal{A}_{i'k'}
$$
여기서 $(\mathcal{E}_{ik})$는 $\mathcal{E}_{ik}$가 생성하는 아이디얼을 나타낸다.
준연접 $\mathcal{O}_X$-대수 $\mathcal{A}_{ki}$는
충분히 큰 $k$에 대해 유한 표시이고 유한이다
(Lemma \ref{spaces-limits-lemma-finite-algebra-directed-colimit-finite-finitely-presented}의 증명 참조).
마지막으로
$$
\colim \mathcal{A}_{ik} = \colim \mathcal{A}_i = \mathcal{A}
$$
이다. 실제로 첫 번째 등식은
Lemma \ref{spaces-limits-lemma-finite-algebra-directed-colimit-finite-finitely-presented}의 증명에서 보였으며,
두 번째 등식은 $\mathcal{A}$가 가군 $\mathcal{F}_i$들의 여극한이기 때문에 성립한다.
\end{proof}

\begin{lemma}
\label{spaces-limits-lemma-extend}
$S$를 스킴이라 하자.
$X$를 $S$ 위의 준콤팩트 준분리 대수공간이라 하자.
$U \subset X$를 준콤팩트 열린집합이라 하자.
$\mathcal{F}$를 준연접 $\mathcal{O}_X$-가군이라 하자.
$\mathcal{G} \subset \mathcal{F}|_U$를 유한형 준연접
$\mathcal{O}_U$-부분가군이라 하자.
그러면 준연접 부분가군 $\mathcal{G}' \subset \mathcal{F}$가 존재하여
유한형이고 $\mathcal{G}'|_U = \mathcal{G}$를 만족한다.
\end{lemma}

\begin{proof}
포함 사상을 $j : U \to X$라 쓰자.
$X$가 준분리이고 $U$가 준콤팩트이므로 사상 $j$는 준콤팩트이다.
따라서 $j_*\mathcal{G} \subset j_*\mathcal{F}|_U$는
$X$ 위의 준연접 가군들이다
(Morphisms of Spaces, Lemma \ref{spaces-morphisms-lemma-pushforward}).
$\mathcal{H} =
\Ker(j_*\mathcal{G} \oplus \mathcal{F} \to j_*\mathcal{F}|_U)$라 하자.
그러면 $\mathcal{H}|_U = \mathcal{G}$이다.
Lemma \ref{spaces-limits-lemma-directed-colimit-finite-type}에 의해
유한형 준연접 부분가군 $\mathcal{H}' \subset \mathcal{H}$를 찾아
$\mathcal{H}'|_U = \mathcal{H}|_U = \mathcal{G}$가 되게 할 수 있다.
$\mathcal{G}' = \Im(\mathcal{H}' \to \mathcal{F})$로 놓으면 결론을 얻는다.
\end{proof}



\section{상대 근사}
\label{spaces-limits-section-relative-approximation}

\noindent
Proposition \ref{spaces-limits-proposition-approximate}을 밑 위에서 일반화한 형태들을 논의한다.

\begin{lemma}
\label{spaces-limits-lemma-approximate-morphism}
$f : X \to Y$를 $\mathbf{Z}$ 위의 준콤팩트 준분리 대수공간의 사상이라 하자.
그러면 유향 집합 $I$와 대수공간의 사상들의 역계 $(f_i : X_i \to Y_i)$가
$I$ 위에 존재하여 다음을 만족한다.
% P09-ERR-0084; P09-ERR-0085: source grammar normalized by translation only.
전이 사상 $X_i \to X_{i'}$ 및 $Y_i \to Y_{i'}$는 아핀이고,
$X_i$와 $Y_i$는 준분리이며 $\mathbf{Z}$ 위에서 유한형이고,
$(X \to Y) = \lim (X_i \to Y_i)$이다.
\end{lemma}

\begin{proof}
Proposition \ref{spaces-limits-proposition-approximate}에서와 같이
$X = \lim_{a \in A} X_a$ 및 $Y = \lim_{b \in B} Y_b$로 쓰자.
즉 $X_a$와 $Y_b$는 준분리이고 $\mathbf{Z}$ 위에서 유한형이며,
전이 사상들은 아핀이다.

\medskip\noindent
$b \in B$를 고정하자.
Lemma \ref{spaces-limits-lemma-better-characterize-relative-limit-preserving}을
$Y_b$와 $X = \lim X_a$에 $\mathbf{Z}$ 위에서 적용하면,
$a \in A$와 사상 $f_{a, b} : X_a \to Y_b$가 존재하여 가환도
$$
\xymatrix{
X \ar[d] \ar[r] & Y \ar[d] \\
X_a \ar[r] & Y_b
}
$$
가 가환하게 됨을 알 수 있다.
$I$를 이 방법으로 얻은 삼중항 $(a, b, f_{a, b})$들의 집합이라 하자.

\medskip\noindent
$(a, b, f_{a, b})$와 $(a', b', f_{a', b'})$가 $I$에 속한다고 하자.
$b'' \leq \min(b, b')$로 놓자.
% P09-ERR-0086: canonical ordering/minimum expression retained.
Lemma \ref{spaces-limits-lemma-better-characterize-relative-limit-preserving}을 다시 적용하면,
$a'' \geq \max(a, a')$가 존재하여 합성
$X_{a''} \to X_a \to Y_b \to Y_{b''}$와
$X_{a''} \to X_{a'} \to Y_{b'} \to Y_{b''}$가 같게 된다.
$I$에 다음 전순서를 준다.
$$
(a, b, f_{a, b}) \geq (a', b', f_{a', b'})
\Leftrightarrow
a \geq a',\ b \geq b',\text{ 및 }
g_{b, b'} \circ f_{a, b} = f_{a', b'} \circ h_{a, a'}
$$
여기서 $h_{a, a'} : X_a \to X_{a'}$ 및
$g_{b, b'} : Y_b \to Y_{b'}$는 전이 사상이다.
위의 관찰에서 $I$가 유향이고, 사상
$I \to A$, $(a, b, f_{a, b}) \mapsto a$ 및
$I \to B$, $(a, b, f_{a, b})$가 공종임을 알 수 있다.
% P09-ERR-0087: incomplete second projection retained without inventing mapsto b.
$i = (a, b, f_{a, b})$에 대해
$X_i = X_a$, $Y_i = Y_b$, $f_i = f_{a, b}$로 놓으면,
$I$ 위의 사상들의 역계를 얻고 다음이 성립한다.
$$
\lim_{i \in I} X_i = \lim_{a \in A} X_a = X
\quad\text{및}\quad
\lim_{i \in I} S_i = \lim_{b \in B} Y_b = Y
$$
% P09-ERR-0088: undefined S_i retained from frozen canonical source.
이는 Categories, Lemma \ref{categories-lemma-initial}에 따른다
($I$ 위의 극한은 실제로 $I$에 대응하는 반대 범주 위의 극한이므로,
공종은 공시작으로 바뀐다는 점을 상기하라).
이로써 증명이 끝난다.
\end{proof}







\begin{lemma}
\label{spaces-limits-lemma-relative-approximation}
$S$를 스킴이라 하자. $f : X \to Y$를 $S$ 위의 대수공간의 사상이라 하자.
다음을 가정하자.
\begin{enumerate}
\item $X$는 준콤팩트이고 준분리이다.
\item $Y$는 준분리이다.
\end{enumerate}
그러면 $X = \lim X_i$는 대수공간 $X_i$의 유향 역계의 극한으로 쓸 수 있으며,
각 대상은 $Y$ 위에서 유한 표시이고 전이 사상들은 $Y$ 위에서 아핀이다.
\end{lemma}

\begin{proof}
$|f|(|X|)$가 준콤팩트이므로 $Y$를 열린 부분공간으로 대체할 수 있다.
그 열린 부분공간은 준콤팩트이고 그 점들의 집합은 $|f|(|X|)$를 포함한다.
따라서 $Y$도 준콤팩트라고 가정할 수 있다.
Lemma \ref{spaces-limits-lemma-approximate-morphism}에 의해
$(X \to Y) = \lim (X_i \to Y_i)$로 쓸 수 있다.
여기서 어떤 유향 역계가 주어지고, 그 대상들은
$\mathbf{Z}$ 위의 유한형 스킴의 사상들이며 전이 사상들은 아핀이다.
% P09-ERR-0089: source calls these schemes though the cited lemma gives spaces.
극한은 극한과 교환하므로
(Categories, Lemma \ref{categories-lemma-colimits-commute}),
$X = \lim X_i \times_{Y_i} Y$이다.
$i \geq i'$에 대해 전이 사상
$X_i \times_{Y_i} Y \to X_{i'} \times_{Y_{i'}} Y$는 아핀이다.
실제로 이는 합성
$$
X_i \times_{Y_i} Y \to X_i \times_{Y_{i'}} Y \to X_{i'} \times_{Y_{i'}} Y
$$
이며, 첫 사상은 닫힌 몰입이다
(Morphisms of Spaces, Lemma \ref{spaces-morphisms-lemma-fibre-product-after-map}).
둘째 사상은 아핀 사상의 밑변환이다
(Morphisms of Spaces, Lemma \ref{spaces-morphisms-lemma-base-change-affine}).
또한 아핀 사상의 합성은 아핀이다
(Morphisms of Spaces, Lemma \ref{spaces-morphisms-lemma-composition-affine}).
사상 $f_i$는 유한 표시이다
(Morphisms of Spaces, Lemmas
\ref{spaces-morphisms-lemma-noetherian-finite-type-finite-presentation} 및
\ref{spaces-morphisms-lemma-finite-presentation-permanence}).
따라서 밑변환 $X_i \times_{f_i, Y_i} Y \to Y$도 유한 표시이다
(Morphisms of Spaces, Lemma
\ref{spaces-morphisms-lemma-base-change-finite-presentation}).
\end{proof}



\section{유한 표시 안에 닫혀 있는 유한형}
\label{spaces-limits-section-finite-type-closed-in-finite-presentation}

\noindent
이 절은 Limits, Section
\ref{limits-section-finite-type-closed-in-finite-presentation}에 대응한다.

\begin{lemma}
\label{spaces-limits-lemma-affine-morphism-is-limit}
$S$를 스킴이라 하자. $f : X \to Y$를 $S$ 위의 대수공간의 아핀 사상이라 하자.
$Y$가 준콤팩트이고 준분리이면, $X$는 유향 극한 $X = \lim X_i$이며
각 $X_i$는 $Y$ 위에서 아핀이고 유한 표시이다.
% P09-ERR-0090: omitted English copula is naturally supplied in Korean.
\end{lemma}

\begin{proof}
준연접 $\mathcal{O}_Y$-가군 $\mathcal{A} = f_*\mathcal{O}_X$를 생각하자.
Lemma \ref{spaces-limits-lemma-algebra-directed-colimit-finite-presentation}에 의해
$\mathcal{A} = \colim \mathcal{A}_i$를 유한 표시
$\mathcal{O}_Y$-대수 $\mathcal{A}_i$들의 유향 여극한으로 쓸 수 있다.
$X_i = \underline{\Spec}_Y(\mathcal{A}_i)$로 놓자.
Morphisms of Spaces, Definition \ref{spaces-morphisms-definition-relative-spec}을 보라.
구성에 의해 $X_i \to Y$는 아핀이고 유한 표시이며, $X = \lim X_i$이다.
\end{proof}

\begin{lemma}
\label{spaces-limits-lemma-integral-limit-finite-and-finite-presentation}
$S$를 스킴이라 하자. $f : X \to Y$를 $S$ 위의 대수공간의 정수적 사상이라 하자.
$Y$가 준콤팩트이고 준분리라고 가정하자.
그러면 $X$는 유향 극한 $X = \lim X_i$로 쓸 수 있고,
$X_i$는 $Y$ 위에서 유한이고 유한 표시이다.
\end{lemma}

\begin{proof}
준연접 $\mathcal{O}_Y$-가군 $\mathcal{A} = f_*\mathcal{O}_X$를 생각하자.
Lemma \ref{spaces-limits-lemma-integral-algebra-directed-colimit-finite}에 의해
$\mathcal{A} = \colim \mathcal{A}_i$를
유한이면서 유한 표시인 $\mathcal{O}_Y$-대수 $\mathcal{A}_i$들의
유향 여극한으로 쓸 수 있다.
$X_i = \underline{\Spec}_Y(\mathcal{A}_i)$로 놓자.
Morphisms of Spaces, Definition \ref{spaces-morphisms-definition-relative-spec}을 보라.
구성에 의해 $X_i \to Y$는 유한이고 유한 표시이며, $X = \lim X_i$이다.
\end{proof}

\begin{lemma}
\label{spaces-limits-lemma-finite-in-finite-and-finite-presentation}
$S$를 스킴이라 하자. $f : X \to Y$를 $S$ 위의 대수공간의 유한 사상이라 하자.
$Y$가 준콤팩트이고 준분리라고 가정하자.
그러면 $X$는 유향 극한 $X = \lim X_i$로 쓸 수 있고,
전이 사상들은 닫힌 몰입이며,
대상 $X_i$는 $Y$ 위에서 유한이고 유한 표시이다.
\end{lemma}

\begin{proof}
유한 준연접 $\mathcal{O}_Y$-가군 $\mathcal{A} = f_*\mathcal{O}_X$를 생각하자.
Lemma \ref{spaces-limits-lemma-finite-algebra-directed-colimit-finite-finitely-presented}에 의해
$\mathcal{A} = \colim \mathcal{A}_i$를
유한이면서 유한 표시인 $\mathcal{O}_Y$-대수 $\mathcal{A}_i$들의 유향 여극한으로
쓸 수 있고 전이 사상들은 전사이다.
$X_i = \underline{\Spec}_Y(\mathcal{A}_i)$로 놓자.
Morphisms of Spaces, Definition \ref{spaces-morphisms-definition-relative-spec}을 보라.
구성에 의해 $X_i \to Y$는 유한이고 유한 표시이며,
전이 사상들은 닫힌 몰입이고, $X = \lim X_i$이다.
\end{proof}

\begin{lemma}
\label{spaces-limits-lemma-closed-is-limit-closed-and-finite-presentation}
\begin{slogan}
qcqs 대수공간의 닫힌 몰입은 유한 표시 닫힌 몰입으로 근사할 수 있다.
\end{slogan}
$S$를 스킴이라 하자. $f : X \to Y$를 $S$ 위의 대수공간의 닫힌 몰입이라 하자.
$Y$가 준콤팩트이고 준분리라고 가정하자.
그러면 $X$는 유향 극한 $X = \lim X_i$로 쓸 수 있고,
전이 사상들은 닫힌 몰입이며,
사상 $X_i \to Y$는 유한 표시 닫힌 몰입이다.
\end{lemma}

\begin{proof}
$\mathcal{I} \subset \mathcal{O}_Y$를 $X$를 $Y$의 닫힌 부분공간으로 정의하는
준연접 아이디얼 층이라 하자.
Lemma \ref{spaces-limits-lemma-directed-colimit-finite-type}에 의해
$\mathcal{I} = \colim \mathcal{I}_i$를
그 유한형 준연접 부분가군들의 여과 여극한으로 쓸 수 있다.
$X_i$를 $X$의 닫힌 부분공간 가운데 $\mathcal{I}_i$가 잘라내는 것으로 하자.
% P09-ERR-0091: source's ambient X retained; I_i lives on Y.
그러면 $X_i \to Y$는 유한 표시 닫힌 몰입이고 $X = \lim X_i$이다.
몇 가지 세부 사항은 생략한다.
\end{proof}

\begin{lemma}
\label{spaces-limits-lemma-quasi-affine-closed-in-finite-presentation}
$S$를 스킴이라 하자. $f : X \to Y$를 $S$ 위의 대수공간의 사상이라 하자.
다음을 가정하자.
\begin{enumerate}
\item $f$는 국소 유한형이고 준아핀이다.
\item $Y$는 준콤팩트이고 준분리이다.
\end{enumerate}
그러면 유한 표시 사상 $f' : X' \to Y$와
닫힌 몰입 $X \to X'$가 존재하며, 후자는 $Y$ 위의 사상이다.
\end{lemma}

\begin{proof}
Morphisms of Spaces, Lemma
\ref{spaces-morphisms-lemma-characterize-quasi-affine}에 의해
분해 $X \to Z \to Y$를 찾을 수 있다.
여기서 $X \to Z$는 준콤팩트 열린 몰입이고 $Z \to Y$는 아핀이다.
$Z = \lim Z_i$로 쓰되 $Z_i$는 $Y$ 위에서 아핀이고 유한 표시라고 하자
(Lemma \ref{spaces-limits-lemma-affine-morphism-is-limit}).
어떤 $0 \in I$에 대해 준콤팩트 열린집합 $U_0 \subset Z_0$를 찾아
$X$가 $U_0$의 역상과 동형이 되게 할 수 있는데, 이 역상은 $Z$ 안에서 취한다
(Lemma \ref{spaces-limits-lemma-descend-opens}).
$U_i$를 $U_0$의 역상이라 하되 그 역상은 $Z_i$ 안에서 취하면 $U = \lim U_i$이다.
% P09-ERR-0092: undefined U in source equality retained.
Lemma \ref{spaces-limits-lemma-finite-type-eventually-closed}에 의해
$X \to U_i$는 충분히 큰 어떤 $i$에 대해 닫힌 몰입이다.
$X' = U_i$로 놓으면 증명이 끝난다.
\end{proof}

\begin{lemma}
\label{spaces-limits-lemma-finite-type-closed-in-finite-presentation}
$S$를 스킴이라 하자. $f : X \to Y$를 $S$ 위의 대수공간의 사상이라 하자.
다음을 가정하자.
\begin{enumerate}
\item $f$는 국소 유한형이다.
% P09-ERR-0093: duplicated English preposition normalized by translation only.
\item $X$는 준콤팩트이고 준분리이다.
\item $Y$는 준콤팩트이고 준분리이다.
\end{enumerate}
그러면 유한 표시 사상 $f' : X' \to Y$와
대수공간의 닫힌 몰입 $X \to X'$가 존재하며, 후자는 $Y$ 위의 사상이다.
\end{lemma}

\begin{proof}
Proposition \ref{spaces-limits-proposition-approximate}에 의해
$X = \lim_i X_i$로 쓸 수 있다.
여기서 $X_i$는 준분리이고 $\mathbf{Z}$ 위에서 유한형이며,
전이 사상 $f_{ii'} : X_i \to X_{i'}$는 아핀이다.
다음 가환도를 생각하자.
$$
\xymatrix{
X \ar[r] \ar[rd] & X_{i, Y} \ar[r] \ar[d] & X_i \ar[d] \\
& Y \ar[r] & \Spec(\mathbf{Z})
}
$$
$X_i$는 $\Spec(\mathbf{Z})$ 위에서 유한 표시임에 주의하자.
Morphisms of Spaces, Lemma
\ref{spaces-morphisms-lemma-noetherian-finite-type-finite-presentation}을 보라.
따라서 밑변환 $X_{i, Y} \to Y$는
Morphisms of Spaces, Lemma
\ref{spaces-morphisms-lemma-base-change-finite-presentation}에 의해 유한 표시이다.
$\lim X_{i, Y} = X \times Y$이고
$X \to X \times Y$가 단사상임을 관찰하자.
% P09-ERR-0094: missing base in both source products retained.
Lemma \ref{spaces-limits-lemma-finite-type-eventually-closed}에 의해
$X \to X_{i, Y}$는 충분히 큰 $i$에 대해 단사상이다.
그러한 $i$를 고정하자.
$X \to X_{i, Y}$는 국소 유한형이고
(Morphisms of Spaces, Lemma \ref{spaces-morphisms-lemma-permanence-finite-type}),
단사상이므로 분리이고 국소 준유한이다
(Morphisms of Spaces, Lemma
\ref{spaces-morphisms-lemma-monomorphism-loc-finite-type-loc-quasi-finite}).
따라서 $X \to X_{i, Y}$는 표현 가능하다.
그러므로 $X \to X_{i, Y}$는 준아핀이다.
실제로 Spaces, Lemma
\ref{spaces-lemma-representable-transformations-property-implication}의 원리와
스킴의 사상에 대한 결과 More on Morphisms, Lemma
\ref{more-morphisms-lemma-quasi-finite-separated-quasi-affine}를 사용할 수 있다.
따라서 Lemma \ref{spaces-limits-lemma-quasi-affine-closed-in-finite-presentation}은
분해 $X \to X' \to X_{i, Y}$를 준다.
여기서 $X \to X'$는 닫힌 몰입이고
$X' \to X_{i, Y}$는 유한 표시이다.
마지막으로 $X' \to Y$는 유한 표시 사상의 합성이므로 유한 표시이다
(Morphisms of Spaces, Lemma
\ref{spaces-morphisms-lemma-composition-finite-presentation}).
\end{proof}

\begin{proposition}
\label{spaces-limits-proposition-separated-closed-in-finite-presentation}
$S$를 스킴이라 하자. $f : X \to Y$를 $S$ 위의 대수공간의 사상이라 하자.
% P09-ERR-0095: missing English Let supplied by Korean syntax only.
다음을 가정하자.
\begin{enumerate}
\item $f$는 유한형이고 분리이다.
\item $Y$는 준콤팩트이고 준분리이다.
\end{enumerate}
그러면 유한 표시인 분리 사상 $f' : X' \to Y$와
닫힌 몰입 $X \to X'$가 존재하며, 후자는 $Y$ 위의 사상이다.
\end{proposition}

\begin{proof}
Lemma \ref{spaces-limits-lemma-finite-type-closed-in-finite-presentation}에 의해
닫힌 몰입 $X \to Z$가 존재하고 $Z/Y$는 유한 표시이다.
$\mathcal{I} \subset \mathcal{O}_Z$를 $X$를 $Y$의 닫힌 부분스킴으로 정의하는
준연접 아이디얼 층이라 하자.
% P09-ERR-0096: source's ambient Y/subscheme wording retained.
Lemma \ref{spaces-limits-lemma-directed-colimit-finite-type}에 의해
$\mathcal{I}$를 유향 여극한
$\mathcal{I} = \colim_{a \in A} \mathcal{I}_a$로 쓸 수 있다.
여기서 각 항은 그 유한형 준연접 아이디얼 층이다.
$X_a \subset Z$를 $\mathcal{I}_a$가 정의하는 닫힌 부분공간이라 하자.
이들은 $A$를 지표로 하는 역계를 이룬다.
전이 사상 $X_a \to X_{a'}$는 닫힌 몰입이므로 아핀이다.
각 $X_a$는 $Z$의 닫힌 부분공간이고,
가정에 의해 $Z$가 준콤팩트이고 준분리이므로 그것들도 준콤팩트이고 준분리이다.
$X = \lim_a X_a$가 성립한다. 이는
$\mathcal{I} = \colim_{a \in A} \mathcal{I}_a$라는 사실에서 직접 따른다.
각 사상 $X_a \to Z$는 유한 표시이다.
Morphisms, Lemma \ref{morphisms-lemma-closed-immersion-finite-presentation}을 보라.
따라서 사상 $X_a \to Y$는 유한 표시이다.
그러므로 $X_a \to Y$가 어떤 $a \in A$에 대해 분리임을 보이면 충분하다.
이는 $X \to Y$가 분리라는 가정과
Lemma \ref{spaces-limits-lemma-eventually-separated}에서 따른다.
\end{proof}





\section{고유 사상의 근사}
\label{spaces-limits-section-approximate-proper}

\begin{lemma}
\label{spaces-limits-lemma-proper-limit-of-proper-finite-presentation}
$S$를 스킴이라 하자. $f : X \to Y$를 $S$ 위의 대수공간의 고유 사상이라 하고,
$Y$가 준콤팩트이고 준분리라고 가정하자. 그러면
$X = \lim X_i$는 대수공간 $X_i$들의 유향 극한으로 쓸 수 있으며,
이들은 $Y$ 위에서 고유하고 유한 표시이고, 전이 사상들과 사상 $X \to X_i$는
닫힌 몰입이다.
\end{lemma}

\begin{proof}
Proposition \ref{spaces-limits-proposition-separated-closed-in-finite-presentation}에 의해
닫힌 몰입 $X \to X'$를 찾을 수 있으며, $X'$는 $Y$ 위에서 분리이고
유한 표시이다. Lemma \ref{spaces-limits-lemma-closed-is-limit-closed-and-finite-presentation}에 의해
$X = \lim X_i$로 쓸 수 있으며, $X_i \to X'$는 유한 표시 닫힌 몰입이다.
충분히 큰 모든 $i$에 대해 사상 $X_i \to Y$가 고유임을 보이면 증명이 끝난다.

\medskip\noindent
이를 보이기 위해 $Y$가 아핀 스킴이라고 가정해도 된다.
Morphisms of Spaces, Lemma \ref{spaces-morphisms-lemma-proper-local}을 보라.
이제 Cohomology of Spaces, Lemma
\ref{spaces-cohomology-lemma-weak-chow}의 Chow 보조정리의 약한 형태를 사용하여
다음 가환도를 얻는다.
$$
\xymatrix{
X' \ar[rd] & X'' \ar[d] \ar[l]^\pi \ar[r] & \mathbf{P}^n_Y \ar[dl] \\
& Y &
}
$$
여기서 $X'' \to \mathbf{P}^n_Y$는 몰입이고,
$\pi : X'' \to X'$는 고유 전사이다.
$X'_i \subset X''$와 $\pi^{-1}(X)$를 각각
$X_i \subset X'$와 $X \subset X'$의 스킴론적 역상이라 하자.
그러면 $\lim X'_i = \pi^{-1}(X)$이다. $\pi^{-1}(X) \to Y$가 고유이므로
(Morphisms of Spaces, Lemmas \ref{spaces-morphisms-lemma-composition-proper}),
$\pi^{-1}(X) \to \mathbf{P}^n_Y$는 닫힌 몰입이다
(Morphisms of Spaces, Lemmas
\ref{spaces-morphisms-lemma-universally-closed-permanence}와
\ref{spaces-morphisms-lemma-immersion-when-closed}).
따라서 충분히 큰 $i$에 대해
$X'_i \to \mathbf{P}^n_Y$는 Lemma \ref{spaces-limits-lemma-eventually-closed-immersion}에 의해
닫힌 몰입이다. 그러므로 $X'_i$는 $Y$ 위에서 고유이다.
그러한 $i$에 대해 사상 $X_i \to Y$는
Morphisms of Spaces, Lemma \ref{spaces-morphisms-lemma-image-proper-is-proper}에 의해
고유이다.
\end{proof}

\begin{lemma}
\label{spaces-limits-lemma-proper-limit-of-proper-finite-presentation-noetherian}
$f : X \to Y$를 $\mathbf{Z}$ 위의 대수공간의 고유 사상이라 하고,
$Y$가 준콤팩트이고 준분리라고 가정하자. 그러면 유향집합 $I$와
대수공간의 사상들로 이루어진 역계 $(f_i : X_i \to Y_i)$가 존재하며,
이 역계는 $I$ 위에 놓인다. 또한
전이 사상 $X_i \to X_{i'}$와 $Y_i \to Y_{i'}$는 아핀이고,
$f_i$는 고유하고 유한 표시이며, $Y_i$는 $\mathbf{Z}$ 위에서 유한 표시이고,
$(X \to Y) = \lim (X_i \to Y_i)$이다.
\end{lemma}

\begin{proof}
Lemma \ref{spaces-limits-lemma-proper-limit-of-proper-finite-presentation}에 의해
$X = \lim_{k \in K} X_k$로 쓸 수 있으며, $X_k \to Y$는 고유하고
유한 표시이다. 다음으로 절대 Noether 근사
(Proposition \ref{spaces-limits-proposition-approximate})에 의해
$Y = \lim_{j \in J} Y_j$로 쓸 수 있으며, $Y_j$는
$\mathbf{Z}$ 위에서 유한 표시이다.
각 $k$에 대해 어떤 $j$와 유한 표시 사상 $X_{k, j} \to Y_j$가 존재하여,
$X_k \cong Y \times_{Y_j} X_{k, j}$가 되고, 이 동형은 $Y$ 위의 대수공간 사이의 동형이다.
Lemma \ref{spaces-limits-lemma-descend-finite-presentation}을 보라.
$j$를 키우면 $X_{k, j} \to Y_j$가 고유라고 가정할 수 있다.
Lemma \ref{spaces-limits-lemma-eventually-proper}을 보라.
% P09-ERR-0097: frozen source grammar normalized by Korean syntax only.
$I$는 이 쌍 $(k, j)$들과 대응하는 사상 $X_{k, j} \to Y_j$로 이루어진다.
% P09-ERR-0098: frozen source's reversed X_{j',k'} and X_{j,k} indices retained.
모든 $k' \geq k$에 대해 어떤 $j' \geq j$와
사상 $X_{j', k'} \to X_{j, k}$를 찾을 수 있으며, 이는 $Y_{j'} \to Y_j$ 위의 사상이다.
이를 $Y$로 밑변환하면 사상 $X_{k'} \to X_k$가 된다
(다시 Lemma \ref{spaces-limits-lemma-descend-finite-presentation}에서 따른다).
이 사상들이 계의 전이 사상을 이룬다. 몇 가지 세부 사항은 생략한다.
\end{proof}

\noindent
유한형 준연접 가군의 스킴론적 지지를 상기하자.
Morphisms of Spaces, Definition
\ref{spaces-morphisms-definition-scheme-theoretic-support}을 보라.

\begin{lemma}
\label{spaces-limits-lemma-eventually-proper-support}
Situation \ref{spaces-limits-situation-descent-property}의 가정과 표기를 사용하자.
$\mathcal{F}_0$를 준연접 $\mathcal{O}_{X_0}$-가군이라 하자.
$\mathcal{F}$와 $\mathcal{F}_i$를 각각 $\mathcal{F}_0$의 $X$와 $X_i$로의
당김이라 하자. 다음을 가정하자.
\begin{enumerate}
\item $f_0$는 국소 유한형이다.
\item $\mathcal{F}_0$는 유한형이다.
\item $\mathcal{F}$의 스킴론적 지지는 $Y$ 위에서 고유이다.
\end{enumerate}
그러면 $\mathcal{F}_i$의 스킴론적 지지가
$Y_i$ 위에서 고유인 어떤 $i$가 존재한다.
\end{lemma}

\begin{proof}
$X_0$를 $\mathcal{F}_0$의 스킴론적 지지로 대체해도 된다.
Morphisms of Spaces, Lemma \ref{spaces-morphisms-lemma-support-finite-type}에 의해
이는 $X_i$가 $\mathcal{F}_i$의 지지이고 $X$가 $\mathcal{F}$의 지지임을 보장한다.
이제 $Z \subset X$가 $\mathcal{F}$의 스킴론적 지지를 나타낸다면,
$Z \to X$는 보편 위상동형이다. 따라서 $X \to Y$가 고유임을 얻는다.
실제로 가정에 의해 $Z \to Y$가 고유이며,
Morphisms, Lemma \ref{morphisms-lemma-image-proper-is-proper}을 적용한다.
% P09-ERR-0099: frozen source's Y targets retained instead of silently changing to Y_i.
Lemma \ref{spaces-limits-lemma-eventually-proper}에 의해 $X_i \to Y$가 고유인 어떤 $i$가 존재한다.
그러면 스킴론적 지지 $Z_i$, 즉 $\mathcal{F}_i$의 지지는 $Y$ 위에서 고유이다.
이는 Morphisms of Spaces, Lemmas
\ref{spaces-morphisms-lemma-closed-immersion-proper}와
\ref{spaces-morphisms-lemma-composition-proper}에서 따른다.
\end{proof}





\section{아핀 공간으로의 매장}
\label{spaces-limits-section-embedding}

\noindent
뒤에서 Chow 보조정리의 증명에 사용할 몇 가지 기술적 보조정리를 제시한다.

\begin{lemma}
\label{spaces-limits-lemma-embedding-into-affine-over-ls-qs}
$S$를 스킴이라 하자. $f : U \to X$를 $S$ 위의 대수공간의 사상이라 하자.
$U$가 아핀 스킴이고, $f$가 국소 유한형이며,
$X$가 준분리이고 국소 분리라고 가정하자.
그러면 몰입 $U \to \mathbf{A}^n_X$가 존재하며, 이는 $X$ 위의 사상이다.
\end{lemma}

\begin{proof}
$U = \Spec(A)$라 하자. $A = \colim A_i$를 유한형
$\mathbf{Z}$-부분대수들의 여과 여극한으로 쓰자. 각 $i$에 대해 사상
$U \to U_i = \Spec(A_i)$는 다음 사상을 유도한다.
$$
U \longrightarrow X \times U_i
$$
이 사상은 $X$ 위의 사상이다. 극한에서 사상 $U \to X \times U$는 몰입이다.
실제로 $X$가 국소 분리이며, Morphisms of Spaces, Lemma
\ref{spaces-morphisms-lemma-semi-diagonal}을 적용한다.
Lemma \ref{spaces-limits-lemma-finite-type-eventually-closed}에 의해
$U \to X \times U_i$가 어떤 $i$에 대해 몰입임을 알 수 있다.
$U_i$는 $\mathbf{A}^n_{\mathbf{Z}}$의 닫힌 부분스킴과 동형이므로
보조정리가 따른다.
\end{proof}

\begin{remark}
\label{spaces-limits-remark-cannot-embed-in-general}
Examples, Section \ref{examples-section-embedding-affines}에서 보았듯이,
Lemma \ref{spaces-limits-lemma-embedding-into-affine-over-ls-qs}는
$X$가 국소 분리라는 가정을 버리면 성립하지 않는다.
이에 따라 다음 질문이 생긴다. $X$가 준분리라는 가정을 버려도
Lemma \ref{spaces-limits-lemma-embedding-into-affine-over-ls-qs}가 성립하는가?
답을 알고 있다면
\href{mailto:stacks.project@gmail.com}{stacks.project@gmail.com}으로 전자우편을 보내 달라.
\end{remark}

\begin{lemma}
\label{spaces-limits-lemma-embedding-into-affine-over-qs}
$S$를 스킴이라 하자. $f : Y \to X$를 $S$ 위의 대수공간의 사상이라 하자.
$X$가 Noether이고 $f$가 유한 표시라고 가정하자.
그러면 조밀한 열린집합 $V \subset Y$와 몰입
$V \to \mathbf{A}^n_X$가 존재한다.
\end{lemma}

\begin{proof}
가정에 의해 $Y$는 Noether이다
(Morphisms of Spaces, Lemma
\ref{spaces-morphisms-lemma-finite-presentation-noetherian}).
그러면 $Y$는 준분리이므로 조밀한 열린 부분스킴을 가진다
(Properties of Spaces, Proposition
\ref{spaces-properties-proposition-locally-quasi-separated-open-dense-scheme}).
따라서 $Y$가 Noether 스킴이라고 가정해도 된다.
$Y$의 기약 성분들의 교차를 제거하면
(Topology, Lemma \ref{topology-lemma-Noetherian}과
Properties, Lemma \ref{properties-lemma-Noetherian-topology}을 사용한다),
$Y$가 기약 Noether 스킴들의 서로소 합이라고 가정해도 된다.
다음 몰입이 존재한다.
$$
\mathbf{A}^n_X \amalg \mathbf{A}^m_X
\longrightarrow
\mathbf{A}^{\max(n, m) + 1}_X
$$
(세부 사항은 생략한다). 따라서 $Y$가 기약인 경우만 증명하면 충분하다.

\medskip\noindent
$Y$가 기약 스킴이라고 가정하자. $T \subset |X|$를
$f : Y \to X$의 상의 폐포라 하자. $|Y|$와 $|X|$는 sober 위상공간이므로
(Properties of Spaces, Lemma
\ref{spaces-properties-lemma-quasi-separated-sober}),
$T$는 기약이고 유일한 일반점 $\xi$를 가지며, 이 점은
일반점 $\eta$의 상이고, 이 일반점은 $Y$의 일반점이다.
% P09-ERR-0100: frozen source's ideal-sheaf inclusion in X retained literally.
$\mathcal{I} \subset X$를 $T$ 위의 축소 유도 공간 구조를 잘라내는
준연접 아이디얼 층이라 하자
(Properties of Spaces, Definition
\ref{spaces-properties-definition-reduced-induced-space}).
$\mathcal{O}_{Y, \eta}$는 Artin 국소환이므로, 어떤 $n > 0$에 대해
$f^{-1}\mathcal{I}^n \mathcal{O}_{Y, \eta} = 0$이다.
$f^{-1}\mathcal{I}\mathcal{O}_Y$는 유한형 준연접 아이디얼이므로,
$f^{-1}\mathcal{I}^n\mathcal{O}_V = 0$인 공집합이 아닌 열린집합
$V \subset Y$가 존재한다.
$Z \subset X$를 $\mathcal{I}^n$이 잘라내는 닫힌 부분공간이라 하자.
구성에 의해 $V \to Y \to X$는 $Z$를 통해 인수분해된다.
$\mathbf{A}^n_Z \to \mathbf{A}^n_X$가 몰입이므로,
$X$를 $Z$로, $Y$를 $V$로 대체해도 된다.
따라서 $Y$와 $X$가 기약이고,
$Y \to X$가 $Y$의 일반점을 $X$의 일반점으로 보내는 상황에 이른다.

\medskip\noindent
$Y$와 $X$가 기약이고, $Y$가 스킴이며,
$Y \to X$가 $Y$의 일반점을 $X$의 일반점으로 보낸다고 가정하자.
Properties of Spaces, Proposition
\ref{spaces-properties-proposition-locally-quasi-separated-open-dense-scheme}에 의해
$X$는 조밀한 열린 부분스킴 $U \subset X$를 가진다.
공집합이 아닌 아핀 열린집합 $V \subset Y$를 택하되,
그 상이 $X$에서 $U$에 포함되게 하자. Morphisms, Lemma
\ref{morphisms-lemma-quasi-affine-finite-type-over-S}에 의해
$V \to U$를 $V \to \mathbf{A}^n_U \to U$로 인수분해할 수 있다.
이를 $\mathbf{A}^n_U \to \mathbf{A}^n_X$와 합성하면 원하는 몰입을 얻는다.
\end{proof}








\section{닫힌 부분집합에 지지를 갖는 절단}
\label{spaces-limits-section-sections-with-support-in-closed}

\noindent
이 절은 Properties, Section
\ref{properties-section-sections-with-support-in-closed}에 대응한다.

\begin{lemma}
\label{spaces-limits-lemma-quasi-coherent-finite-type-ideals}
$S$를 스킴이라 하자.
$X$를 준콤팩트 준분리 대수공간이라 하자.
$U \subset X$를 열린 부분공간이라 하자. 다음은 서로 동치이다.
\begin{enumerate}
\item $U \to X$는 준콤팩트이다.
\item $U$는 준콤팩트이다.
\item 유한형 준연접 아이디얼 층
$\mathcal{I} \subset \mathcal{O}_X$가 존재하여
$|X| \setminus |U| = |V(\mathcal{I})|$이다.
\end{enumerate}
\end{lemma}

\begin{proof}
$W$를 아핀 스킴이라 하고, $\varphi : W \to X$를 전사 에탈 사상이라 하자.
Properties of Spaces, Lemma
\ref{spaces-properties-lemma-quasi-compact-affine-cover}를 보라.
(1)이 성립하면 $\varphi^{-1}(U) \to W$는 준콤팩트이므로
$\varphi^{-1}(U)$는 준콤팩트이고, 따라서 $U$도 준콤팩트이다
($|\varphi^{-1}(U)| \to |U|$가 전사이기 때문이다).
(2)가 성립하면 $\varphi^{-1}(U)$는 준콤팩트이다. 실제로 $\varphi$는 준콤팩트하다.
이는 $X$가 준분리이기 때문이다
(Morphisms of Spaces, Lemma
\ref{spaces-morphisms-lemma-quasi-compact-quasi-separated-permanence}).
따라서 $\varphi^{-1}(U) \to W$는
Properties, Lemma \ref{properties-lemma-quasi-coherent-finite-type-ideals}에 의해
스킴의 준콤팩트 사상이다. 그러므로
$U \to X$는 Morphisms of Spaces, Lemma
\ref{spaces-morphisms-lemma-quasi-compact-local}에 의해 준콤팩트이다.
따라서 (1)과 (2)는 동치이다.

\medskip\noindent
(1)과 (2)를 가정하자. Properties of Spaces, Lemma
\ref{spaces-properties-lemma-reduced-closed-subspace}에 의해,
유일한 준연접 아이디얼 층 $\mathcal{J}$가 존재하며,
이는 $|X| \setminus |U|$ 위의 축소 유도 닫힌 부분공간 구조를 잘라낸다.
% P09-ERR-0101: frozen source's singular/plural grammar normalized in Korean.
$\mathcal{J}|_U = \mathcal{O}_U$이고, 이는 유한형
$\mathcal{O}_U$-가군임에 주의하자.
$U$가 준콤팩트이므로 Lemma \ref{spaces-limits-lemma-directed-colimit-finite-type}에 의해,
유한형 준연접 부분층 $\mathcal{I} \subset \mathcal{J}$가 존재하여
$\mathcal{I}|_U = \mathcal{J}|_U$가 된다.
그러면 $|X| \setminus |U| = |V(\mathcal{I})|$이므로 (3)을 얻는다.
반대로 $\mathcal{I}$가 (3)에서와 같다면,
$\varphi^{-1}(U) \subset W$는 스킴에 대한 보조정리
(Properties, Lemma \ref{properties-lemma-quasi-coherent-finite-type-ideals})에 의해
준콤팩트 열린집합이다. 여기서는 $\varphi^{-1}\mathcal{I}$에
그 보조정리를 적용하며, 이는 $W$ 위의 층이다. 따라서 (2)가 성립한다.
\end{proof}

\begin{lemma}
\label{spaces-limits-lemma-sections-annihilated-by-ideal}
$S$를 스킴이라 하자. $X$를 $S$ 위의 대수공간이라 하자.
$\mathcal{I} \subset \mathcal{O}_X$를 준연접 아이디얼 층이라 하자.
$\mathcal{F}$를 준연접 $\mathcal{O}_X$-가군이라 하자.
$\mathcal{O}_X$-가군의 층 $\mathcal{F}'$를 생각하자. 이 층은 모든 대상
$U$에 다음 가군을 대응시키며, 그 대상은 $X_\etale$에 속한다.
$$
\mathcal{F}'(U)
=
\{s \in \mathcal{F}(U) \mid
\mathcal{I}s = 0\}
$$
$\mathcal{I}$가 유한형이라고 가정하자. 그러면
\begin{enumerate}
\item $\mathcal{F}'$는 준연접 $\mathcal{O}_X$-가군의 층이다.
\item 아핀 대상 $U$가 $X_\etale$에 속하면
$\mathcal{F}'(U) = \{s \in \mathcal{F}(U) \mid \mathcal{I}(U)s = 0\}$이다.
\item $\mathcal{F}'_x = \{s \in \mathcal{F}_x \mid \mathcal{I}_x s = 0\}$이다.
\end{enumerate}
\end{lemma}

\begin{proof}
$\mathcal{F}'$를 정의하는 규칙이 $\mathcal{F}$의 부분층을 준다는 것은 명백하다.
따라서 나머지 명제를 확인할 때 $X$ 위에서 에탈 국소적으로 작업해도 된다.
그러므로 보조정리는 스킴의 경우인 Properties, Lemma
\ref{properties-lemma-sections-annihilated-by-ideal}로 환원된다.
\end{proof}

\begin{definition}
\label{spaces-limits-definition-subsheaf-sections-annihilated-by-ideal}
$S$를 스킴이라 하자. $X$를 $S$ 위의 대수공간이라 하자.
$\mathcal{I} \subset \mathcal{O}_X$를 유한형 준연접 아이디얼 층이라 하자.
$\mathcal{F}$를 준연접 $\mathcal{O}_X$-가군이라 하자.
위의 Lemma \ref{spaces-limits-lemma-sections-annihilated-by-ideal}에서 정의한 부분층
$\mathcal{F}' \subset \mathcal{F}$를
{\it $\mathcal{I}$에 의해 소멸되는 절단들의 부분층}이라 한다.
\end{definition}

\begin{lemma}
\label{spaces-limits-lemma-push-sections-annihilated-by-ideal}
$S$를 스킴이라 하자.
$f : X \to Y$를 $S$ 위의 대수공간의 준콤팩트 준분리 사상이라 하자.
$\mathcal{I} \subset \mathcal{O}_Y$를 유한형 준연접 아이디얼 층이라 하자.
$\mathcal{F}$를 준연접 $\mathcal{O}_X$-가군이라 하자.
$\mathcal{F}' \subset \mathcal{F}$를
$f^{-1}\mathcal{I}\mathcal{O}_X$에 의해 소멸되는 절단들의 부분층이라 하자.
그러면 $f_*\mathcal{F}' \subset f_*\mathcal{F}$는
$\mathcal{I}$에 의해 소멸되는 절단들의 부분층이다.
\end{lemma}

\begin{proof}
생략한다. 힌트: $f$가 준콤팩트이고 준분리라는 가정에 의해
$f_*\mathcal{F}$는 준연접이다
(Morphisms of Spaces, Lemma \ref{spaces-morphisms-lemma-pushforward}).
따라서 Lemma \ref{spaces-limits-lemma-sections-annihilated-by-ideal}을
$\mathcal{I}$와 $f_*\mathcal{F}$에 적용할 수 있다.
\end{proof}

\noindent
이제 닫힌 부분집합에 지지를 갖는 절단들의 층을 살펴보자.
이 층 역시 언제나 준연접인 것은 아니지만, 주어진 대수공간에서
닫힌 부분집합의 여집합이 ``retrocompact''이면 준연접이다.

\begin{lemma}
\label{spaces-limits-lemma-sections-supported-on-closed-subset}
$S$를 스킴이라 하자. $X$를 $S$ 위의 대수공간이라 하자.
$T \subset |X|$를 닫힌 부분집합이라 하고, $U \subset X$를
$T \amalg |U| = |X|$가 되게 하는 열린 부분공간이라 하자.
$\mathcal{F}$를 준연접 $\mathcal{O}_X$-가군이라 하자.
$\mathcal{O}_X$-가군의 층 $\mathcal{F}'$를 생각하자. 이 층은 각 대상
$\varphi : W \to X$에 다음 가군을 대응시키며, 그 대상은 $X_\etale$에 속한다.
$$
\mathcal{F}'(W)
=
\{s \in \mathcal{F}(W) \mid
\text{절단 }s\text{의 지지가 다음 집합에 포함됨: }|\varphi|^{-1}(T)\}
$$
$U \to X$가 준콤팩트이면 다음이 성립한다.
\begin{enumerate}
% P09-ERR-0102: frozen source's subject-verb disagreement normalized in Korean.
\item $W$가 아핀이면 유한 생성 아이디얼
$I \subset \mathcal{O}_X(W)$가 존재하여 $|\varphi|^{-1}(T) = V(I)$가 된다.
\item (1)에서와 같은 $W$와 $I$에 대해
$\mathcal{F}'(W) = \{x \in \mathcal{F}(W) \mid
I^nx = 0 \text{인 어떤 } n\}$이다.
\item $\mathcal{F}'$는 준연접 $\mathcal{O}_X$-가군의 층이다.
\end{enumerate}
\end{lemma}

\begin{proof}
$\mathcal{F}'$를 정의하는 규칙이 $\mathcal{F}$의 부분층을 준다는 것은 명백하다.
따라서 나머지 명제를 확인할 때 $X$ 위에서 에탈 국소적으로 작업해도 된다.
그러므로 보조정리는 스킴의 경우인 Properties, Lemma
\ref{properties-lemma-sections-supported-on-closed-subset}로 환원된다.
\end{proof}

\begin{definition}
\label{spaces-limits-definition-subsheaf-sections-supported-on-closed}
$S$를 스킴이라 하자. $X$를 $S$ 위의 대수공간이라 하자.
$T \subset |X|$를 닫힌 부분집합이라 하자. 그 여집합이 열린 부분공간
$U \subset X$에 대응하고 포함 사상 $U \to X$가 준콤팩트라고 가정하자.
$\mathcal{F}$를 준연접 $\mathcal{O}_X$-가군이라 하자.
위의 Lemma \ref{spaces-limits-lemma-sections-supported-on-closed-subset}에서 정의한
준연접 부분층 $\mathcal{F}' \subset \mathcal{F}$를
{\it $T$ 위에 지지를 갖는 절단들의 부분층}이라 한다.
\end{definition}

\begin{lemma}
\label{spaces-limits-lemma-push-sections-supported-on-closed-subset}
$S$를 스킴이라 하자.
$f : X \to Y$를 $S$ 위의 대수공간의 준콤팩트 준분리 사상이라 하자.
$T \subset |Y|$를 닫힌 부분집합이라 하자.
$|Y| \setminus T$가 열린 부분공간 $V \subset Y$에 대응하고,
$V \to Y$가 준콤팩트라고 가정하자.
$\mathcal{F}$를 준연접 $\mathcal{O}_X$-가군이라 하자.
$\mathcal{F}' \subset \mathcal{F}$를 $|f|^{-1}T$ 위에 지지를 갖는
절단들의 부분층이라 하자. 그러면
$f_*\mathcal{F}' \subset f_*\mathcal{F}$는 $T$ 위에 지지를 갖는
절단들의 부분층이다.
\end{lemma}

\begin{proof}
생략한다. 힌트:
% P09-ERR-0103: frozen source's "support of the open subspace" retained literally.
$|X| \setminus |f|^{-1}T$는 열린 부분공간 $U = f^{-1}V \subset X$의 지지이다.
$V \to Y$가 준콤팩트이므로 $U \to X$도 밑변환에 의해 준콤팩트이다.
$f$가 준콤팩트이고 준분리라는 가정에 의해 $f_*\mathcal{F}$는 준연접이다.
따라서 Lemma \ref{spaces-limits-lemma-sections-supported-on-closed-subset}을
$T$와 $f_*\mathcal{F}$에 적용할 수 있고,
$|f|^{-1}T$와 $\mathcal{F}$에도 적용할 수 있다.
주어진 준연접 가군들이 같다는 것은 정의에서 즉시 따른다.
\end{proof}















\section{아핀 공간의 특성화}
\label{spaces-limits-section-affine}

\noindent
이 절은 Limits, Section \ref{limits-section-affine}에 대응한다.

\begin{lemma}
\label{spaces-limits-lemma-affine}
$S$를 스킴이라 하자. $f : X \to Y$를 $S$ 위의 대수공간의 사상이라 하자.
$f$가 전사 유한 사상이고 $X$가 아핀이라고 가정하자.
그러면 $Y$는 아핀이다.
\end{lemma}

\begin{proof}
$f : X \to Y$를 $\Spec(\mathbf{Z})$ 위의 대수공간의 사상으로 보자
(Spaces, Definition \ref{spaces-definition-base-change}을 보라).
유한 사상은 아핀이고 보편 닫힌 사상임에 주의하자.
Morphisms of Spaces, Lemma
\ref{spaces-morphisms-lemma-integral-universally-closed}을 보라.
Morphisms of Spaces, Lemma
\ref{spaces-morphisms-lemma-image-universally-closed-separated}에 의해
$Y$는 분리 대수공간이다.
$f$가 전사이고 $X$가 준콤팩트이므로 $Y$가 준콤팩트임을 알 수 있다.

\medskip\noindent
Lemma \ref{spaces-limits-lemma-finite-in-finite-and-finite-presentation}에 의해
$X = \lim X_a$로 쓸 수 있으며, 각 $X_a \to Y$는 유한하고 유한 표시이다.
Lemma \ref{spaces-limits-lemma-limit-is-affine}에 의해
$X_a$가 아핀인 충분히 큰 $a$가 존재한다.
따라서 $f : X \to Y$가 유한 표시인 유한 전사라고 가정해도 된다.

\medskip\noindent
Proposition \ref{spaces-limits-proposition-approximate}에 의해
$Y = \lim Y_i$를 $\mathbf{Z}$ 위에서 유한 표시인 대수공간들의
유향 극한으로 쓸 수 있다.
Lemma \ref{spaces-limits-lemma-descend-finite-presentation}에 의해
$0 \in I$와 유한 표시 사상 $X_0 \to Y_0$를 찾아,
$X_i = X_0 \times_{Y_0} Y_i$가 $i \geq 0$에 대해 성립하고
$X = \lim_i X_i$가 되게 할 수 있다.
Lemma \ref{spaces-limits-lemma-descend-finite}에 의해
$X_i \to Y_i$가 유한인 충분히 큰 $i$가 존재한다.
Lemma \ref{spaces-limits-lemma-descend-surjective}에 의해
$X_i \to Y_i$가 전사인 충분히 큰 $i$가 존재한다.
Lemma \ref{spaces-limits-lemma-limit-is-affine}에 의해 $X_i$가 아핀인 충분히 큰 $i$가 존재한다.
따라서 충분히 큰 $i$에 대해 Cohomology of Spaces, Lemma
\ref{spaces-cohomology-lemma-image-affine-finite-morphism-affine-Noetherian}을 적용하여
$Y_i$가 아핀임을 얻는다. 이는 $Y$가 아핀임을 뜻하므로 증명이 끝난다.
\end{proof}

\begin{proposition}
\label{spaces-limits-proposition-affine}
$S$를 스킴이라 하자. $f : X \to Y$를 $S$ 위의 대수공간의 사상이라 하자.
$X$가 아핀이고 $f$가 전사이며 보편 닫힌 사상이라고 가정하자
\footnote{정수적 사상은 보편 닫힌 사상이다. Morphisms of Spaces, Lemma
\ref{spaces-morphisms-lemma-integral-universally-closed}을 보라.}.
그러면 $Y$는 아핀이다.
\end{proposition}

\begin{proof}
$f : X \to Y$를 $\Spec(\mathbf{Z})$ 위의 대수공간의 사상으로 보자
(Spaces, Definition \ref{spaces-definition-base-change}을 보라).
Morphisms of Spaces, Lemma
\ref{spaces-morphisms-lemma-image-universally-closed-separated}에 의해
$Y$는 분리 대수공간이다. 이어서 Morphisms of Spaces, Lemma
\ref{spaces-morphisms-lemma-affine-permanence}에 의해 $f$가 아핀임을 얻는다.
그러면 Morphisms of Spaces, Lemma
\ref{spaces-morphisms-lemma-integral-universally-closed}에 의해
$f$는 정수적이다.

\medskip\noindent
앞 문단에 의해 $f : X \to Y$가 전사 정수적이고,
$X$가 아핀이며 $Y$가 분리라고 가정해도 된다.
$f$가 전사이고 $X$가 준콤팩트이므로 $Y$도 준콤팩트임을 얻는다.

\medskip\noindent
층 $\mathcal{A} = f_*\mathcal{O}_X$를 생각하자.
이는 준연접 $\mathcal{O}_Y$-대수 층이다.
Morphisms of Spaces, Lemma \ref{spaces-morphisms-lemma-pushforward}를 보라.
Lemma \ref{spaces-limits-lemma-colimit-finitely-presented}에 의해
$\mathcal{A} = \colim_i \mathcal{F}_i$를 유한형
$\mathcal{O}_Y$-가군들의 여과 여극한으로 쓸 수 있다.
$\mathcal{A}_i \subset \mathcal{A}$를 $\mathcal{O}_Y$-부분대수 중
$\mathcal{F}_i$가 생성하는 것으로 하자. 대수의 사상
$\mathcal{O}_Y \to \mathcal{A}$가 정수적이므로, 각 $\mathcal{A}_i$는
유한 준연접 $\mathcal{O}_Y$-대수이다. 따라서
$$
X_i = \underline{\Spec}_Y(\mathcal{A}_i) \longrightarrow Y
$$
는 대수공간의 유한 사상이다. 여기서
$\underline{\Spec}$는 Morphisms of Spaces, Lemma
\ref{spaces-morphisms-lemma-affine-equivalence-algebras}의 구성이다.
$X = \lim_i X_i$라는 것은 명백하다. 따라서
Lemma \ref{spaces-limits-lemma-limit-is-affine}에 의해 충분히 큰 $i$에 대해
스킴 $X_i$는 아핀이다. 또한 $X \to Y$가 각 $X_i$를 통해
인수분해되므로 $X_i \to Y$는 전사이다. 따라서
Lemma \ref{spaces-limits-lemma-affine}에 의해 $Y$가 아핀임을 얻는다.
\end{proof}

\noindent
위 결과의 다음 따름정리는 \cite{CLO}에서 찾을 수 있다.

\begin{lemma}
\label{spaces-limits-lemma-reduction-scheme}
\begin{reference}
\cite[3.1.12]{CLO}
\end{reference}
$S$를 스킴이라 하자. $X$를 $S$ 위의 대수공간이라 하자.
$X_{red}$가 스킴이면 $X$도 스킴이다.
\end{lemma}

\begin{proof}
$U' \subset X_{red}$를 아핀 열린 부분스킴이라 하자.
$U \subset X$를 열린집합
$|U'| \subset |X_{red}| = |X|$에 대응하는 열린 부분공간이라 하자.
그러면 $U' \to U$는 전사 정수적이다. 따라서
Proposition \ref{spaces-limits-proposition-affine}에 의해 $U$는 아핀이다.
그러므로 모든 점이 $X$의 열린 부분스킴에 포함된다. 즉,
$X$는 스킴이다.
\end{proof}

\begin{lemma}
\label{spaces-limits-lemma-integral-universally-bijective-scheme}
$S$를 스킴이라 하자. $f : X \to Y$를 $S$ 위의 대수공간의 사상이라 하자.
$f$가 정수적이고 전단사 $|X| \to |Y|$를 유도한다고 가정하자.
그러면 $X$가 스킴일 필요충분조건은 $Y$가 스킴인 것이다.
\end{lemma}

\begin{proof}
정수적 사상은 정의에 의해 표현 가능하므로, $Y$가 스킴이면 $X$도 스킴이다.
반대로 $X$가 스킴이라고 가정하자. $U \subset X$를 아핀 열린집합이라 하자.
정수적 사상은 닫힌 사상이고 $|f|$는 전단사이므로,
$|f|(|U|) \subset |Y|$는 $|f|(|X| \setminus |U|)$의 여집합으로서 열린집합이다.
$V \subset Y$를 $|V| = |f|(|U|)$인 열린 부분공간이라 하자.
Properties of Spaces, Lemma \ref{spaces-properties-lemma-open-subspaces}를 보라.
그러면 $U \to V$는 정수적 전사이므로,
Proposition \ref{spaces-limits-proposition-affine}에 의해 $V$는 아핀 스킴이다.
이로써 증명이 끝난다.
\end{proof}

\begin{lemma}
\label{spaces-limits-lemma-check-closed-infinitesimally}
$S$를 스킴이라 하자.
$f : X \to B$와 $B' \to B$를 $S$ 위의 대수공간의 사상이라 하자.
다음을 가정하자.
\begin{enumerate}
\item $B' \to B$는 닫힌 몰입이다.
\item $|B'| \to |B|$는 전단사이다.
\item $X \times_B B' \to B'$는 닫힌 몰입이다.
\item $X \to B$는 유한형이거나 $B' \to B$는 유한 표시이다.
\end{enumerate}
그러면 $f : X \to B$는 닫힌 몰입이다.
\end{lemma}

\begin{proof}
가정 (1)과 (2)에 의해 $B_{red} = B'_{red}$이다.
$X' = X \times_B B'$로 놓자.
% P09-ERR-0104: frozen source's missing article is immaterial in Korean.
그러면 $X' \to X$는 닫힌 몰입이고 $X'_{red} = X_{red}$이다.
$U \to B$를 에탈 사상이라 하고 $U$가 아핀이라고 하자.
그러면 $X' \times_B U \to X \times_B U$는
기저 축소 공간 위에서 동형을 유도하는 대수공간의 닫힌 몰입이다.
$X' \times_B U$는 스킴이다
($B' \to B$와 $X' \to B'$가 표현 가능하기 때문이다).
따라서 Lemma \ref{spaces-limits-lemma-reduction-scheme}에 의해 $X \times_B U$도 스킴이다.
그러므로 $X \to B$도 표현 가능하다. 이에 따라 스킴의 경우로 환원되며,
Morphisms, Lemma \ref{morphisms-lemma-check-closed-infinitesimally}을 보라.
\end{proof}








\section{스킴에 의한 유한 피복}
\label{spaces-limits-section-finite-cover}

\noindent
이 장의 극한 결과를 응용하여 다음을 증명한다. 임의의 준콤팩트 준분리
대수공간 $X$가 주어지면 스킴 $Y$와 전사 유한 사상 $Y \to X$가 존재한다.
% P09-ERR-0105: frozen source's erroneous finite integral cover claim retained.
이를 위해 스킴에 의한 유한 정수적 피복을 찾을 수 있다는 이미 증명된 결과에 의존한다.
이 결과는 Decent Spaces, Section
\ref{decent-spaces-section-integral-cover}에서 증명되었다.

\begin{proposition}
\label{spaces-limits-proposition-there-is-a-scheme-finite-over}
$S$를 스킴이라 하자. $X$를 $S$ 위의 준콤팩트 준분리 대수공간이라 하자.
\begin{enumerate}
\item 유한 표시인 전사 유한 사상 $Y \to X$가 존재하며, $Y$는 스킴이다.
\item 전사 에탈 사상 $U \to X$가 주어지면 $Y \to X$를 택하여,
모든 $y \in Y$에 대해 열린 근방 $V \subset Y$가 존재하고
$V \to X$가 $U$를 통해 인수분해되게 할 수 있다.
\end{enumerate}
\end{proposition}

\begin{proof}
(1)은 $U = X$인 (2)의 특수한 경우이다.
$Y \to X$를 Decent Spaces, Lemma
\ref{decent-spaces-lemma-there-is-a-scheme-integral-over}에서와 같이 택하자.
$Y = \bigcup V_j$인 유한 아핀 열린 피복을 택하여
$V_j \to X$가 $U$를 통해 인수분해되게 하자.
$Y = \lim Y_i$로 쓸 수 있으며, $Y_i \to X$는 유한하고 유한 표시이다.
Lemma \ref{spaces-limits-lemma-integral-limit-finite-and-finite-presentation}을 보라.
충분히 큰 $i$에 대해 대수공간 $Y_i$는 스킴이다.
Lemma \ref{spaces-limits-lemma-limit-is-scheme}을 보라.
충분히 큰 $i$에 대해 아핀 열린집합 $V_{i, j} \subset Y_i$를 찾아,
$Y$에서의 역상이 $V_j$가 되게 할 수 있다.
Lemma \ref{spaces-limits-lemma-descend-opens}를 보라.
$i$를 더 키우면 사상 $V_j \to U$는 $X$ 위의 사상이고,
사상 $V_{i, j} \to U$에서 오며 이 사상도 $X$ 위에 있다.
Proposition \ref{spaces-limits-proposition-characterize-locally-finite-presentation}을 보라.
이로써 증명이 끝난다.
\end{proof}

\begin{lemma}
\label{spaces-limits-lemma-integral-limit-finite-and-finite-presentation-refined}
$S$를 스킴이라 하자. $f : X \to Y$를 $S$ 위의 대수공간의 정수적 사상이라 하자.
$Y$가 준콤팩트이고 준분리라고 가정하자.
$V \subset Y$를 준콤팩트 열린 부분공간이라 하고,
$f^{-1}(V) \to V$가 유한하고 유한 표시라고 가정하자.
그러면 $X$를 유향 극한 $X = \lim X_i$로 쓸 수 있으며,
$f_i : X_i \to Y$는 유한하고 유한 표시이고,
사상 $f^{-1}(V) \to f_i^{-1}(V)$는 모든 $i$에 대해 동형이다.
\end{lemma}

\begin{proof}
이 보조정리는 Proposition
\ref{spaces-limits-proposition-there-is-a-scheme-finite-over}을 약간 정교화한 것이다.
정수적인 준연접 $\mathcal{O}_Y$-대수
$\mathcal{A} = f_*\mathcal{O}_X$를 생각하자. 다음 문단에서
$\mathcal{A} = \colim \mathcal{A}_i$를 유한하고 유한 표시인
$\mathcal{O}_Y$-대수 $\mathcal{A}_i$들의 유향 여극한으로 쓰되,
$\mathcal{A}_i|_V = \mathcal{A}|_V$가 되게 할 것이다.
이렇게 한 뒤 $X_i = \underline{\Spec}_Y(\mathcal{A}_i)$로 놓는다.
Morphisms of Spaces, Definition
\ref{spaces-morphisms-definition-relative-spec}을 보라.
구성에 의해 $X_i \to Y$는 유한하고 유한 표시이며,
$X = \lim X_i$이고 $f_i^{-1}(V) = f^{-1}(V)$이다.

\medskip\noindent
대수에 대한 이 주장의 증명은 Lemma
\ref{spaces-limits-lemma-integral-algebra-directed-colimit-finite}의 (2)의 증명과 비슷하다.
먼저 Lemma \ref{spaces-limits-lemma-colimit-finitely-presented}을 사용하여
$\mathcal{A} = \colim \mathcal{F}_i$를 유한 표시
$\mathcal{O}_Y$-가군들의 여극한으로 쓰자.
$\mathcal{A}|_V$는 유한형 $\mathcal{O}_V$-가군이므로,
$\mathcal{F}_i|_V \to \mathcal{A}|_V$가 모든 $i$에 대해 전사라고
가정해도 된다. 각 $i$에 대해 $\mathcal{J}_i$를 다음 사상의 핵이라 하자.
$$
\text{Sym}^*_{\mathcal{O}_X}(\mathcal{F}_i) \longrightarrow \mathcal{A}
$$
% P09-ERR-0106: every frozen O_X base in this paragraph is retained literally.
$i' \geq i$이면 유도된 사상 $\mathcal{J}_i \to \mathcal{J}_{i'}$가 존재한다.
다음이 성립한다.
$$
\mathcal{A} =
\colim \text{Sym}^*_{\mathcal{O}_X}(\mathcal{F}_i)/\mathcal{J}_i
$$
또한 준연접 $\mathcal{O}_X$-대수
$\text{Sym}^*_{\mathcal{O}_X}(\mathcal{F}_i)/\mathcal{J}_i$는 유한이다
(이는 정수적 준연접 $\mathcal{O}_Y$-대수 $\mathcal{A}$의
유한형 준연접 부분대수이며, 그 대수는 $\mathcal{O}_X$ 위에 있기 때문이다).
위의 전사성에 의해
$\text{Sym}^*_{\mathcal{O}_X}(\mathcal{F}_i)/\mathcal{J}_i$를 $V$에
제한하면 $\mathcal{A}|_V$가 된다. 따라서 $\mathcal{J}_i|_V$는
$\text{Sym}^*_{\mathcal{O}_X}(\mathcal{F}_i)|_V$의 아이디얼 층으로서
유한 생성이다. 이는 $\mathcal{A}|_V$가 유한 표시
$\mathcal{O}_Y$-대수라는 사실에 의한다. $\mathcal{J}_i = \colim \mathcal{E}_{ik}$를
유한 표시 $\mathcal{O}_X$-가군들의 여극한으로 쓰자.
% P09-ERR-0107: frozen source grammar is normalized in Korean only.
위의 유한 생성성에 의해 $\mathcal{E}_{ik}|_V$가
$\mathcal{J}_i|_V$를 생성한다고 가정해도 된다. 여기서 생성은
$\text{Sym}^*_{\mathcal{O}_X}(\mathcal{F}_i)|_V$의 아이디얼 층으로서이다.
$i' \geq i$와 $k$가 주어지면 어떤 $k'$가 존재하여,
다음 가환도가 가환하게 하는 사상
$\mathcal{E}_{ik} \to \mathcal{E}_{i'k'}$가 존재한다.
$$
\xymatrix{
\mathcal{J}_i \ar[r] & \mathcal{J}_{i'} \\
\mathcal{E}_{ik} \ar[u] \ar[r] & \mathcal{E}_{i'k'} \ar[u]
}
$$
이는 Cohomology of Spaces, Lemma
\ref{spaces-cohomology-lemma-finite-presentation-quasi-compact-colimit}에서 따른다.
이로부터 다음 사상을 얻는다.
$$
\mathcal{A}_{ik} =
\text{Sym}^*_{\mathcal{O}_X}(\mathcal{F}_i)/(\mathcal{E}_{ik})
\longrightarrow
\text{Sym}^*_{\mathcal{O}_X}(\mathcal{F}_{i'})/(\mathcal{E}_{i'k'}) =
\mathcal{A}_{i'k'}
$$
여기서 $(\mathcal{E}_{ik})$는 $\mathcal{E}_{ik}$가 생성하는 아이디얼을 뜻한다.
% P09-ERR-0108: frozen source's reversed A_{ki} index is retained.
준연접 $\mathcal{O}_X$-대수 $\mathcal{A}_{ki}$는 충분히 큰 $k$에 대해
유한 표시이고 유한이다
(Lemma \ref{spaces-limits-lemma-finite-algebra-directed-colimit-finite-finitely-presented}의
증명을 보라). 또한 구성에 의해
$\mathcal{A}_{ik}|_V = \mathcal{A}|_V$이다. 마지막으로 다음이 성립한다.
$$
\colim \mathcal{A}_{ik} = \colim \mathcal{A}_i = \mathcal{A}
$$
실제로 첫 번째 등식은 Lemma
\ref{spaces-limits-lemma-finite-algebra-directed-colimit-finite-finitely-presented}의 증명에서 보였고,
두 번째 등식은 $\mathcal{A}$가 가군 $\mathcal{F}_i$들의 여극한이기 때문이다.
\end{proof}

\begin{lemma}
\label{spaces-limits-lemma-there-is-a-scheme-finite-over-refined}
$S$를 스킴이라 하자. $X$를 $S$ 위의 준콤팩트 준분리 대수공간이라 하고,
$|X|$가 유한 개의 기약 성분을 가진다고 가정하자.
\begin{enumerate}
\item 유한 표시인 전사 유한 사상 $f : Y \to X$가 존재하며,
$Y$는 스킴이고, $f$는 준콤팩트 조밀한 열린집합 $U \subset X$ 위에서
유한 에탈 사상이다.
\item 전사 에탈 사상 $V \to X$가 주어지면 $Y \to X$를 택하여,
모든 $y \in Y$에 대해 열린 근방 $W \subset Y$가 존재하고
$W \to X$가 $V$를 통해 인수분해되게 할 수 있다.
\end{enumerate}
\end{lemma}

\begin{proof}
(1)은 $V = X$인 (2)의 특수한 경우이다.

\medskip\noindent
(2)를 증명하자. $\pi : Y \to X$를 Decent Spaces, Lemma
\ref{decent-spaces-lemma-there-is-a-scheme-integral-over-refined}에서와 같이 택하고,
$U \subset X$를 준콤팩트 조밀한 열린집합으로 택하여
$\pi^{-1}(U) \to U$가 유한 에탈이게 하자.
$Y = \bigcup W_j$인 유한 아핀 열린 피복을 택하여
$W_j \to X$가 $V$를 통해 인수분해되게 하자.
$Y = \lim Y_i$로 쓸 수 있으며, $\pi_i : Y_i \to X$는 유한하고
유한 표시이고, $\pi^{-1}(U) \to \pi_i^{-1}(U)$는 동형이다.
Lemma \ref{spaces-limits-lemma-integral-limit-finite-and-finite-presentation-refined}를 보라.
충분히 큰 $i$에 대해 대수공간 $Y_i$는 스킴이다.
Lemma \ref{spaces-limits-lemma-limit-is-scheme}을 보라.
충분히 큰 $i$에 대해 아핀 열린집합 $W_{i, j} \subset Y_i$를 찾아
$Y$에서의 역상이 $W_j$가 되게 할 수 있다.
Lemma \ref{spaces-limits-lemma-descend-opens}를 보라.
$i$를 더 키우면 사상 $W_j \to V$는 $X$ 위의 사상이고,
% P09-ERR-0109: frozen source's U target is retained rather than silently changed to V.
사상 $W_{i, j} \to U$에서 오며 이 사상도 $X$ 위에 있다.
Proposition \ref{spaces-limits-proposition-characterize-locally-finite-presentation}을 보라.
이로써 증명이 끝난다.
\end{proof}

\begin{lemma}
\label{spaces-limits-lemma-there-is-a-scheme-finite-over-filtered}
$S$를 스킴이라 하자. $X$를 $S$ 위의 준콤팩트 준분리 대수공간이라 하자.
어떤 $t \geq 0$와 닫힌 부분공간들의 열
$$
X \supset Z_0 \supset Z_1 \supset \ldots \supset Z_t = \emptyset
$$
이 존재하여, $Z_i \to X$는 유한 표시이고,
$Z_0 \subset X$는 두껍게 하기이며, 각 $i = 0, \ldots t - 1$에 대해
스킴 $Y_i$와 유한 표시인 전사 유한 사상 $Y_i \to Z_i$가 존재하고,
이 사상은 $Z_i \setminus Z_{i + 1}$ 위에서 유한 에탈이다.
\end{lemma}

\begin{proof}
$X$를 $\Spec(\mathbf{Z})$ 위의 대수공간으로 볼 수 있다.
Spaces, Definition \ref{spaces-definition-base-change}과
Properties of Spaces, Definition \ref{spaces-properties-definition-separated}를 보라.
따라서 Proposition \ref{spaces-limits-proposition-approximate}을 적용할 수 있다.
이에 따라 아핀 사상 $X \to X_0$을 찾을 수 있으며,
$X_0$는 $\mathbf{Z}$ 위에서 유한 표시이다.
$X_0$에 대해 보조정리를 증명할 수 있다면, 층화와 사상들을 $X$로 당겨
$X$에 대한 결과를 얻을 수 있다. 몇 가지 세부 사항은 생략한다.
이로써 다음 문단에서 다룰 경우로 환원된다.

\medskip\noindent
$X$가 $\mathbf{Z}$ 위에서 유한 표시라고 가정하자.
그러면 $X$는 Noether이고 $|X|$는 유한 개의 기약 성분을 가지는
유한 차원 Noether 위상공간이다.
따라서 $\dim(|X|)$에 관한 귀납법을 사용할 수 있다.
$X$로 가는 모든 유한 사상은 유한 표시이므로,
증명의 나머지 부분에서는 그 요구를 무시할 수 있다.
Lemma \ref{spaces-limits-lemma-there-is-a-scheme-finite-over-refined}에 의해
전사 유한 사상 $Y \to X$가 존재하며, 이는 조밀한 열린집합
$U \subset X$ 위에서 유한 에탈이다.
$Z_0 = X$로 놓고, $Z_1 \subset X$를
$|Z_1| = |X| \setminus |U|$인 축소 닫힌 부분공간이라 하자.
귀납법에 의해 정수 $t \geq 0$와 닫힌 부분공간들에 의한 여과
$$
Z_1 \supset Z_{1, 0} \supset Z_{1, 1} \supset \ldots
\supset Z_{1, t} = \emptyset
$$
를 얻는다. 여기서 $Z_{1, 0} \to Z_1$은 두껍게 하기이고,
유한 전사 사상 $Y_{1, i} \to Z_{1, i}$가 존재하여
$Z_{1, i} \setminus Z_{1, i + 1}$ 위에서 유한 에탈이다.
$Z_1$은 축소되어 있으므로 $Z_1 = Z_{1, 0}$이다.
따라서 $Z_i = Z_{1, i - 1}$와 $Y_i = Y_{1, i - 1}$로 놓을 수 있다.
이는 $i \geq 1$에 대한 것이며, 보조정리가 증명된다.
\end{proof}











\section{스킴을 얻는 방법}
\label{spaces-limits-section-representable}

\noindent
대수공간이 스킴임을 보이는 몇 가지 기법을 더 제시한다.
첫 번째 기법은 스킴이 아닌 최소 닫힌 부분공간이 존재함을 보이는 것이다.

\begin{lemma}
\label{spaces-limits-lemma-minimal-closed-subspace}
$S$를 스킴이라 하자. $X$를 $S$ 위의 준콤팩트 준분리 대수공간이라 하자.
$X$가 스킴이 아니면 닫힌 부분공간 $Z \subset X$가 존재하고,
$Z$는 스킴이 아니지만 모든 진 닫힌 부분공간 $Z' \subset Z$는 스킴이다.
\end{lemma}

\begin{proof}
Zorn 보조정리로 증명한다. $\mathcal{Z}$를 스킴이 아닌 닫힌 부분공간
$Z$들의 집합이라 하고 포함관계로 순서화하자.
가정에 의해 $\mathcal{Z}$는 $X$를 포함하므로 공집합이 아니다.
$Z_\alpha$가 $\mathcal{Z}$의 전순서 부분집합이면
$Z = \bigcap Z_\alpha$도 $\mathcal{Z}$에 속한다. 실제로
$$
Z = \lim Z_\alpha
$$
이고 전이 사상들은 아핀이다.
따라서 Lemma \ref{spaces-limits-lemma-limit-is-scheme}을 적용하면,
$Z$가 스킴인 경우 어떤 $Z_\alpha$도 스킴이 됨을 알 수 있다.
($Z = \emptyset$인 경우에도 이 논증은 성립하지만,
Lemma \ref{spaces-limits-lemma-limit-nonempty}에 의해 이 경우는 일어날 수 없다.)
따라서 Zorn 보조정리에 의해 $\mathcal{Z}$는 최소 원소를 가진다.
\end{proof}

\noindent
이제 이러한 최소 비스킴에 관해 조금 더 증명할 수 있다.

\begin{lemma}
\label{spaces-limits-lemma-minimal-nonscheme}
$S$를 스킴이라 하자. $X$를 $S$ 위의 준콤팩트 준분리 대수공간이라 하자.
모든 진 닫힌 부분공간 $Z \subset X$가 스킴이지만 $X$는 스킴이 아니라고
가정하자. 그러면 $X$는 축소되어 있고 기약이다.
\end{lemma}

\begin{proof}
Lemma \ref{spaces-limits-lemma-reduction-scheme}에 의해 $X$는 축소되어 있다.
닫힌 부분집합 $T_1 \subset |X|$와 $T_2 \subset |X|$를 택하여
$|X| = T_1 \cup T_2$라 하자. $T_1$과 $T_2$가 진 닫힌 부분집합이면,
이에 대응하는 축소 유도 닫힌 부분공간 $Z_1, Z_2 \subset X$는
(Properties of Spaces, Definition
\ref{spaces-properties-definition-reduced-induced-space})
스킴이다. 또한 $Z = Z_1 \times_X Z_2 = Z_1 \cap Z_2$도
$Z_1$이나 $Z_2$의 닫힌 부분스킴으로서 스킴이다.
여곱 $Z_1 \amalg_Z Z_2$는 스킴의 범주에서 존재함에 주의하자.
More on Morphisms, Lemma
\ref{more-morphisms-lemma-pushout-along-closed-immersions}을 보라.
% P09-ERR-0110: frozen source's comma before “is” is immaterial in Korean.
한 가지 방법은 $Z_1 \amalg_Z Z_2$가 $X$와 동형임을 보이는 것이지만,
대수공간의 밀어내기에 관한 내용은 이 이론의 뒤에서 나오므로 여기서는
이 방법을 사용할 수 없다. 대신 Lemma \ref{spaces-limits-lemma-affine}을 사용하여
각 점의 아핀 근방을 찾는다. 구체적으로 $x \in |X|$라 하자.
$x \not \in Z_1$이면 $x$는 스킴인 근방 $X \setminus Z_1$을 가진다.
$x \not \in Z_2$인 경우도 마찬가지이다.
$x \in Z = Z_1 \cap Z_2$이면 아핀 열린집합
$U \subset Z_1 \amalg_Z Z_2$를 택하자. 이 열린집합은 $z$를 포함한다.
% P09-ERR-0111: frozen source's undefined point z is retained literally.
그러면 $U_1 = Z_1 \cap U$와 $U_2 = Z_2 \cap U$는 아핀 열린집합이고,
$Z$와의 교차가 서로 같다. $|Z_1| = T_1$과 $|Z_2| = T_2$는
$|X|$의 닫힌 부분집합이며 $|Z|$에서 교차하므로,
열린집합 $W \subset |X|$를 찾아
$W \cap T_1 = |U_1|$와 $W \cap T_2 = |U_2|$가 되게 할 수 있다.
$W$를 이에 대응하는 $X$의 열린 부분공간이라고도 쓰자.
그러면 $x \in |W|$이고 사상 $U_1 \amalg U_2 \to W$는 전사 유한 사상이며,
그 정의역은 아핀 스킴이다. 따라서 Lemma \ref{spaces-limits-lemma-affine}에 의해
$W$는 아핀 스킴이다.
\end{proof}

\noindent
다음 보조정리의 핵심은 $X$의 점들의 상에서만 조건을 확인하면 된다는 것이다.

\begin{lemma}
\label{spaces-limits-lemma-enough-local}
$f: X \to S$를 대수공간에서 스킴 $S$로 가는 준콤팩트 준분리 사상이라 하자.
모든 $x \in |X|$에 대해 그 상을 $s = f(x) \in S$라 할 때
대수공간 $X \times_S \Spec(\mathcal{O}_{S,s})$가 스킴이라고 가정하자.
그러면 $X$는 스킴이다.
\end{lemma}

\begin{proof}
$x \in |X|$라 하자. 열린 근방 $U$를 찾아, $s = f(x)$의 근방이 되며
$X \times_S U$가 스킴이게 하면 충분하다.
$X \times_S \Spec(\mathcal{O}_{S, s})$는 스킴이고,
$\mathcal{O}_{S, s} = \colim \mathcal{O}_S(U)$이다. 여기서 여극한은
$s$의 아핀 열린 근방들, 즉 $S$ 안의 그러한 근방들을 따라 취한다. 따라서
$$
X \times_S \Spec(\mathcal{O}_{S, s}) = \lim X \times_S U
$$
이다. Lemma \ref{spaces-limits-lemma-limit-is-scheme}에 의해
% P09-ERR-0112: frozen source's malformed paired connective is normalized only by Korean syntax.
$X \times_S U$가 스킴인 어떤 $U$가 존재함을 알 수 있다.
\end{proof}

\noindent
Lemma \ref{spaces-limits-lemma-enough-local}에서처럼 국소환으로 제한하는 대신,
밑의 닫힌 부분스킴으로 제한할 수 있다.

\begin{lemma}
\label{spaces-limits-lemma-maximal-ideal}
$\varphi : X \to \Spec(A)$를 대수공간에서 아핀 스킴으로 가는
준콤팩트 준분리 사상이라 하자. $X$가 스킴이 아니면
아이디얼 $I \subset A$가 존재하여 밑변환 $X_{A/I}$는 스킴이 아니지만,
모든 $I \subset I'$ 가운데 $I \not = I'$인 것에 대해
밑변환 $X_{A/I'}$는 스킴이다.
\end{lemma}

\begin{proof}
Zorn 보조정리로 증명한다. $\mathcal{I}$를 아이디얼 $I$들의 집합이라 하자.
이들은 $X_{A/I}$가 스킴이 아닌 아이디얼들이다.
가정에 의해 $\mathcal{I}$는 $(0)$을 포함한다.
$I_\alpha$가 $\mathcal{I}$에 속하는 아이디얼들의 사슬이면
$I = \bigcup I_\alpha$도 $\mathcal{I}$에 속한다. 실제로
$A/I = \colim A/I_\alpha$이므로
$$
X_{A/I} = \lim X_{A/I_\alpha}
$$
이다. 따라서 Lemma \ref{spaces-limits-lemma-limit-is-scheme}을 적용하면,
% P09-ERR-0113: frozen source's auxiliary inversion is normalized only by Korean syntax.
$X_{A/I}$가 스킴인 경우 어떤 $X_{A/I_\alpha}$도 스킴이 됨을 알 수 있다.
그러므로 Zorn 보조정리에 의해 $\mathcal{I}$는 극대 원소를 가진다.
\end{proof}






\section{닫힌 올에서의 붙이기}
\label{spaces-limits-section-change-over-closed-points}

\noindent
위의 이론을 국소환의 스펙트럼에 적용하면 상대 대수공간에 관한
몇 가지 유용한 붙이기 결과를 얻는다. 먼저 보조정리를 하나 증명한다.
이 결과는 나중에 Bootstrap, Section
\ref{bootstrap-section-applications}에서 크게 일반화될 것이다.

\begin{lemma}
\label{spaces-limits-lemma-relative-glueing}
$S = U \cup W$를 스킴의 열린 피복이라 하자. 그러면 밑변환으로 주어지는 함자
$$
FP_S \longrightarrow FP_U \times_{FP_{U \cap W}} FP_W
$$
는 범주의 동치이다. 여기서 $FP_T$는 스킴 $T$ 위에서 유한 표시인
대수공간들의 범주이다.
\end{lemma}

\begin{proof}
먼저 $S = U \cup W$가 Zariski 피복이므로,
$(\Sch/S)_{fppf}$ 위의 층들의 범주는 삼중항
$(\mathcal{F}_U, \mathcal{F}_W, \varphi)$들의 범주와 동치이다.
여기서 $\mathcal{F}_U$는 $(\Sch/U)_{fppf}$ 위의 층이고,
$\mathcal{F}_W$는 $(\Sch/W)_{fppf}$ 위의 층이며,
$$
\varphi :
\mathcal{F}_U|_{(\Sch/U \cap W)_{fppf}}
\longrightarrow
\mathcal{F}_W|_{(\Sch/U \cap W)_{fppf}}
$$
는 동형이다. Sites, Lemma \ref{sites-lemma-mapping-property-glue}를 보라.
(다른 붙이기 자료는 필요하지 않는다. 실제로
$U \times_S U = U$, $W \times_S W = W$이고,
같은 이유로 코사이클 조건은 자동으로 성립한다.)
이제 층 $\mathcal{F}$가 $(\Sch/S)_{fppf}$ 위에 놓이고, 이 동치를 통해
$(\mathcal{F}_U, \mathcal{F}_W, \varphi)$로 간다고 하자.
그러면 $\mathcal{F}$가 대수공간일 필요충분조건은
$\mathcal{F}_U$와 $\mathcal{F}_W$가 대수공간인 것이다.
이는 Algebraic Spaces, Lemma
\ref{spaces-lemma-glueing-algebraic-spaces}에서 즉시 따른다.
실제로 $\mathcal{F}_U \to \mathcal{F}$와
$\mathcal{F}_W \to \mathcal{F}$는 열린 몰입으로 표현 가능하며
$\mathcal{F}$를 덮는다. 마지막으로 이 경우 대수공간 $\mathcal{F}$가
$S$ 위에서 유한 표시일 필요충분조건은
$\mathcal{F}_U$가 $U$ 위에서 유한 표시이고
$\mathcal{F}_W$가 $W$ 위에서 유한 표시인 것이다.
이는 Morphisms of Spaces, Lemmas
\ref{spaces-morphisms-lemma-quasi-compact-local},
\ref{spaces-morphisms-lemma-separated-local},
\ref{spaces-morphisms-lemma-finite-presentation-local}에서 따른다.
\end{proof}

\begin{lemma}
\label{spaces-limits-lemma-glueing-near-closed-point}
$S$를 스킴이라 하자. $s \in S$를 닫힌점이라 하고,
$U = S \setminus \{s\} \to S$가 준콤팩트라고 가정하자.
$V = \Spec(\mathcal{O}_{S, s}) \setminus \{s\}$로 놓으면 범주의 동치
$$
FP_S \longrightarrow FP_U \times_{FP_V} FP_{\Spec(\mathcal{O}_{S, s})}
$$
가 존재한다. 여기서 $FP_T$는 $T$ 위에서 유한 표시인 대수공간들의 범주이다.
\end{lemma}

\begin{proof}
$W \subset S$를 $s$의 열린 근방이라 하자. 함자
$$
FP_S \to FP_U \times_{FP_{W \setminus \{s\}}} FP_W
$$
는 Lemma \ref{spaces-limits-lemma-relative-glueing}에 의해 범주의 동치이다.
$\mathcal{O}_{S, s} = \colim \mathcal{O}_W(W)$이고, 여기서 $W$는
$s$의 아핀 열린 근방들을 움직인다. 따라서
$\Spec(\mathcal{O}_{S, s}) = \lim W$이고, 여기서도 $W$는
$s$의 아핀 열린 근방들을 움직인다.
% P09-ERR-0114: frozen source's singular varying-category wording is rendered grammatically in Korean.
그러므로 $\Spec(\mathcal{O}_{S, s})$ 위에서 유한 표시인 대수공간들의 범주는
$W$ 위에서 유한 표시인 대수공간들의 범주들의 극한이다.
여기서 $W$는 $s$의 아핀 열린 근방들을 움직인다.
Lemma \ref{spaces-limits-lemma-descend-finite-presentation}을 보라.
% P09-ERR-0115: frozen source's reversed affine-open membership relation is retained.
각 아핀 열린집합 $s \in W$에 대해 $U \cap W$는 준콤팩트이다.
실제로 $U \to S$가 준콤팩트이기 때문이다.
따라서 $V = \lim W \cap U = \lim W \setminus \{s\}$는
준콤팩트 준분리 스킴들의 극한이다
(Limits, Lemma \ref{limits-lemma-directed-inverse-system-has-limit}을 보라).
따라서 $V$ 위에서 유한 표시인 대수공간들의 범주도
$W \cap U$ 위에서 유한 표시인 대수공간들의 범주들의 극한이다.
여기서 $W$는 $s$의 아핀 열린 근방들을 움직인다.
이 결과들을 결합하면 보조정리가 형식적으로 따른다.
\end{proof}

\begin{lemma}
\label{spaces-limits-lemma-glueing-near-point}
$S$를 스킴이라 하자. $U \subset S$를 retrocompact 열린집합이라 하자.
$s \in S$를 $U$의 여집합에 속하는 점이라 하자.
$V = \Spec(\mathcal{O}_{S, s}) \cap U$로 놓으면 범주의 동치
$$
\colim_{s \in U' \supset U\text{ open}} FP_{U'}
\longrightarrow
FP_U \times_{FP_V} FP_{\Spec(\mathcal{O}_{S, s})}
$$
가 존재한다. 여기서 $FP_T$는 $T$ 위에서 유한 표시인 대수공간들의 범주이다.
\end{lemma}

\begin{proof}
$W \subset S$를 $s$의 열린 근방이라 하자.
Lemma \ref{spaces-limits-lemma-relative-glueing}에 의해 함자
$$
FP_{U \cup W}
\longrightarrow
FP_U \times_{FP_{U \cap W}} FP_W
$$
는 범주의 동치이다. 다음이 성립한다.
$\mathcal{O}_{S, s} = \colim \mathcal{O}_W(W)$.
여기서 $W$는 $s$의 아핀 열린 근방들을 움직인다.
따라서 $\Spec(\mathcal{O}_{S, s}) = \lim W$이고, 여기서 $W$는
$s$의 아핀 열린 근방들을 움직인다.
그러므로 $\Spec(\mathcal{O}_{S, s})$ 위에서 유한 표시인 대수공간들의 범주는
$W$ 위에서 유한 표시인 대수공간들의 범주들의 극한이다.
여기서 $W$는 $s$의 아핀 열린 근방들을 움직인다.
Lemma \ref{spaces-limits-lemma-descend-finite-presentation}을 보라.
각 아핀 열린집합 $s \in W$에 대해 $U \cap W$는 준콤팩트이다.
실제로 $U \to S$가 준콤팩트이기 때문이다.
따라서 $V = \lim W \cap U$는 준콤팩트 준분리 스킴들의 극한이다
(Limits, Lemma \ref{limits-lemma-directed-inverse-system-has-limit}을 보라).
그러므로 $V$ 위에서 유한 표시인 대수공간들의 범주도
$W \cap U$ 위에서 유한 표시인 대수공간들의 범주들의 극한이다.
여기서 $W$는 $s$의 아핀 열린 근방들을 움직인다.
이 결과들을 결합하면 보조정리가 형식적으로 따른다.
\end{proof}

\begin{lemma}
\label{spaces-limits-lemma-glueing-near-multiple-closed-points}
$S$를 스킴이라 하자. $s_1, \ldots, s_n \in S$를 서로 다른 닫힌점들이라 하고,
$U = S \setminus \{s_1, \ldots, s_n\} \to S$가 준콤팩트라고 가정하자.
$S_i = \Spec(\mathcal{O}_{S, s_i})$ 및 $U_i = S_i \setminus \{s_i\}$로 놓으면
범주의 동치
$$
FP_S \longrightarrow
FP_U \times_{(FP_{U_1} \times \ldots \times FP_{U_n})}
(FP_{S_1} \times \ldots \times FP_{S_n})
$$
가 존재한다. 여기서 $FP_T$는 $T$ 위에서 유한 표시인 대수공간들의 범주이다.
\end{lemma}

\begin{proof}
$n = 1$이면 이는 Lemma \ref{spaces-limits-lemma-glueing-near-closed-point}이다.
$n > 1$이면 보조정리를 정확히 같은 방법으로 증명하거나 그 결과에서 유도할 수 있다.
예를 들어 $f_i : X_i \to S_i$가 $FP_{S_i}$의 대상들이고,
$f : X \to U$가 $FP_U$의 대상이며, 동형
$X_i \times_{S_i} U_i = X \times_U U_i$들이 주어졌다고 하자.
Lemma \ref{spaces-limits-lemma-glueing-near-closed-point}에 의해
유한 표시 사상 $f' : X' \to U' = S \setminus \{s_1, \ldots, s_{n - 1}\}$를
찾을 수 있다. 이 사상은 $X_i$와 동형이며 그 동형은 $S_i$ 위에 있고,
% P09-ERR-0116: frozen source's X_i/S_i and X_i over S_n indices are retained literally.
$X$와도 동형이며 그 동형은 $U$ 위에 있다. 이 동형들은 주어진 동형
$X_i \times_{S_n} U_n = X \times_U U_n$과 양립한다.
이제 $f_i : X_i \to S_i$, $i \leq n - 1$과
$f' : X' \to U'$, 그리고 유도된 동형
$X_i \times_{S_i} U_i = X' \times_{U'} U_i$, $i \leq n - 1$에
귀납법을 적용할 수 있다. 이는 본질적 전사성을 보인다.
% P09-ERR-0117: frozen source's malformed noun phrase is normalized only by Korean syntax.
충만 충실성의 증명은 생략한다.
\end{proof}

\section{수정에의 응용}
\label{spaces-limits-section-modifications-at-a-point}

\noindent
극한을 사용하면 닫힌점 위에서 decent 대수공간의 수정들의 범주를
헨젤 국소환을 이용하여 기술할 수 있다.

\begin{lemma}
\label{spaces-limits-lemma-excision-modifications}
$S$를 스킴이라 하자. 대수공간들의 분리 에탈 사상
$f : V \to W$를 생각하자. 이 사상은 $S$ 위에 있다.
닫힌 부분공간 $T \subset W$가 존재하여
$f^{-1}T \to T$가 동형이라고 가정하자. 그러면
$W^0 = W \setminus T$ 및 $V^0 = f^{-1}W^0$로 놓을 때 밑변환 함자
$$
\left\{
\begin{matrix}
g : X \to W\text{는 대수공간들의 사상} \\
g^{-1}(W^0) \to W^0\text{는 동형}
\end{matrix}
\right\}
\longrightarrow
\left\{
\begin{matrix}
h : Y \to V\text{는 대수공간들의 사상} \\
h^{-1}(V^0) \to V^0\text{는 동형}
\end{matrix}
\right\}
$$
는 범주의 동치이다.
\end{lemma}

\begin{proof}
$V \to W$가 분리이므로
$V \times_W V = \Delta(V) \amalg U$이고, 여기서 $U$는
$V \times_W V$의 열린닫힌 부분공간이다.
$f^{-1}T \to T$가 동형이라는 가정에 의해
$U \times_W T = \emptyset$이다. 즉, 두 사영 $U \to V$는 $V^0$로 간다.
% P09-ERR-0118: frozen source's subject-verb disagreement is normalized only by Korean syntax.

\medskip\noindent
오른쪽 범주의 $h : Y \to V$가 주어졌다고 하자.
반변 함자 $X$를 $(\Sch/S)_{fppf}$ 위에서 다음 규칙으로 정의한다.
$$
X(T) = \{(w, y) \mid
w :  T \to W,\ y : T \times_{w, W} V \to Y\text{는 사상이며 그 밑은 }V\}
$$
% P09-ERR-0119: frozen source's missing “by” is immaterial in Korean syntax.
$g : X \to W$를 $(w, y) \in X(T)$를 $w \in W(T)$로 보내는 사상이라 하자.
$h^{-1}V^0 \to V^0$가 동형이므로 $w : T \to W$가 $W^0$로 가면
$h$의 선택은 유일하다.
% P09-ERR-0120: frozen source's choice of h, rather than y, is retained literally.
다시 말해 $X \times_{g, W} W^0 = W^0$이다. 한편
$T$-값 점 $(w, y, v)$를 $X \times_{g, W, f} V$에서 생각하자.
그러면 $w = f \circ v$이고
$$
y : T \times_{f \circ v, W} V \longrightarrow V
$$
% P09-ERR-0121: frozen source's displayed codomain V is retained rather than silently changed to Y.
는 $V$ 위의 사상이다. 사상
$$
T \times_{f \circ v, W} V \xrightarrow{(v, \text{id}_V)}
V \times_W V = V \amalg U
$$
을 생각하자. $V$의 역상은
$T$이고, 이는 $(\text{id}_T, v) : T \to T \times_{f \circ v, W} V$로 매장된다.
합성 $y' = y \circ (\text{id}_T, v) : T \to Y$는
$v = h \circ y'$를 만족하는 사상이며, 이것이 $y$를 결정한다.
실제로 다른 부분에 대한 $y$의 제한은 두 번째 사영에 의해
$U$가 $V^0$로 가므로 유일하게 결정된다. 따라서
$X \times_{g, W, f} V \to Y$, $(w, y, v) \mapsto y'$는 동형이다.

\medskip\noindent
따라서 $X$가 대수공간임을 보이면 증명이 끝난다.
$V \to W$는 분리 에탈 사상이므로 Morphisms of Spaces, Lemma
\ref{spaces-morphisms-lemma-locally-quasi-finite-separated-representable}
(및 Morphisms of Spaces, Lemma
\ref{spaces-morphisms-lemma-etale-locally-quasi-finite})에 의해 표현 가능하다.
물론 $W^0 \to W$는 열린 몰입이므로 표현 가능하고 에탈이다. 따라서
$$
W^0 \amalg Y = X \times_{g, W} W^0 \amalg X \times_{g, W, f} V
= X \times_{g, W} (W^0 \amalg V) \longrightarrow X
$$
는 Spaces, Lemmas
\ref{spaces-lemma-base-change-representable-transformations}와
\ref{spaces-lemma-base-change-representable-transformations-property}에 의해
표현 가능하고 전사이며 에탈이다. 그러므로 Spaces, Lemma
\ref{spaces-lemma-etale-locally-representable-by-space-gives-space}에 의해
$X$는 대수공간이다.
\end{proof}

\begin{lemma}
\label{spaces-limits-lemma-excision-modifications-properties}
Lemma \ref{spaces-limits-lemma-excision-modifications}의 표기와 가정을 사용하자.
$g : X \to W$가 그 동치를 통해 $h : Y \to V$에 대응한다고 하자.
그러면 $g$가 준콤팩트, 준분리, 분리, 국소 유한 표시, 유한 표시,
국소 유한형, 유한형, 고유, 정수적, 유한, 그리고 여기에 더 추가할 것인
필요충분조건은 $h$도 그러한 것이다.
% P09-ERR-0122: frozen source's editorial placeholder “and add more here” is retained literally.
\end{lemma}

\begin{proof}
$g$가 준콤팩트, 준분리, 분리, 국소 유한 표시, 유한 표시,
국소 유한형, 유한형, 고유, 유한이면, $h$도 $g$의 밑변환으로서 그러하다.
% P09-ERR-0123: frozen source's omission of integral from this forward list is retained.
이는 Morphisms of Spaces, Lemmas
\ref{spaces-morphisms-lemma-base-change-quasi-compact},
\ref{spaces-morphisms-lemma-base-change-separated},
\ref{spaces-morphisms-lemma-base-change-finite-presentation},
\ref{spaces-morphisms-lemma-base-change-finite-type},
\ref{spaces-morphisms-lemma-base-change-proper},
\ref{spaces-morphisms-lemma-base-change-integral}에서 따른다.
역으로, 대수공간의 사상들의 성질 $P$가 밑에서 에탈 국소적이고
모든 대수공간의 항등사상에 대해 성립한다고 하자.
$\{W^0 \to W, V \to W\}$가 에탈 피복이므로 $g$가 $P$를 가짐을 보이려면
$h$가 $P$를 가짐을 보이는 것으로 충분하다. 따라서 Morphisms of Spaces, Lemmas
\ref{spaces-morphisms-lemma-quasi-compact-local},
\ref{spaces-morphisms-lemma-separated-local},
\ref{spaces-morphisms-lemma-finite-presentation-local},
\ref{spaces-morphisms-lemma-finite-type-local},
\ref{spaces-morphisms-lemma-proper-local},
\ref{spaces-morphisms-lemma-integral-local}을 사용하여 결론을 얻는다.
\end{proof}

\begin{lemma}
\label{spaces-limits-lemma-modifications}
$S$를 스킴이라 하자. $X$를 $S$ 위의 decent 대수공간이라 하자.
$x \in |X|$를 닫힌점이라 하고 $U = X \setminus \{x\} \to X$가
준콤팩트라고 가정하자.
$V = \Spec(\mathcal{O}_{X, x}^h) \setminus \{\mathfrak m_x^h\}$로 놓으면
밑변환 함자
$$
\left\{
\begin{matrix}
f : Y \to X\text{는 유한 표시} \\
f^{-1}(U) \to U\text{는 동형}
\end{matrix}
\right\}
\longrightarrow
\left\{
\begin{matrix}
g : Y \to \Spec(\mathcal{O}_{X, x}^h)\text{는 유한 표시} \\
g^{-1}(V) \to V\text{는 동형}
\end{matrix}
\right\}
$$
는 범주의 동치이다.
\end{lemma}

\begin{proof}
기본 에탈 근방 $a : (W, w) \to (X, x)$를 $x$에 대해 택하되,
$W$가 아핀이 되게 Decent Spaces, Lemma
\ref{decent-spaces-lemma-decent-space-elementary-etale-neighbourhood}에서처럼 택하자.
$x$는 $X$의 닫힌점이고 $w$는 $W$에서 $x$ 위에 놓이는
유일한 점이므로, $w$는 $W$의 닫힌점이다.
$a$는 에탈이고 $x$와 $w$에서 잉여체를 동일시하므로,
$a$는 동형 $a^{-1}x \to x$를 유도한다
($X$와 $W$의 닫힌 부분공간으로서). 따라서 Lemmas
\ref{spaces-limits-lemma-excision-modifications}와
\ref{spaces-limits-lemma-excision-modifications-properties}를 적용하여
$X$가 아핀 스킴인 경우로 환원할 수 있다.

\medskip\noindent
$X$가 아핀 스킴이라고 가정하자. $\mathcal{O}_{X, x}^h$는
$\Gamma(U, \mathcal{O}_U)$들의 여극한임을 상기하자. 여기서 여극한은
아핀 기본 에탈 근방 $(U, u) \to (X, x)$들을 따라 취한다.
이 근방들의 범주는 공여과되어 있다. Decent Spaces, Lemma
\ref{decent-spaces-lemma-elementary-etale-neighbourhoods} 또는
More on Morphisms, Lemma
\ref{more-morphisms-lemma-elementary-etale-neighbourhoods}를 보라.
그러면 같은 범주를 따라 극한을 취할 때
$\Spec(\mathcal{O}_{X, x}^h) = \lim U$이고
$V = \lim U \setminus \{u\}$이다
(Lemma \ref{spaces-limits-lemma-directed-inverse-system-has-limit}). 따라서 Lemma
\ref{spaces-limits-lemma-descend-finite-presentation}에 의해 오른쪽 범주는
쌍 $(U, u)$에 대한 범주들의 여극한이다.
% P09-ERR-0124: frozen source's broken sentence boundary is normalized only by Korean syntax.
그리고 첫 문단의 내용에 의해 이 범주들은 각각 쌍 $(X, x)$에 대한 범주와
동치이다. 이로써 증명이 끝난다.
\end{proof}






\section{보편 닫힌 사상}
\label{spaces-limits-section-universally-closed}

\noindent
이 절에서는 준콤팩트이지만 반드시 분리일 필요는 없는 사상이
언제 보편 닫힌 사상인지 논의한다. 먼저 국소 유한 표시인 밑변환 뒤에
보편 닫힘을 확인할 수 있게 해 주는 보조정리를 증명한다.

\begin{lemma}
\label{spaces-limits-lemma-separate}
$S$를 스킴이라 하자. $f : X \to Y$와 $g : Z \to Y$를
$S$ 위의 대수공간들의 사상이라 하자. $z \in |Z|$라 하고,
$T \subset |X \times_Y Z|$를
$z \not \in \Im(T \to |Z|)$인 닫힌 부분집합이라 하자.
$f$가 준콤팩트이면, 에탈 근방 $(V, v) \to (Z, z)$와 가환 그림
$$
\xymatrix{
V \ar[d] \ar[r]_a & Z' \ar[d]^b \\
Z \ar[r]^g & Y,
}
$$
및 닫힌 부분집합 $T' \subset |X \times_Y Z'|$가 존재하여 다음을 만족한다.
\begin{enumerate}
\item 사상 $b : Z' \to Y$는 국소 유한 표시이고,
\item $z' = a(v)$로 놓으면 $z' \not \in \Im(T' \to |Z'|)$이며,
\item $T$의 $|X \times_Y V|$ 안에서의 역상은 $T'$로 가며,
이때 $|X \times_Y V| \to |X \times_Y Z'|$를 사용한다.
\end{enumerate}
또한 $V$와 $Z'$를 아핀 스킴으로 가정할 수 있으며,
$Z$가 스킴이면 $V$를 $z$의 아핀 열린 근방으로 가정할 수 있다.
\end{lemma}

\begin{proof}
스킴의 사상에 대한 대응하는 결과에서 이를 유도한다.
$y \in |Y|$를 $z$의 상이라 하자. 먼저 아핀 에탈 근방
$(U, u) \to (Y, y)$를 택하고, 이어서 아핀 에탈 근방
$(V, v) \to (Z, z)$를 택하되 사상 $V \to Y$가 $U$를 통해
인수분해되게 한다.
그러면 다음과 같이 대체할 수 있다.
\begin{enumerate}
\item $X \to Y$를 $X \times_Y U \to U$로,
\item $Z \to Y$를 $V \to U$로,
\item $z$를 $v$로,
\item $T$를
$|(X \times_Y U) \times_U V| = |X \times_Y V|$ 안의 그 역상으로 대체한다.
\end{enumerate}
실제로 아래에서 $V$를 $v$의 아핀 열린 근방으로 대체한 뒤,
사상 $a : V \to Z'$를 얻을 것이며, 여기서 $Z' \to U$는 유한 표시이다.
또한 닫힌 부분집합 $T'$를
$|(X \times_Y U) \times_U Z'| = |X \times_Y Z'|$ 안에서 찾아
$T$가 $T'$로 가고 $a(v) \not \in \Im(T' \to |Z'|)$가 되게 할 것이다.
따라서 $Z$와 $Y$가 아핀 스킴이라고 가정해도 된다. 다만 해답에서
$V$는 $z$의 열린 근방이어야 한다.

\medskip\noindent
$f$가 준콤팩트이고 $Y$가 아핀이므로 대수공간 $X$는 준콤팩트이다.
아핀 스킴 $W$와 전사 에탈 사상 $W \to X$를 택하자.
$T_W \subset |W \times_Y Z|$를 $T$의 역상이라 하자.
그러면 $z$는 $T_W$의 상에 속하지 않는다. 스킴의 경우
(Limits, Lemma \ref{limits-lemma-separate})에 의해
열린 근방 $V \subset Z$를 $z$에 대해 택하고, 스킴들의 가환 그림
% P09-ERR-0125: frozen source's missing list comma is immaterial in Korean syntax.
$$
\xymatrix{
V \ar[d] \ar[r]_a & Z' \ar[d]^b \\
Z \ar[r]^g & Y,
}
$$
및 닫힌 부분집합 $T' \subset |W \times_Y Z'|$를 찾아 다음이 성립하게 할 수 있다.
\begin{enumerate}
\item 사상 $b : Z' \to Y$는 국소 유한 표시이고,
\item $z' = a(z)$로 놓으면 $z' \not \in \Im(T' \to Z')$이며,
% P09-ERR-0126: frozen source's unbarred Z' image target is retained literally.
\item $T_1 = T_W \cap |W \times_Y V|$는 $T'$로 가며,
이때 $|W \times_Y V| \to |W \times_Y Z'|$를 사용한다.
\end{enumerate}
가환 그림
$$
\xymatrix{
W \times_Y Z \ar[d] &
W \times_Y V \ar[l] \ar[rr]_{a_1} \ar[d]_c & &
W \times_Y Z' \ar[d]^q \\
X \times_Y Z &
X \times_Y V \ar[l] \ar[rr]^{a_2} & &
X \times_Y Z'
}
$$
의 정사각형들은 카르테시안이고 수직 사상들은 전사 에탈이며,
따라서 특히 열린 사상이다.
% P09-ERR-0127: frozen source's spurious comma after “are” is immaterial in Korean.
왼쪽 정사각형을 보면
$T_1 = T_W \cap |W \times_Y V|$는
$T_2 = T \cap |X \times_Y V|$의 $c$에 의한 역상이다. Properties of Spaces, Lemma
\ref{spaces-properties-lemma-points-cartesian}에 의해
$a_1(T_1) = q^{-1}(a_2(T_2))$를 얻는다.
Topology, Lemma \ref{topology-lemma-open-morphism-quotient-topology}에 의해
$$
q^{-1}\left(\overline{a_2(T_2)}\right) =
\overline{q^{-1}(a_2(T_2))} =
\overline{a_1(T_1)} \subset T'
$$
이다. $q$가 전사이므로 $\overline{a_2(T_2)} \to |Z'|$의 상은
$z'$를 포함하지 않는다. 실제로 $T'$에 대해 같은 사실이 성립한다.
따라서 위의 $Z', V, a, b$로 이루어진 그림과 닫힌 부분집합
$\overline{a_2(T_2)} \subset |X \times_Y Z'|$를
보조정리에서 요구한 해답으로 취할 수 있다.
\end{proof}

\begin{lemma}
\label{spaces-limits-lemma-test-universally-closed}
$S$를 스킴이라 하자. $f : X \to Y$를
$S$ 위의 대수공간들의 준콤팩트 사상이라 하자. 다음은 동치이다.
\begin{enumerate}
\item $f$는 보편 닫힌 사상이다.
\item 국소 유한 표시인 모든 사상 $Z \to Y$에 대해
$|X \times_Y Z| \to |Z|$는 닫힌 사상이다.
\item 스킴 $V$와 전사 에탈 사상 $V \to Y$가 존재하여
$|\mathbf{A}^n \times (X \times_Y V)| \to |\mathbf{A}^n \times V|$가
모든 $n \geq 0$에 대해 닫힌 사상이다.
\end{enumerate}
\end{lemma}

\begin{proof}
(1)이 (2)를 함의함은 분명하다.
$|X \times_Y Z| \to |Z|$가 닫힌 사상이 아니라고 가정하자.
여기서 $Z \to Y$는 $S$ 위의 대수공간들의 어떤 사상이다.
이는 닫힌 부분집합 $T \subset |X \times_Y Z|$가 존재하여
$\Im(T \to |Z|)$가 닫히지 않았다는 뜻이다.
$z \in |Z|$를 택하되, 이는 $T$의 상의 폐포에는 속하지만
그 상에는 속하지 않게 하자. Lemma \ref{spaces-limits-lemma-separate}를 적용하면,
에탈 근방 $(V, v) \to (Z, z)$와 가환 그림
$$
\xymatrix{
V \ar[d] \ar[r]_a & Z' \ar[d]^b \\
Z \ar[r]^g & Y,
}
$$
및 닫힌 부분집합 $T' \subset |X \times_Y Z'|$를 얻으며 다음이 성립한다.
\begin{enumerate}
\item 사상 $b : Z' \to Y$는 국소 유한 표시이고,
\item $z' = a(v)$로 놓으면 $z' \not \in \Im(T' \to |Z'|)$이며,
\item $T$의 $|X \times_Y V|$ 안에서의 역상은 $T'$로 가며,
이때 $|X \times_Y V| \to |X \times_Y Z'|$를 사용한다.
\end{enumerate}
$z'$가 $\Im(T' \to |Z'|)$의 폐포에 속한다고 주장한다.
그러면 $|X \times_Y Z'| \to |Z'|$가 닫힌 사상이 아니게 된다.
이 주장은 (2)가 (1)을 함의함을 보인다. 이를 확인하기 위해
다음 가환 그림을 생각하자.
% P09-ERR-0128: frozen source's missing definite article is immaterial in Korean syntax.
$$
\xymatrix{
X \times_Y Z \ar[d] &
X \times_Y V \ar[l] \ar[d] \ar[r] &
X \times_Y Z' \ar[d] \\
Z & V \ar[l] \ar[r]^a & Z'
}
$$
$T_V \subset |X \times_Y V|$를 $T$의 역상이라 하자.
Properties of Spaces, Lemma \ref{spaces-properties-lemma-points-cartesian}에 의해
$T_V$의 $|V|$ 안에서의 상은 $T$의 $|Z|$ 안에서의 상의 역상이다.
이제 $z$는 $T \to |Z|$의 상의 폐포에 속하고 $|V| \to |Z|$는
열린 사상이므로, $v$는 $T_V \to |V|$의 상의 폐포에 속한다.
$T_V$의 $|X \times_Y Z'|$ 안에서의 상은 $|T'|$에 포함되므로,
% P09-ERR-0129: frozen source's underlying-space notation |T'| is retained literally.
$z' = a(v)$는 $T'$의 상의 폐포에 속한다.

\medskip\noindent
(1)이 (3)을 함의함은 분명하다. $V \to Y$가 (3)에서와 같다고 하자.
$X \times_Y V \to V$가 보편 닫힌 사상임을 보일 수 있으면,
Morphisms of Spaces, Lemma
\ref{spaces-morphisms-lemma-universally-closed-local}에 의해
$f$는 보편 닫힌 사상이다. 따라서 $f : X \to Y$가 (2)를 만족함을
보이면 충분하다. 여기서 $f$는 대수공간들의 준콤팩트 사상이고,
$Y$는 스킴이며, $|\mathbf{A}^n \times X| \to |\mathbf{A}^n \times Y|$는
모든 $n$에 대해 닫힌 사상이라고 가정한다.
국소 유한 표시인 사상 $Z \to Y$를 택하자.
$|X \times_Y Z| \to |Z|$가 닫힌 사상임을 보여야 한다.
이 문제는 $Z$에서 에탈 국소적이므로 $Z$가 아핀이라고 가정해도 된다
(몇 가지 세부사항은 생략한다). $Y$가 스킴이고 $Z$가 아핀이며
$Z \to Y$가 국소 유한 표시이므로 몰입
$Z \to \mathbf{A}^n \times Y$를 찾을 수 있다. Morphisms, Lemma
\ref{morphisms-lemma-quasi-affine-finite-type-over-S}를 보라.
카르테시안 그림
$$
\vcenter{
\xymatrix{
X \times_Y Z \ar[d] \ar[r] & \mathbf{A}^n \times X \ar[d] \\
Z \ar[r] & \mathbf{A}^n \times Y
}
}
\quad
\begin{matrix}
\text{이는 다음을} \\
\text{유도한다}
\end{matrix}
\quad
\vcenter{
\xymatrix{
|X \times_Y Z| \ar[d] \ar[r] & |\mathbf{A}^n \times X| \ar[d] \\
|Z| \ar[r] & |\mathbf{A}^n \times Y|
}
}
$$
을 생각하자. 이 위상공간 그림의 수평 화살표들은 국소 닫힌 부분집합 위의
위상동형이다(Properties of Spaces, Lemma
\ref{spaces-properties-lemma-subspace-induced-topology}). 따라서
$T$를 $|X \times_Y Z|$의 임의의 닫힌 부분집합이라 하자. 이는
닫힌 부분집합 $T'$에서 당겨오며, 그 부분집합은
$|\mathbf{A}^n \times Y|$ 안에 있다.
또한 가정에 의해 $T'$의 $|\mathbf{A}^n \times X|$ 안에서의 상은 닫혀 있다.
% P09-ERR-0130: frozen source's reversed ambient spaces for T' and its image are retained literally.
그러므로 원하는 대로 $T$의 $|Z|$ 안에서의 상은 닫혀 있다.
\end{proof}

\begin{lemma}
\label{spaces-limits-lemma-limited-base-change}
$S$를 스킴이라 하자. $f : X \to Y$를 $S$ 위의 대수공간들의
사상이라 하자. $f$가 분리이고 유한형이라고 가정하자.
% P09-ERR-0131: frozen source's missing copula is immaterial in Korean syntax.
다음은 동치이다.
\begin{enumerate}
\item 사상 $f$는 고유이다.
\item 국소 유한 표시인 임의의 사상 $Y \to Z$에 대해
$|X \times_Y Z| \to |Z|$는 닫힌 사상이다.
% P09-ERR-0132: frozen source's reversed Y-to-Z base-change direction is retained literally.
\item 스킴 $V$와 전사 에탈 사상 $V \to Y$가 존재하여
$|\mathbf{A}^n \times (X \times_Y V)| \to |\mathbf{A}^n \times V|$가
모든 $n \geq 0$에 대해 닫힌 사상이다.
\end{enumerate}
\end{lemma}

\begin{proof}
고유 사상은 분리이고 유한형이며 보편 닫힌 사상인 것과 같은 개념이므로,
이 보조정리는 Lemma \ref{spaces-limits-lemma-test-universally-closed}의 특수한 경우이다.
\end{proof}


\section{Noether 부치환 판정법}
\label{spaces-limits-section-Noetherian-valuative-criterion}

\noindent
Cohomology of Spaces, Section
\ref{spaces-cohomology-section-Noetherian-valuative-criterion}에서 이미
몇 가지 결과를 증명했다. 스킴에 대한 대응하는 절은
Limits, Section \ref{limits-section-Noetherian-valuative-criterion}이다.

\medskip\noindent
이 절의 많은 결과는 다음 보조정리를 인용하여 증명할 수 있고
(어쩌면 그렇게 해야 하지만), 항상 그렇게 하지는 않았다.

\begin{lemma}
\label{spaces-limits-lemma-reach-point-closure-Noetherian}
$S$를 스킴이라 하자. $f : X \to Y$를 $S$ 위의 대수공간들의
사상이라 하자. $f$가 유한형이고 $Y$가 국소 Noether라고 가정하자.
% P09-ERR-0133: frozen source's missing copula is immaterial in Korean syntax.
$y \in |Y|$를 $|f|$의 상의 폐포에 속하는 점이라 하자.
그러면 가환 그림
$$
\xymatrix{
\Spec(K) \ar[r] \ar[d] & X \ar[d]^f \\
\Spec(A) \ar[r] & Y
}
$$
가 존재한다. 여기서 $A$는 이산 부치환환이고 $K$는 그 분수체이며,
$\Spec(A)$의 닫힌점을 $y$로 보낸다. 또한
$x \in |X|$를 $\Spec(K) \to X$에 대응하는 점이라 하면,
이 점이 여차원 $0$인 점이라고
가정할 수 있다\footnote{Properties of Spaces, Section
\ref{spaces-properties-section-generic-points}의 논의를 보라.}.
또한 $K$는 $X$ 위에서 에탈인 어떤 스킴의 점의 잉여체라고 가정할 수 있다.
\end{lemma}

\begin{proof}
아핀 스킴 $V$, 점 $v \in V$, 그리고 에탈 사상 $V \to Y$를 택하자.
이 사상은 $v$를 $y$로 보낸다. 사상 $|V| \to |Y|$는 열려 있고,
Properties of Spaces, Lemma \ref{spaces-properties-lemma-points-cartesian}에 의해
$|X \times_Y V| \to |V|$의 상은 $|f|$의 상의 역상이다.
따라서 점 $v$는 $|X \times_Y V| \to |V|$의 상의 폐포에 속한다.
$X \times_Y V \to V$와 점 $v$에 대해 보조정리를 증명하면,
$f$와 $y$에 대해서도 보조정리가 따른다. 이와 같이 다음 문단에서
설명하는 상황으로 환원한다.

\medskip\noindent
$f : X \to Y$와 $y \in |Y|$가 보조정리에서와 같고
$Y$가 아핀 스킴이라고 가정하자. $f$가 준콤팩트이므로 $X$는 준콤팩트이다.
따라서 아핀 스킴 $W$와 전사 에탈 사상 $W \to X$를 택할 수 있다.
그러면 $|f|$의 상은 $W \to Y$의 상과 같다. 이로써 스킴의 경우로
환원되며, 그 경우는 Limits, Lemma
\ref{limits-lemma-reach-point-closure-Noetherian}이다.
\end{proof}

\noindent
먼저 분리에 관한 결과를 진술한다. 밑스킴 $S$ 위의 대수공간들의
사상으로 이루어진 다음 모양의 실선 가환 그림을 자주 사용할 것이다.
\begin{equation}
\label{spaces-limits-equation-valuative}
\vcenter{
\xymatrix{
\Spec(K) \ar[r] \ar[d] & X \ar[d] \\
\Spec(A) \ar[r] \ar@{-->}[ru] & Y
}
}
\end{equation}
여기서 $A$는 부치환환이고 $K$는 그 분수체이다.

\begin{lemma}
\label{spaces-limits-lemma-Noetherian-dvr-valuative-separation}
$S$를 스킴이라 하자. $f : X \to Y$를 $S$ 위의 대수공간들의
사상이라 하자. $f$가 준분리이고 국소 유한형이며
$Y$가 국소 Noether라고 가정하자. 다음은 동치이다.
\begin{enumerate}
\item 사상 $f$는 분리이다.
\item 임의의 그림 (\ref{spaces-limits-equation-valuative})에서 점선 화살표는 많아야 하나이다.
\item $A$가 이산 부치환환인 모든 그림 (\ref{spaces-limits-equation-valuative})에서
점선 화살표는 많아야 하나이다.
\item $A$가 이산 부치환환이고 $\Spec(K) \to X$의 상이
여차원 $0$인 $X$의 점인 모든 그림 (\ref{spaces-limits-equation-valuative})에서
점선 화살표는 많아야 하나이다.
\end{enumerate}
\end{lemma}

\begin{proof}
Morphisms of Spaces, Lemma
\ref{spaces-morphisms-lemma-separated-implies-valuative}에 의해
(1) $\Rightarrow$ (2)이다. (2) $\Rightarrow$ (3)과
(3) $\Rightarrow$ (4)는 자명하다. (4)가 (1)을 함의함을 보이면 된다.

\medskip\noindent
(4)를 가정하자. 대각사상
$\Delta : X \to X \times_Y X$가 닫힌 몰입임을 보여야 한다.
$\Delta$가 표현 가능하고 분리이며 단사상이고 국소 유한형임은 이미 안다.
Morphisms of Spaces, Lemma
\ref{spaces-morphisms-lemma-properties-diagonal}를 보라.
아핀 스킴 $U$와 에탈 사상 $U \to X \times_Y X$를 택하자.
$V = X \times_{\Delta, X \times_Y X} U$로 놓는다.
$V \to U$가 닫힌 몰입임을 보이면 충분하다
(Morphisms of Spaces, Lemma
\ref{spaces-morphisms-lemma-closed-immersion-local}).
$X \times_Y X$는 $Y$ 위에서 국소 유한형이므로 $U$는 Noether이다
(Morphisms of Spaces, Lemmas
\ref{spaces-morphisms-lemma-composition-finite-type},
\ref{spaces-morphisms-lemma-base-change-finite-type}, 및
\ref{spaces-morphisms-lemma-locally-finite-type-locally-noetherian}을 사용한다).
$V$는 스킴임에 주의하자. 이는 $\Delta$가 표현 가능하기 때문이다.
또한 $V$는 준콤팩트이다. 이는 $f$가 준분리이기 때문이다.
따라서 $V \to U$는 분리이고 유한형이다. 스킴의 사상들의 가환 그림
$$
\xymatrix{
\Spec(K) \ar[r] \ar[d] & V \ar[d] \\
\Spec(A) \ar[r] \ar@{-->}[ru] & U
}
$$
을 생각하자. 여기서 $A$는 분수체 $K$를 갖는 이산 부치환환이고,
$K$는 Noether 스킴 $V$의 일반점 하나의 잉여체이다.
$V \to X$는 에탈이므로(에탈 사상 $U \to X \times_Y X$의 밑변환),
$\Spec(K) \to V \to X$의 상은 여차원 $0$인 점이다.
Properties of Spaces, Section
\ref{spaces-properties-section-dimension-local-ring}을 보라.
합성 $\Spec(A) \to U \to X \times_Y X$를 사상들의 쌍
$a, b : \Spec(A) \to X$로 해석할 수 있다. 이 사상들은 $Y$로 가는
사상으로서 일치하고, $\Spec(K)$로 제한하면 같으며, 이 제한은
여차원 $0$인 점으로 간다. 따라서 가정 (4)에 의해 $a = b$이고,
그림의 점선 화살표를 얻는다. Limits, Lemma
\ref{limits-lemma-Noetherian-dvr-valuative-proper}에 의해
$V \to U$는 고유이다. 다시 말해 $\Delta$는 고유이다.
$\Delta$가 단사상이므로 $\Delta$는 원하는 대로 닫힌 몰입이다
(\'Etale Morphisms, Lemma
\ref{etale-lemma-characterize-closed-immersion}).
\end{proof}

\begin{lemma}
\label{spaces-limits-lemma-Noetherian-dvr-valuative-proper}
$S$를 스킴이라 하자. $f : X \to Y$를 $S$ 위의 대수공간들의
사상이라 하자. $f$가 준분리이고 유한형이며 $Y$가 국소 Noether라고
가정하자. 다음은 동치이다.
\begin{enumerate}
\item $f$는 고유이다.
\item $f$는 부치환 판정법을 만족한다. Morphisms of Spaces, Definition
\ref{spaces-morphisms-definition-valuative-criterion}을 보라.
\item 임의의 그림 (\ref{spaces-limits-equation-valuative})에 점선 화살표가 정확히 하나 존재한다.
\item $A$가 이산 부치환환인 모든 그림 (\ref{spaces-limits-equation-valuative})에
점선 화살표가 정확히 하나 존재한다.
\item $A$가 이산 부치환환이고 $\Spec(K) \to X$의 상이
여차원 $0$인 $X$의 점인 모든 그림 (\ref{spaces-limits-equation-valuative})에
점선 화살표가 정확히 하나 존재한다\footnote{존재성 부분에서
이산 부치환환을 확대한 뒤 점선 화살표가 존재하기만 하면 된다는
더 정밀한 정식화가 있다.}.
% P09-ERR-0134: Korean uses one terminal punctuation mark; the frozen source doubles it around this footnote.
\end{enumerate}
\end{lemma}

\begin{proof}
Morphisms of Spaces, Lemma
\ref{spaces-morphisms-lemma-characterize-proper}에 의해
(1) $\Leftrightarrow$ (2) $\Leftrightarrow$ (3)이다.
(3) $\Rightarrow$ (4) $\Rightarrow$ (5)임은 분명하다.
증명을 끝내기 위해 이제 (5)가 (1)을 함의함을 보인다.

\medskip\noindent
(5)를 가정하자. Lemma \ref{spaces-limits-lemma-Noetherian-dvr-valuative-separation}에 의해
$f$는 분리이다. 증명을 끝내려면 $f$가 보편 닫힌 사상임을 보이면 충분하다.
$V \to Y$를 에탈 사상이라 하고 $V$는 아핀 스킴이라고 하자.
밑변환 $V \times_Y X \to V$가 보편 닫힌 사상임을 보이면 충분하다.
Morphisms of Spaces, Lemma
\ref{spaces-morphisms-lemma-universally-closed-local}을 보라.
다음 그림
$$
\xymatrix{
\Spec(K) \ar[r] \ar[d] & V \times_Y X \ar[d] \ar[r] & X \ar[d] \\
\Spec(A) \ar[r] \ar@{-->}[ru] \ar@{..>}[rru] & V \ar[r] & Y
}
$$
을 생각하자. 이는 $S$ 위의 대수공간들의 가환 그림이다.
여기서 $A$는 분수체 $K$를 갖는 이산 부치환환이고,
$\Spec(K) \to V \times_Y X$는 여차원 $0$인 대수공간
$V \times_Y X$의 점으로 간다. $V \times_Y X \to X$가 에탈이므로
$\Spec(K) \to X$의 상은 여차원 $0$인 $X$의 점이다.
따라서 (5)에 의해 그림의 두 점선 화살표 중 더 긴 것을 얻고,
물론 더 짧은 것도 얻는다. 그러므로 사상 $V \times_Y X \to V$에 대해
가정들이 성립하며, 다음 문단에서 논의하는 경우로 환원된다.

\medskip\noindent
$Y$가 Noether 아핀 스킴이라고 가정하자.
% P09-ERR-0135: frozen source's “Aassume” typo is normalized in Korean.
이 경우 $X$는 분리 Noether 대수공간이며
($f$가 분리임은 이미 안다) $Y$ 위에서 유한형이다.
(특히 대수공간 $X$는 스킴인 조밀 열린 부분공간을 가진다. Properties of Spaces,
Proposition
\ref{spaces-properties-proposition-locally-quasi-separated-open-dense-scheme}을 보라.
엄밀히 말해 이 사실은 필요하지 않다.)
준사영 스킴 $X'$를 $Y$ 위에서 택하고, 고유 전사 사상 $X' \to X$도
Chow 보조정리의 약한 형태
(Cohomology of Spaces, Lemma \ref{spaces-cohomology-lemma-weak-chow})에서처럼 택하자.
$X'$를 $X$의 한 기약성분을 지배하는 기약성분들의 분리합으로 대체해도 된다.
세부사항은 생략한다. 특히 스킴 $X'$의 일반점들이 여차원 $0$인
$X$의 점들로 가게 할 수 있다(이 경우 이 점들은 정확히 $X$의 일반점들이다).
$X' \to Y$가 고유라고 주장한다. 이 주장은 Morphisms of Spaces, Lemma
\ref{spaces-morphisms-lemma-image-proper-is-proper}에 의해
$X$가 $Y$ 위에서 고유임을 함의한다. 이를 증명하기 위해 Limits, Lemma
\ref{limits-lemma-Noetherian-dvr-valuative-proper}에 따라 모든 실선 가환 그림
$$
\xymatrix{
\Spec(K) \ar[r] \ar[d] & X' \ar[r] & X \ar[d] \\
\Spec(A) \ar[rr] \ar@{-->}[ru]^a \ar@{-->}[rru]_b & & Y
}
$$
에서 다음을 증명하면 충분하다. 여기서 $A$는 분수체 $K$를 갖는
이산 부치환환이고, $K$는 $X'$의 일반점 하나의 잉여체이다.
점선 화살표 $a$를 찾으면 된다($X'$가 분리이므로 유일성은 이미 안다).
가정 (5)에 의해 점선 화살표 $b$를 찾을 수 있다.
그러면 사상 $X' \times_{X, b} \Spec(A) \to \Spec(A)$는
스킴들의 고유 사상이고, 스킴 사상에 대한 부치환 판정법에 의해
$b$를 원하는 사상 $a$로 올릴 수 있다.
\end{proof}

\begin{lemma}
\label{spaces-limits-lemma-check-universally-closed-Noetherian}
$S$를 스킴이라 하자. $f : X \to Y$를 $S$ 위의 대수공간들의
사상이라 하자. $Y$가 국소 Noether이고 $f$가 유한형이라고 가정하자.
그러면 다음은 동치이다.
\begin{enumerate}
\item $f$는 보편 닫힌 사상이다.
\item $f$는 부치환 판정법의 존재성 부분을 만족한다.
\item 스킴 $V$와 전사 에탈 사상 $V \to Y$가 존재하여
$|\mathbf{A}^n \times X \times_Y V| \to |\mathbf{A}^n \times V|$가
모든 $n \geq 0$에 대해 닫힌 사상이다.
\item $A$가 이산 부치환환인 모든 그림 (\ref{spaces-limits-equation-valuative})에 대해
체의 유한 분리 확대 $K'/K$, 이산 부치환환 $A' \subset K'$ 중
$A$를 지배하는 것, 그리고 다음 그림을 가환으로 만드는 사상
$\Spec(A') \to X$가 존재한다.
% P09-ERR-0136: frozen source's duplicated “there” is normalized only by Korean syntax.
$$
\xymatrix{
\Spec(K') \ar[r] \ar[d] & \Spec(K) \ar[r] & X \ar[d] \\
\Spec(A') \ar[r] \ar[rru] & \Spec(A) \ar[r] & Y
}
$$
\item $A$가 이산 부치환환인 모든 그림 (\ref{spaces-limits-equation-valuative})에 대해
체 확대 $K'/K$, 부치환환 $A' \subset K'$ 중 $A$를 지배하는 것, 그리고
다음 그림을 가환으로 만드는 사상 $\Spec(A') \to X$가 존재한다.
% P09-ERR-0137: frozen source's duplicated “there” is normalized only by Korean syntax.
$$
\xymatrix{
\Spec(K') \ar[r] \ar[d] & \Spec(K) \ar[r] & X \ar[d] \\
\Spec(A') \ar[r] \ar[rru] & \Spec(A) \ar[r] & Y
}
$$
\end{enumerate}
\end{lemma}

\begin{proof}
Lemma \ref{spaces-limits-lemma-test-universally-closed}와 Morphisms of Spaces, Lemma
\ref{spaces-morphisms-lemma-quasi-compact-existence-universally-closed}에 의해
(1), (2), (3)은 동치이다. 이 동치 조건들은 (4)를 함의한다.
실제로 Morphisms of Spaces, Lemma
\ref{spaces-morphisms-lemma-finite-separable-enough}에 의해
부치환 판정법의 존재성 부분에서 $K'/K$를 항상 유한 분리 확대가 되게
택할 수 있고, 그러면 Krull--Akizuki에 의해 $A'$는 자동으로
이산 부치환환이다(Algebra, Lemma \ref{algebra-lemma-krull-akizuki}).
(4) $\Rightarrow$ (5)는 자명하다. 나머지 증명에서는 (5)가 (1)을
함의함을 보인다.

\medskip\noindent
(5)를 가정하자. 아핀 스킴 $V$와 에탈 사상 $V \to Y$를 택하자.
% P09-ERR-0138: frozen source's “Chose” typo is normalized in Korean.
$f$의 $V$로의 밑변환이 보편 닫힌 사상임을 보이면 충분하다.
Morphisms of Spaces, Lemma
\ref{spaces-morphisms-lemma-universally-closed-local}을 보라.
Lemma \ref{spaces-limits-lemma-Noetherian-dvr-valuative-proper}의 증명에서와 정확히 같이,
가정 (5)는 이 밑변환으로 계승된다. 세부사항은 생략한다.
이로써 다음 문단에서 논의하는 경우로 환원된다.

\medskip\noindent
$Y$가 Noether 아핀 스킴이고 (5)가 성립한다고 가정하자.
$f$가 보편 닫힌 사상임을 증명하려면
$|X \times \mathbf{A}^n| \to |Y \times \mathbf{A}^n|$가 모든 $n$에 대해
닫힌 사상임을 보이면 충분하다(위의 논의에 의해).
가정 (5)는 곱사상 $X \times \mathbf{A}^n \to Y \times \mathbf{A}^n$으로
계승되므로(세부사항은 생략한다),
$|X| \to |Y|$가 닫힌 사상임을 증명하는 것으로 환원된다.

\medskip\noindent
$Y$가 Noether 아핀 스킴이고 (5)가 성립한다고 가정하자.
$T \subset |X|$를 닫힌 부분집합이라 하자. $T$의 $|Y|$ 안에서의 상이
닫혔음을 보여야 한다. $X$를 $T$ 위의 축소 유도 닫힌 부분공간 구조로
대체해도 된다. 이 대체가 성질 (5)를 보존한다는 확인은 생략한다.
따라서 $|X| \to |Y|$의 상이 닫혔음을 증명하는 것으로 환원된다.

\medskip\noindent
$y \in |Y|$를 $|X| \to |Y|$의 상의 폐포에 속하는 점이라 하자.
Lemma \ref{spaces-limits-lemma-reach-point-closure-Noetherian}에 의해 가환 그림
$$
\xymatrix{
\Spec(K) \ar[r] \ar[d] & X \ar[d]^f \\
\Spec(A) \ar[r] & Y
}
$$
을 택할 수 있다. 여기서 $A$는 이산 부치환환이고 $K$는 그 분수체이며,
$\Spec(A)$의 닫힌점을 $y$로 보낸다. 성질 (5)에 의해 즉시
$y$는 $|X| \to |Y|$의 상에 속한다. 증명이 끝났다.
\end{proof}


















\section{정밀한 Noether 부치환 판정법}
\label{spaces-limits-section-refined-valuative-criteria}

\noindent
이 절은 Limits, Section
\ref{limits-section-refined-valuative-criteria}에 대응한다.
부치환 판정법을 확인할 때 가능한 모든 부치환환 그림을
생각할 필요는 보통 없다.

\begin{lemma}
\label{spaces-limits-lemma-refined-valuative-criterion-proper}
$S$를 스킴이라 하자. $f : X \to Y$와 $h : U \to X$를
$S$ 위의 대수공간들의 사상이라 하자. $Y$가 국소 Noether이고,
$f$와 $h$가 유한형이며, $f$가 분리이고,
$|h| : |U| \to |X|$의 상이 $|X|$에서 조밀하다고 가정하자.
임의의 실선 가환 그림
$$
\xymatrix{
\Spec(K) \ar[r] \ar[d] & U \ar[r]^h & X \ar[d]^f \\
\Spec(A) \ar[rr] \ar@{-->}[rru] & & Y
}
$$
이 주어졌을 때, 여기서 $A$는 분수체 $K$를 갖는 이산 부치환환이고,
그림을 가환으로 만드는 점선 화살표가 존재한다면 $f$는 고유이다.
\end{lemma}

\begin{proof}
$f$가 보편 닫힌 사상임을 보이면 충분하다.
$V \to Y$를 에탈 사상이라 하고 $V$는 아핀 스킴이라고 하자.
Morphisms of Spaces, Lemma
\ref{spaces-morphisms-lemma-universally-closed-local}에 의해 밑변환
$X \times_Y V \to V$가 보편 닫힌 사상임을 보이면 충분하다.
Properties of Spaces, Lemma
\ref{spaces-properties-lemma-points-cartesian}에 의해 $I$를
$|U \times_Y V| \to |X \times_Y V|$의 상이라 하면,
이는 $|h|$의 상의 역상이다.
$|X \times_Y V| \to |X|$가 열린 사상이므로
(Properties of Spaces, Lemma \ref{spaces-properties-lemma-etale-open}),
$I$는 $|X \times_Y V|$에서 조밀하다. 따라서 보조정리의 가정들은
사상들 $U \times_Y V \to X \times_Y V \to V$에 대해 성립한다.
그러므로 $Y$가 아핀 스킴이라고 가정해도 된다.

\medskip\noindent
$Y$가 아핀 스킴이라고 가정하자. 그러면 $U$는 준콤팩트이다.
아핀 스킴 하나와 전사 에탈 사상 $W \to U$를 택하자.
% P09-ERR-0139: frozen source leaves the affine scheme unnamed before introducing W; retained literally.
이제 $U$를 $W$로 대체하고 $U$가 아핀이라고 가정해도 된다.
Chow 보조정리의 약한 형태
(Cohomology of Spaces, Lemma \ref{spaces-cohomology-lemma-weak-chow})에 의해
전사 고유 사상 $X' \to X$를 택할 수 있으며, 여기서 $X'$는 스킴이다.
그러면 $U' = X' \times_X U$는 스킴이고 $U' \to X'$는 유한형이다.
$X'$를 $h' : U' \to X'$의 스킴론적 상으로 대체할 수 있으며,
따라서 $h'(U')$는 $X'$에서 조밀하다. 모든 그림
$$
\xymatrix{
\Spec(K) \ar[r] \ar[d] & U' \ar[r]^h & X' \ar[d]^{f'} \\
\Spec(A) \ar[rr] \ar@{-->}[rru] & & Y
}
$$
% P09-ERR-0140: frozen diagram label h is retained rather than silently changed to h'.
에서 다음이 성립한다고 주장한다. 여기서 $A$는 분수체 $K$를 갖는
이산 부치환환이며, 그림을 가환으로 만드는 점선 화살표가 존재한다.
실제로 먼저 보조정리의 가정에 의해 화살표 $\Spec(A) \to X$를 얻고,
이어서 고유성의 부치환 판정법을 사용하여 이를 화살표
$\Spec(A) \to X'$로 올린다(Morphisms of Spaces, Lemma
\ref{spaces-morphisms-lemma-characterize-proper}).
사상 $X' \to Y$는 고유 사상과 분리 사상의 합성이므로 분리이다.
따라서 스킴의 경우에 의해 $X' \to Y$는 고유이다
(Limits, Lemma
\ref{limits-lemma-refined-valuative-criterion-proper}).
Morphisms of Spaces, Lemma
\ref{spaces-morphisms-lemma-image-proper-is-proper}에 의해
$X \to Y$가 고유라고 결론짓는다.
\end{proof}

\begin{lemma}
\label{spaces-limits-lemma-refined-valuative-criterion-separated}
$S$를 스킴이라 하자. $f : X \to Y$와 $h : U \to X$를
$S$ 위의 대수공간들의 사상이라 하자. $Y$가 국소 Noether이고,
$f$가 국소 유한형이고 준분리이며, $h$가 유한형이고,
$|h| : |U| \to |X|$의 상이 $|X|$에서 조밀하다고 가정하자.
임의의 실선 가환 그림
$$
\xymatrix{
\Spec(K) \ar[r] \ar[d] & U \ar[r]^h & X \ar[d]^f \\
\Spec(A) \ar[rr] \ar@{-->}[rru] & & Y
}
$$
이 주어졌을 때, 여기서 $A$는 분수체 $K$를 갖는 이산 부치환환이고,
그림을 가환으로 만드는 점선 화살표가 많아야 하나라면 $f$는 분리이다.
\end{lemma}

\begin{proof}
Lemma \ref{spaces-limits-lemma-refined-valuative-criterion-proper}를 사상들
$U \to X$와 $\Delta : X \to X \times_Y X$에 적용한다.
조건들을 확인하자. $\Delta$는 준콤팩트이다. 이는 $f$가 준분리이기 때문이다.
물론 $\Delta$는 국소 유한형이고 분리이다(이는 모든 대각사상에 참이다).
마지막으로 실선 가환 그림
$$
\xymatrix{
\Spec(K) \ar[r] \ar[d] & U \ar[r]^h & X \ar[d]^\Delta \\
\Spec(A) \ar[rr]^{(a, b)} \ar@{-->}[rru] & & X \times_Y X
}
$$
이 주어졌다고 하자. 여기서 $A$는 분수체 $K$를 갖는 이산 부치환환이다.
그러면 $a$와 $b$는 보조정리의 그림에서 두 점선 화살표를 주므로 같아야 한다.
따라서 점선 화살표로 $a = b$를 사용할 수 있으며, 이것이 존재성을 준다.
이로써 증명이 끝난다.
\end{proof}

\begin{lemma}
\label{spaces-limits-lemma-refined-valuative-criterion-universally-closed}
$S$를 스킴이라 하자. $f : X \to Y$와 $h : U \to X$를
$S$ 위의 대수공간들의 사상이라 하자. $Y$가 국소 Noether이고,
$f$와 $h$가 유한형이며, $f$가 준분리이고,
$h(U)$가 $X$에서 조밀하다고 가정하자. 임의의 실선 가환 그림
$$
\xymatrix{
\Spec(K) \ar[r] \ar[d] & U \ar[r]^h & X \ar[d]^f \\
\Spec(A) \ar[rr] \ar@{-->}[rru] & & Y
}
$$
이 주어졌을 때, 여기서 $A$는 분수체 $K$를 갖는 이산 부치환환이고,
그림을 가환으로 만드는 점선 화살표가 유일하게 존재한다면 $f$는 고유이다.
\end{lemma}

\begin{proof}
Lemmas \ref{spaces-limits-lemma-refined-valuative-criterion-separated}와
\ref{spaces-limits-lemma-refined-valuative-criterion-proper}를 결합하라.
\end{proof}







\section{유한형 공간의 하강}
\label{spaces-limits-section-finite-type-quasi-separated}

\noindent
이 절에서는 Section \ref{spaces-limits-section-finite-type-closed-in-finite-presentation}의
주제를 Section \ref{spaces-limits-section-descending-relative}에서 논의한 결과들의
정신에 따라 잇는다.
또한 대수공간에 대한 Limits, Section
\ref{limits-section-finite-type-quasi-separated}의 유사물이다.

\begin{situation}
\label{spaces-limits-situation-limit-noetherian}
$S$를 스킴이라 하자. 예를 들어 $\Spec(\mathbf{Z})$로 둘 수 있다.
$B = \lim_{i \in I} B_i$를 $S$ 위의 Noether 공간들의 유향 역계의
극한이라 하자. 전이 사상 $B_{i'} \to B_i$는 $i' \geq i$일 때 아핀이다.
\end{situation}

\begin{lemma}
\label{spaces-limits-lemma-good-diagram}
Situation \ref{spaces-limits-situation-limit-noetherian}의 상황이라 하자.
$X \to B$를 준분리 유한형인 대수공간들의 사상이라 하자.
그러면 $i \in I$와 그림
\begin{equation}
\label{spaces-limits-equation-good-diagram}
\vcenter{
\xymatrix{
X \ar[r] \ar[d] & W \ar[d] \\
B \ar[r] & B_i
}
}
\end{equation}
가 존재하여 $W \to B_i$는 유한형이고, 유도된 사상
$X \to B \times_{B_i} W$는 닫힌 몰입이다.
\end{lemma}

\begin{proof}
Lemma \ref{spaces-limits-lemma-finite-type-closed-in-finite-presentation}에 의해
닫힌 몰입 $X \to X'$를 $B$ 위에서 찾을 수 있으며,
$X'$는 $B$ 위의 유한 표시 대수공간이다.
Lemma \ref{spaces-limits-lemma-descend-finite-presentation}에 의해 $i$와
유한 표시 사상 $X'_i \to B_i$를 찾을 수 있고, 그 당김은 $X'$이다.
$W = X'_i$로 놓자.
\end{proof}

\begin{lemma}
\label{spaces-limits-lemma-limit-from-good-diagram}
Situation \ref{spaces-limits-situation-limit-noetherian}의 상황이라 하자.
$X \to B$를 준분리 유한형인 대수공간들의 사상이라 하자.
$i \in I$와 (\ref{spaces-limits-equation-good-diagram})에서와 같은 그림
$$
\vcenter{
\xymatrix{
X \ar[r] \ar[d] & W \ar[d] \\
B \ar[r] & B_i
}
}
$$
가 주어졌다고 하자. $i' \geq i$에 대해 $X_{i'}$를
$X \to B_{i'} \times_{B_i} W$의 스킴론적 상이라 하자.
그러면 $X = \lim_{i' \geq i} X_{i'}$이다.
\end{lemma}

\begin{proof}
$X$가 준콤팩트이고 준분리이므로
$X \to B_{i'} \times_{B_i} W$의 스킴론적 상의 형성은 에탈 국소화와 가환한다
(Morphisms of Spaces, Lemma
\ref{spaces-morphisms-lemma-quasi-compact-scheme-theoretic-image}).
따라서 $W$가 아핀이고 어떤 아핀 $U_i$로 가며, 이 아핀은
$B_i$ 위에서 에탈이라고 가정해도 된다. 그러면
$$
B_{i'} \times_{B_i} W =
B_{i'} \times_{B_i} U_i \times_{U_i} W =
U_{i'} \times_{U_i} W
$$
이다. 여기서 $U_{i'} = B_{i'} \times_{B_i} U_i$는 전이 사상들이
아핀이므로 아핀이다. 따라서 보조정리는 스킴의 경우, 즉 Limits, Lemma
\ref{limits-lemma-limit-from-good-diagram}에서 따른다.
\end{proof}

\begin{lemma}
\label{spaces-limits-lemma-morphism-good-diagram}
Situation \ref{spaces-limits-situation-limit-noetherian}의 상황이라 하자.
$f : X \to Y$를 $B$ 위에서 준분리이고 유한형인 대수공간들의 사상이라 하자.
다음 두 그림을 택하자.
$$
\vcenter{
\xymatrix{
X \ar[r] \ar[d] & W \ar[d] \\
B \ar[r] & B_{i_1}
}
}
\quad\text{그리고}\quad
\vcenter{
\xymatrix{
Y \ar[r] \ar[d] & V \ar[d] \\
B \ar[r] & B_{i_2}
}
}
$$
이들은 (\ref{spaces-limits-equation-good-diagram})에서와 같은 그림들이다.
Lemma \ref{spaces-limits-lemma-limit-from-good-diagram}에서와 같은 대응하는 극한 표현을
$X = \lim_{i \geq i_1} X_i$ 및
$Y = \lim_{i \geq i_2} Y_i$라 하자.
그러면 $i_0 \geq \max(i_1, i_2)$가 되게 택할 수 있고, 사상
$$
(f_i)_{i \geq i_0} : (X_i)_{i \geq i_0} \to (Y_i)_{i \geq i_0}
$$
이 존재한다. 이는 역계 $(B_i)_{i \geq i_0}$ 위의 사상이며,
$f = \lim_{i \geq i_0} f_i$이다.
% P09-ERR-0141: frozen source's duplicated “such that” is normalized in Korean.
두 번째 역계 사상
$(g_i)_{i \geq i_0} : (X_i)_{i \geq i_0} \to (Y_i)_{i \geq i_0}$가
$(B_i)_{i \geq i_0}$ 위에 있고 $f = \lim_{i \geq i_0} g_i$라면,
% P09-ERR-0142: frozen source's second duplicated “such that” is normalized in Korean.
$f_i = g_i$가 모든 $i \gg i_0$에 대해 성립한다.
\end{lemma}

\begin{proof}
$V \to B_{i_2}$가 유한 표시이고
$X = \lim_{i \geq i_1} X_i$이므로 Proposition
\ref{spaces-limits-proposition-characterize-locally-finite-presentation}과 이를 개선한 Lemma
\ref{spaces-limits-lemma-better-characterize-relative-limit-preserving}를 적용할 수 있다.
따라서 $i_0 \geq \max(i_1, i_2)$가 되게 택하고, 사상
$h : X_{i_0} \to V$를 $B_{i_2}$ 위에서 찾아,
$X \to X_{i_0} \to V$가 $X \to Y \to V$와 같게 할 수 있다.
$i \geq i_0$에 대해 실선 가환 그림
$$
\xymatrix{
X \ar[d] \ar[r] &
X_i \ar[r] \ar@{..>}[d] \ar@/_2pc/[dd] |!{[d];[ld]}\hole &
X_{i_0} \ar[d]^h \\
Y \ar[r] \ar[d] & Y_i \ar[r] \ar[d] & V \ar[d] \\
B \ar[r] & B_i \ar[r] & B_{i_0}
}
$$
을 얻는다. $X \to X_i$의 상은 스킴론적으로 조밀하고,
$Y_i$는 $Y \to B_i \times_{B_{i_2}} V$의 스킴론적 상이므로,
그림이 유도하는 사상 $X_i \to B_i \times_{B_{i_2}} V$는 $Y_i$를 통해
인수분해된다(Morphisms of Spaces, Lemma
\ref{spaces-morphisms-lemma-factor-factor}). 이것이 존재성을 증명한다.

\medskip\noindent
유일성을 보이자. $E_i \to X_i$를 $f_i$와 $g_i$의 동등자라 하자
($i \geq i_0$). 그러면
$E_i = Y_i \times_{\Delta, Y_i \times_{B_i} Y_i, (f_i, g_i)} X_i$이다.
따라서 $E_i \to X_i$는 $Y_i$의 $B_i$ 위 대각사상의 밑변환으로서
유한 표시 단사상이다. Morphisms of Spaces, Lemmas
\ref{spaces-morphisms-lemma-properties-diagonal}와
\ref{spaces-morphisms-lemma-diagonal-morphism-finite-type}를 보라.
$X_i$는 $B_i \times_{B_{i_0}} X_{i_0}$의 닫힌 부분공간이고
$Y_i$도 마찬가지이므로
$$
E_i =
X_i \times_{(B_i \times_{B_{i_0}} X_{i_0})} (B_i \times_{B_{i_0}} E_{i_0}) =
X_i \times_{X_{i_0}} E_{i_0}
$$
이다. 마찬가지로 $X = X \times_{X_{i_0}} E_{i_0}$이다.
따라서 Lemma \ref{spaces-limits-lemma-descend-isomorphism}에 의해
$E_i = X_i$가 충분히 큰 $i$에 대해 성립한다고 결론짓는다.
\end{proof}

\begin{remark}
\label{spaces-limits-remark-finite-type-gives-well-defined-system}
Situation \ref{spaces-limits-situation-limit-noetherian}에서 Lemmas
\ref{spaces-limits-lemma-good-diagram}, \ref{spaces-limits-lemma-limit-from-good-diagram},
\ref{spaces-limits-lemma-morphism-good-diagram}은 $B$ 위에서 준분리이고 유한형인
대수공간들의 범주가 $(B_i)_{i \in I}$ 위 대수공간들의 특정한 유형의
역계 범주와 동치임을 말한다. 여기서 그 역계들은 Lemma
\ref{spaces-limits-lemma-limit-from-good-diagram}을
(\ref{spaces-limits-equation-good-diagram}) 꼴의 그림에 적용하여 얻는 것들이다.
예를 들어 유한형이고 준분리인 $X \to B$가 주어졌을 때,
서로 다른 두 그림 $X \to V_1 \to B_{i_1}$와
$X \to V_2 \to B_{i_2}$를 택하고, 이들이
(\ref{spaces-limits-equation-good-diagram})에서와 같다고 하자.
그러면 $\text{id}_X$에 Lemma \ref{spaces-limits-lemma-morphism-good-diagram}을
(두 방향으로) 적용함으로써, 이에 대응하는 $X$의 극한 표현들이
표준적으로 동형임을 알 수 있다(유향 집합 $I$를 축소하는 것까지 허용한다).
이와 같은 방식으로 계속된다.
\end{remark}

\begin{lemma}
\label{spaces-limits-lemma-morphism-good-diagram-flat}
표기와 가정은 Lemma \ref{spaces-limits-lemma-morphism-good-diagram}에서와 같다.
$f$가 평탄하고 유한 표시이면,
% P09-ERR-0980: frozen source uses the strict threshold $i_3 > i_0$; preserved.
$i_3 > i_0$인 것이 존재하여 $i \geq i_3$이면
$f_i$는 평탄하고, $X_i = Y_i \times_{Y_{i_3}} X_{i_3}$이며,
$X = Y \times_{Y_{i_3}} X_{i_3}$이다.
\end{lemma}

\begin{proof}
Lemma \ref{spaces-limits-lemma-descend-finite-presentation}에 의해
$i \geq i_2$와 유한 표시 사상
$U \to Y_i$를 택하여 $X = Y \times_{Y_i} U$로 만들 수 있다
(여기에서 $f$가 유한 표시라는 가정을 사용한다).
$i$를 증가시킨 뒤에는 $U \to Y_i$가 평탄하다고 가정할 수 있다.
Lemma \ref{spaces-limits-lemma-descend-flat}을 보라.
Remark \ref{spaces-limits-remark-finite-type-gives-well-defined-system}에서 논의했듯이,
계 $(X_i)_{i \geq i_1}$를 정의하는 데 사용한 처음 그림을
$X \to U \to B_i$에 대응하는 계로 바꾸어도 되고, 그렇게 하겠다.
따라서 $X_{i'}$는 $i' \geq i$일 때
$X \to B_{i'} \times_{B_i} U$의 스킴론적 상으로 정의된다.

\medskip\noindent
$U \to Y_i$가 평탄하고(여기에서 $f$가 평탄하다는 가정을 사용한다),
$X = Y \times_{Y_i} U$이며, $Y \to Y_i$의 스킴론적 상이 $Y_i$이므로,
$X \to U$의 스킴론적 상은 $U$이다
(Morphisms of Spaces, Lemma
\ref{spaces-morphisms-lemma-flat-base-change-scheme-theoretic-image}).
$Y_{i'} \to B_{i'} \times_{B_i} Y_i$는 $i' \geq i$일 때
$Y_j$의 계를 구성한 방식에 의해 닫힌 몰입임을 주목하자.
그러면 위와 같은 논증으로 $X \to B_{i'} \times_{B_i} U$의
스킴론적 상이 닫힌 부분공간 $Y_{i'} \times_{Y_i} U$와 같음을 알 수 있다.
따라서 $X_{i'} = Y_{i'} \times_{Y_i} U$이고, 이는 모든
$i' \geq i$에 대해 성립한다. 그러므로 $i_3 = i$로 두면
보조정리가 성립한다.
\end{proof}

\begin{lemma}
\label{spaces-limits-lemma-morphism-good-diagram-smooth}
표기와 가정은 Lemma \ref{spaces-limits-lemma-morphism-good-diagram}에서와 같다.
$f$가 매끄러우면,
% P09-ERR-0980: frozen source uses the strict threshold $i_3 > i_0$; preserved.
$i_3 > i_0$인 것이 존재하여 $i \geq i_3$이면 $f_i$는 매끄럽다.
\end{lemma}

\begin{proof}
Lemmas \ref{spaces-limits-lemma-morphism-good-diagram-flat}과
\ref{spaces-limits-lemma-descend-smooth}를 결합하라.
\end{proof}

\begin{lemma}
\label{spaces-limits-lemma-morphism-good-diagram-proper}
표기와 가정은 Lemma \ref{spaces-limits-lemma-morphism-good-diagram}에서와 같다.
$f$가 고유이면, $i_3 \geq i_0$인 것이 존재하여
$i \geq i_3$이면 $f_i$는 고유이다.
\end{lemma}

\begin{proof}
Remark \ref{spaces-limits-remark-finite-type-gives-well-defined-system}의 논의에 따르면,
(\ref{spaces-limits-equation-good-diagram})에서와 같은 그림에 들어가는
$i_1$과 $W$의 선택은 보조정리의 참 여부에 영향을 주지 않는다.
따라서 $W$를 다음과 같이 택한다. 먼저 닫힌 몰입 $X \to X'$를 택하여
$X' \to Y$가 고유이고 유한 표시가 되게 한다. Lemma
\ref{spaces-limits-lemma-proper-limit-of-proper-finite-presentation}을 보라.
그다음 $i_3 \geq i_2$와 고유 사상 $W \to Y_{i_3}$를 택하여
$X' = Y \times_{Y_{i_3}} W$로 만든다. 이는
$Y = \lim_{i \geq i_2} Y_i$이고
% P09-ERR-0143: frozen source omits “by” before these lemmas; Korean normalizes grammar.
Lemmas \ref{spaces-limits-lemma-relative-approximation}과
\ref{spaces-limits-lemma-eventually-proper}에 의해 가능하다.
이렇게 $W$를 택하면, 구성으로부터 곧바로 $i \geq i_3$일 때
대수공간 $X_i$가
$Y_i \times_{Y_{i_3}} W \subset B_i \times_{B_{i_3}} W$의
닫힌 부분공간이고, 따라서 $Y_i$ 위에서 고유임을 알 수 있다.
\end{proof}

\begin{lemma}
\label{spaces-limits-lemma-good-diagram-fibre-product}
Situation \ref{spaces-limits-situation-limit-noetherian}에서 다음 올곱 그림이 주어졌다고 하자.
$$
\xymatrix{
X^1 \ar[r]_p \ar[d]_q & X^3 \ar[d]^a \\
X^2 \ar[r]^b & X^4
}
$$
여기서 각 대수공간은 $B$ 위에서 준분리이고 유한형이다.
각 $j = 1, 2, 3, 4$에 대해 $i_j \in I$와 그림
$$
\xymatrix{
X^j \ar[r] \ar[d] & W^j \ar[d] \\
B \ar[r] & B_{i_j}
}
$$
을 (\ref{spaces-limits-equation-good-diagram})에서와 같이 택하자. 대응하는 극한 표현을
$X^j = \lim_{i \geq i_j} X^j_i$라 하자. 이는 Lemma
\ref{spaces-limits-lemma-morphism-good-diagram}에서와 같다.
% P09-ERR-0144: frozen source cites lemma-morphism-good-diagram here; preserved.
Lemma \ref{spaces-limits-lemma-morphism-good-diagram}에서 구성한 대응하는 역계 사상들을
$(a_i)_{i \geq i_5}$, $(b_i)_{i \geq i_6}$,
$(p_i)_{i \geq i_7}$, $(q_i)_{i \geq i_8}$라 하자.
그러면 $i_9 \geq \max(i_5, i_6, i_7, i_8)$인 것이 존재하여
$i \geq i_9$이면 $a_i \circ p_i = b_i \circ q_i$이고,
$$
(q_i, p_i) : X^1_i \longrightarrow X^2_i \times_{b_i, X^4_i, a_i} X^3_i
$$
는 닫힌 몰입이다.
$a$와 $b$가 평탄하고 유한 표시이면,
$i_{10} \geq \max(i_5, i_6, i_7, i_8, i_9)$인 것이 존재하여
$i \geq i_{10}$이면 마지막 표시 사상은 동형이다.
\end{lemma}

\begin{proof}
Remark \ref{spaces-limits-remark-finite-type-gives-well-defined-system}의 논의에 따르면,
(\ref{spaces-limits-equation-good-diagram})에서와 같은 그림에 들어가는 $W^1$의 선택은
보조정리의 참 여부에 영향을 주지 않는다. 따라서
$W^1 = W^2 \times_{W^4} W^3$로 택해도 된다.
그러면 $X^1_i$의 구성으로부터 곧바로
$a_i \circ p_i = b_i \circ q_i$이고
$$
(q_i, p_i) : X^1_i \longrightarrow X^2_i \times_{b_i, X^4_i, a_i} X^3_i
$$
가 닫힌 몰입임을 알 수 있다.

\medskip\noindent
$a$와 $b$가 평탄하고 유한 표시이면, $p$와 $q$도
$a$와 $b$의 밑변환으로서 그러하다. 따라서 Lemma
\ref{spaces-limits-lemma-morphism-good-diagram-flat}을
$a$, $b$, $p$, $q$, 그리고 $a \circ p = b \circ q$ 각각에 적용할 수 있다.
% P09-ERR-0145: frozen source reuses threshold $i_9$ instead of statement's $i_{10}$; preserved.
그 결과 $i_9 \in I$가 존재하여
$$
(q_i, p_i) : X^1_i \to X^2_i \times_{X^4_i} X^3_i
$$
가 $(q_{i_9}, p_{i_9})$의 사상
% P09-ERR-0146: frozen source repeats “by the morphism”; Korean normalizes grammar.
$X^4_i \to X^4_{i_9}$에 의한 밑변환이고, 이는 모든
$i \geq i_9$에 대해 성립한다. Lemma \ref{spaces-limits-lemma-descend-isomorphism}에 의해
$(q_i, p_i)$가 충분히 큰 모든 $i$에 대해 동형이라고 결론짓는다.
\end{proof}
